\documentclass[a4paper,12pt]{book}

% === Encodage et langue ===
\usepackage[utf8]{inputenc}
\usepackage[T1]{fontenc}
\usepackage[french]{babel}

% === Mathématiques ===
\usepackage{amsmath}
\usepackage{amssymb}
\usepackage{amsthm}
\usepackage{mathtools}
\usepackage{stmaryrd}

% === Boîtes colorées ===
\usepackage[theorems,skins,breakable]{tcolorbox}

% === Diagrammes et graphiques ===
\usepackage{tikz}
\usetikzlibrary{calc}
\usepackage{tikz-cd}
\usepackage{pgfplots}
\pgfplotsset{compat=1.18}

% === Figures et tableaux ===
\usepackage{graphicx}
\usepackage{float}
\usepackage{caption}
\usepackage{booktabs}
\usepackage{array}
\usepackage{longtable}

% === Listes ===
\usepackage{enumitem}

% === Références et index ===
\usepackage{hyperref}
\usepackage{cleveref}
\usepackage{makeidx}
\usepackage{minitoc}

\makeindex

% ============================================================
% Opérateurs mathématiques
% ============================================================
\DeclareMathOperator{\Hom}{Hom}
\DeclareMathOperator{\End}{End}
\DeclareMathOperator{\Aut}{Aut}
\DeclareMathOperator{\im}{im}
\DeclareMathOperator{\Id}{Id}
\DeclareMathOperator{\Ker}{Ker}
\DeclareMathOperator{\coker}{coker}
\DeclareMathOperator{\Ob}{Ob}
\DeclareMathOperator{\Mor}{Mor}
\DeclareMathOperator{\Fun}{Fun}
\DeclareMathOperator{\Nat}{Nat}
\DeclareMathOperator{\colim}{colim}
\DeclareMathOperator{\eq}{eq}
\DeclareMathOperator{\coeq}{coeq}

% Catégories usuelles
\DeclareMathOperator{\Set}{\mathbf{Set}}
\DeclareMathOperator{\Grp}{\mathbf{Grp}}
\DeclareMathOperator{\Ab}{\mathbf{Ab}}
\DeclareMathOperator{\Ring}{\mathbf{Ring}}
\DeclareMathOperator{\Top}{\mathbf{Top}}
\DeclareMathOperator{\Vect}{\mathbf{Vect}}
\DeclareMathOperator{\Mod}{\mathbf{Mod}}
\DeclareMathOperator{\Cat}{\mathbf{Cat}}

% Commandes utiles
\newcommand{\op}{\mathrm{op}}
\newcommand{\id}{\mathrm{id}}
\newcommand{\N}{\mathbb{N}}
\newcommand{\Z}{\mathbb{Z}}
\newcommand{\Q}{\mathbb{Q}}
\newcommand{\R}{\mathbb{R}}
\newcommand{\C}{\mathbb{C}}
\newcommand{\K}{\mathbb{K}}

% ============================================================
% Environnements de théorèmes
% ============================================================
\newtheorem{theoreme}{Théorème}[chapter]
\newtheorem{proposition}[theoreme]{Proposition}
\newtheorem{lemme}[theoreme]{Lemme}
\newtheorem{corollaire}[theoreme]{Corollaire}
\newtheorem{definition}[theoreme]{Définition}
\newtheorem{exemple}[theoreme]{Exemple}
\newtheorem{exercice}{Exercice}[chapter]
\newtheorem*{remarque}{Remarque}
\newtheorem*{notation}{Notation}

% === Styles tcolorbox ===
\tcolorboxenvironment{definition}{
  enhanced, breakable,
  colback=blue!5, colframe=blue!70!black,
  boxrule=0.8pt, arc=2pt,
  left=6pt, right=6pt, top=4pt, bottom=4pt
}
\tcolorboxenvironment{theoreme}{
  enhanced, breakable,
  colback=red!5, colframe=red!70!black,
  boxrule=0.8pt, arc=2pt,
  left=6pt, right=6pt, top=4pt, bottom=4pt
}
\tcolorboxenvironment{proposition}{
  enhanced, breakable,
  colback=red!3, colframe=red!60!black,
  boxrule=0.8pt, arc=2pt,
  left=6pt, right=6pt, top=4pt, bottom=4pt
}
\tcolorboxenvironment{lemme}{
  enhanced, breakable,
  colback=yellow!10, colframe=yellow!70!black,
  boxrule=0.8pt, arc=2pt,
  left=6pt, right=6pt, top=4pt, bottom=4pt
}
\tcolorboxenvironment{corollaire}{
  enhanced, breakable,
  colback=orange!5, colframe=orange!70!black,
  boxrule=0.8pt, arc=2pt,
  left=6pt, right=6pt, top=4pt, bottom=4pt
}
\tcolorboxenvironment{exemple}{
  enhanced, breakable,
  colback=green!5, colframe=green!60!black,
  boxrule=0.8pt, arc=2pt,
  left=6pt, right=6pt, top=4pt, bottom=4pt
}
\tcolorboxenvironment{exercice}{
  enhanced, breakable,
  colback=purple!5, colframe=purple!70!black,
  boxrule=0.8pt, arc=2pt,
  left=6pt, right=6pt, top=4pt, bottom=4pt
}
\tcolorboxenvironment{remarque}{
  enhanced, breakable,
  colback=gray!5, colframe=gray!60!black,
  boxrule=0.8pt, arc=2pt,
  left=6pt, right=6pt, top=4pt, bottom=4pt
}

% ============================================================
\title{\Huge\textbf{Théorie des Catégories} \\ \Large Cours de Master}
\author{Notes de cours}
\date{2025--2026}

\begin{document}

\begin{titlepage}
\begin{tikzpicture}[remember picture, overlay]
  \fill[left color=blue!15!white, right color=blue!5!white]
    (current page.south west) rectangle (current page.north east);
  \fill[color=orange!70!yellow]
    ([yshift=-2cm]current page.north west) rectangle ([yshift=-2.3cm]current page.north east);
  \fill[color=blue!80!black, rounded corners=8pt]
    ([xshift=3cm, yshift=-4cm]current page.north west) rectangle
    ([xshift=-3cm, yshift=-10cm]current page.north east);
  \node[anchor=center, text=white, font=\Huge\bfseries]
    at ([yshift=-6cm]current page.north) {Th\'eorie des Cat\'egories};
  \node[anchor=center, text=orange!80!yellow, font=\Large]
    at ([yshift=-7.5cm]current page.north) {Notes de Cours};
  \node[anchor=center, text=white!80, font=\large]
    at ([yshift=-8.8cm]current page.north) {M1--Doctorat --- 2025--2026};
  \node[anchor=center, text=blue!80!black, font=\Large\itshape]
    at ([yshift=-12cm]current page.north) {Ya\'e Ulrich Gaba};
  \draw[orange!70!yellow, line width=2pt]
    ([xshift=5cm, yshift=-13.5cm]current page.north west) --
    ([xshift=-5cm, yshift=-13.5cm]current page.north east);
  \node[anchor=center, text width=12cm, align=center, text=blue!60!black, font=\itshape]
    at ([yshift=-16cm]current page.north) {<<~J'ai pens\'e que la notion de cat\'egorie \'etait n\'ecessaire pour axiomatiser les transformations naturelles.~>> --- Saunders Mac\,Lane};
  \node[anchor=south, text=gray, font=\small]
    at ([yshift=1.5cm]current page.south) {\today};
\end{tikzpicture}
\end{titlepage}

\dominitoc
\tableofcontents

% ============================================================
% PRÉFACE
% ============================================================
\chapter*{Préface}
\addcontentsline{toc}{chapter}{Préface}

La théorie des catégories\index{catégorie}, née des travaux fondateurs d'Eilenberg et Mac\,Lane dans les années 1940, s'est imposée comme le langage unificateur des mathématiques modernes. Conçue initialement pour formaliser la notion de \emph{transformation naturelle}\index{transformation naturelle} en topologie algébrique, elle a trouvé des applications profondes dans pratiquement toutes les branches des mathématiques.

C'est avant tout dans l'\oe uvre monumentale d'Alexandre Grothendieck\index{Grothendieck} que la théorie des catégories a révélé sa pleine puissance. En refondant la géométrie algébrique sur des bases catégoriques --- schémas, faisceaux, topos --- Grothendieck a démontré qu'un cadre abstrait bien choisi permet non seulement d'unifier des théories existantes, mais de découvrir des structures invisibles dans les formulations classiques.

Ce cours s'adresse aux étudiants de Master souhaitant acquérir une maîtrise solide des concepts fondamentaux : catégories, foncteurs, transformations naturelles, limites, adjonctions, et au-delà. L'accent est mis sur la rigueur des définitions et la richesse des exemples, issus de l'algèbre, de la topologie et de l'analyse fonctionnelle.

\bigskip

\chapter*{Notations}
\addcontentsline{toc}{chapter}{Notations}

\begin{notation}
Tout au long de ce texte, on adopte les conventions suivantes :
\begin{itemize}[leftmargin=2cm]
  \item[$\Set$] --- catégorie des ensembles et applications.
  \item[$\Grp$] --- catégorie des groupes et homomorphismes.
  \item[$\Ab$] --- catégorie des groupes abéliens et homomorphismes.
  \item[$\Ring$] --- catégorie des anneaux (unitaires) et homomorphismes d'anneaux.
  \item[$\Top$] --- catégorie des espaces topologiques et applications continues.
  \item[$\Vect_\K$] --- catégorie des $\K$-espaces vectoriels et applications linéaires.
  \item[$R\text{-}\Mod$] --- catégorie des $R$-modules à gauche et homomorphismes de modules.
  \item[$\Cat$] --- catégorie des petites catégories et foncteurs.
  \item[$\mathcal{C}^{\op}$] --- catégorie opposée\index{catégorie!opposée} de $\mathcal{C}$.
  \item[$\Hom_{\mathcal{C}}(A,B)$] --- ensemble des morphismes de $A$ vers $B$ dans $\mathcal{C}$.
  \item[$\id_A$] --- morphisme identité de l'objet $A$.
  \item[$F \circ G$] --- composition de foncteurs.
  \item[$\alpha : F \Rightarrow G$] --- transformation naturelle de $F$ vers $G$.
  \item[{$[\mathcal{C}, \mathcal{D}]$}] --- catégorie des foncteurs de $\mathcal{C}$ vers $\mathcal{D}$.
\end{itemize}
\end{notation}

% ============================================================
% CHAPITRE 1
% ============================================================
\chapter{Catégories, Foncteurs, Transformations Naturelles}
\minitoc

\section{Catégories}

\subsection{Définition d'une catégorie}

\begin{definition}[Catégorie]\label{def:categorie}
\index{catégorie!définition}
Une \textbf{catégorie} $\mathcal{C}$ est la donnée de :
\begin{enumerate}[label=(\roman*)]
  \item une classe $\Ob(\mathcal{C})$ dont les éléments sont appelés \textbf{objets}\index{objet} de $\mathcal{C}$ ;
  \item pour tout couple $(A, B)$ d'objets, un ensemble $\Hom_{\mathcal{C}}(A, B)$ dont les éléments sont appelés \textbf{morphismes}\index{morphisme} (ou \textbf{flèches}) de $A$ vers $B$ ;
  \item pour tout triplet $(A, B, C)$ d'objets, une loi de \textbf{composition}\index{composition}
  \[
    \circ : \Hom_{\mathcal{C}}(B, C) \times \Hom_{\mathcal{C}}(A, B) \longrightarrow \Hom_{\mathcal{C}}(A, C), \quad (g, f) \longmapsto g \circ f ;
  \]
  \item pour tout objet $A$, un morphisme $\id_A \in \Hom_{\mathcal{C}}(A, A)$ appelé \textbf{identité}\index{morphisme!identité} de $A$ ;
\end{enumerate}
satisfaisant les axiomes suivants :
\begin{enumerate}[label=(C\arabic*)]
  \item \textbf{Associativité}\index{associativité} : pour tous morphismes $f : A \to B$, $g : B \to C$, $h : C \to D$,
  \[ h \circ (g \circ f) = (h \circ g) \circ f. \]
  \item \textbf{Identité} : pour tout morphisme $f : A \to B$,
  \[ f \circ \id_A = f = \id_B \circ f. \]
\end{enumerate}
On exige de plus que les ensembles $\Hom_{\mathcal{C}}(A, B)$ soient deux à deux disjoints pour des couples $(A,B)$ distincts (condition de « typage » des morphismes).
\end{definition}

\begin{remarque}
Si $f \in \Hom_{\mathcal{C}}(A,B)$, on note aussi $f : A \to B$ ou $A \xrightarrow{f} B$. L'objet $A$ est la \textbf{source}\index{source} (ou domaine) de $f$ et $B$ est le \textbf{but}\index{but} (ou codomaine).
\end{remarque}

\subsection{Exemples fondamentaux}

\begin{exemple}[Catégories concrètes]\label{ex:categories-concretes}
\index{catégorie!concrète}
\begin{enumerate}[label=(\alph*)]
  \item $\Set$\index{Set@$\Set$} : objets = ensembles, morphismes = applications, composition usuelle.
  \item $\Grp$\index{Grp@$\Grp$} : objets = groupes, morphismes = homomorphismes de groupes.
  \item $\Ab$\index{Ab@$\Ab$} : objets = groupes abéliens, morphismes = homomorphismes de groupes.
  \item $\Ring$\index{Ring@$\Ring$} : objets = anneaux unitaires, morphismes = homomorphismes d'anneaux (préservant l'unité).
  \item $\Top$\index{Top@$\Top$} : objets = espaces topologiques, morphismes = applications continues.
  \item $\Vect_\K$\index{Vect@$\Vect_\K$} : objets = $\K$-espaces vectoriels, morphismes = applications $\K$-linéaires.
  \item $R$-$\Mod$\index{Mod@$R$-$\Mod$} : objets = $R$-modules à gauche, morphismes = homomorphismes de $R$-modules.
\end{enumerate}
\end{exemple}

\begin{exemple}[Ensembles ordonnés comme catégories]\label{ex:poset-categorie}
\index{catégorie!ensemble ordonné}
Soit $(P, \leq)$ un ensemble partiellement ordonné\index{ensemble ordonné}. On définit une catégorie $\mathcal{C}_P$ par :
\begin{itemize}
  \item $\Ob(\mathcal{C}_P) = P$ ;
  \item $\Hom_{\mathcal{C}_P}(x, y) = \begin{cases} \{(x,y)\} & \text{si } x \leq y, \\ \varnothing & \text{sinon.} \end{cases}$
\end{itemize}
La transitivité de $\leq$ assure l'existence de la composition, la réflexivité fournit les identités, et l'associativité est automatique (unicité des morphismes).
\end{exemple}

\begin{exemple}[Monoïdes comme catégories]\label{ex:monoide-categorie}
\index{catégorie!monoïde}\index{monoïde}
Tout monoïde $(M, \cdot, e)$ définit une catégorie $\mathcal{B}M$ à un seul objet $\star$ :
\[
  \Hom_{\mathcal{B}M}(\star, \star) = M, \quad g \circ f = g \cdot f, \quad \id_\star = e.
\]
Inversement, toute catégorie à un seul objet est (les données d'un) monoïde.
\end{exemple}

\begin{exemple}[Groupes comme catégories]\label{ex:groupe-categorie}
\index{catégorie!groupe}
Un groupe $G$ définit une catégorie $\mathcal{B}G$ à un seul objet où \emph{tous} les morphismes sont inversibles. Réciproquement, une catégorie à un seul objet dont tous les morphismes sont inversibles est (les données d')un groupe.
\end{exemple}

\subsection{Petitesse et taille}

\begin{definition}[Petite catégorie, catégorie localement petite]\label{def:petite-categorie}
\index{catégorie!petite}\index{catégorie!localement petite}
\begin{enumerate}[label=(\roman*)]
  \item Une catégorie $\mathcal{C}$ est \textbf{petite} si $\Ob(\mathcal{C})$ est un ensemble (et non une classe propre).
  \item Une catégorie $\mathcal{C}$ est \textbf{localement petite} si pour tous objets $A, B$, l'ensemble $\Hom_{\mathcal{C}}(A,B)$ est un ensemble (au sens de la théorie des ensembles fixée).
\end{enumerate}
\end{definition}

\begin{remarque}
Les catégories $\Set$, $\Grp$, $\Top$, etc.\ ne sont pas petites (la classe de tous les ensembles n'est pas un ensemble), mais elles sont localement petites. La catégorie des ensembles finis $\mathbf{FinSet}$ est un exemple de catégorie à la fois petite et localement petite.
\end{remarque}

\subsection{Monomorphismes, épimorphismes, isomorphismes}

\begin{definition}[Monomorphisme]\label{def:mono}
\index{monomorphisme}\index{morphisme!mono}
Un morphisme $f : A \to B$ dans $\mathcal{C}$ est un \textbf{monomorphisme} (ou \textbf{mono}) si pour tous morphismes $g, h : C \to A$,
\[
  f \circ g = f \circ h \implies g = h.
\]
On dit que $f$ est \emph{annulable à gauche}.
\end{definition}

\begin{definition}[Épimorphisme]\label{def:epi}
\index{épimorphisme}\index{morphisme!épi}
Un morphisme $f : A \to B$ dans $\mathcal{C}$ est un \textbf{épimorphisme} (ou \textbf{épi}) si pour tous morphismes $g, h : B \to C$,
\[
  g \circ f = h \circ f \implies g = h.
\]
On dit que $f$ est \emph{annulable à droite}.
\end{definition}

\begin{definition}[Isomorphisme]\label{def:iso}
\index{isomorphisme}\index{morphisme!iso}
Un morphisme $f : A \to B$ est un \textbf{isomorphisme} s'il existe $g : B \to A$ tel que
\[
  g \circ f = \id_A \quad \text{et} \quad f \circ g = \id_B.
\]
Un tel $g$ est unique et noté $f^{-1}$. On écrit $A \cong B$\index{$\cong$} si un tel isomorphisme existe.
\end{definition}

\begin{proposition}\label{prop:iso-mono-epi}
Tout isomorphisme est à la fois un monomorphisme et un épimorphisme.
\end{proposition}

\begin{proof}
Soit $f : A \to B$ un isomorphisme d'inverse $f^{-1}$. Si $f \circ g = f \circ h$, alors $g = f^{-1} \circ f \circ g = f^{-1} \circ f \circ h = h$, donc $f$ est mono. De même pour épi par composition à droite avec $f^{-1}$.
\end{proof}

\begin{remarque}
La réciproque est fausse en général. Dans $\Ring$, l'inclusion $\Z \hookrightarrow \Q$ est à la fois mono et épi, mais n'est pas un isomorphisme.
\end{remarque}

\begin{proposition}\label{prop:comp-mono-epi}
\begin{enumerate}[label=(\roman*)]
  \item La composée de deux monomorphismes est un monomorphisme.
  \item La composée de deux épimorphismes est un épimorphisme.
  \item Si $g \circ f$ est mono, alors $f$ est mono.
  \item Si $g \circ f$ est épi, alors $g$ est épi.
\end{enumerate}
\end{proposition}

\begin{proof}
(i) Soient $f : A \to B$ et $g : B \to C$ monos. Si $(g \circ f) \circ h_1 = (g \circ f) \circ h_2$, alors $g \circ (f \circ h_1) = g \circ (f \circ h_2)$, d'où $f \circ h_1 = f \circ h_2$ ($g$ mono), puis $h_1 = h_2$ ($f$ mono).

(iii) Si $g \circ f$ est mono et $f \circ h_1 = f \circ h_2$, alors $(g \circ f) \circ h_1 = (g \circ f) \circ h_2$, d'où $h_1 = h_2$. Les points (ii) et (iv) sont duaux.
\end{proof}

\begin{lemme}\label{lem:unicite-inverse}
\index{inverse!unicité}
L'inverse d'un isomorphisme est unique.
\end{lemme}

\begin{proof}
Si $g$ et $g'$ sont deux inverses de $f$, alors $g = g \circ \id_B = g \circ (f \circ g') = (g \circ f) \circ g' = \id_A \circ g' = g'$.
\end{proof}

\subsection{Objets initiaux, terminaux, nuls}

\begin{definition}[Objet initial, terminal, nul]\label{def:init-term-nul}
\index{objet!initial}\index{objet!terminal}\index{objet!nul}
Soit $\mathcal{C}$ une catégorie.
\begin{enumerate}[label=(\roman*)]
  \item Un objet $I$ est \textbf{initial} si pour tout objet $A$, il existe un unique morphisme $I \to A$.
  \item Un objet $T$ est \textbf{terminal} si pour tout objet $A$, il existe un unique morphisme $A \to T$.
  \item Un objet est \textbf{nul} (ou \textbf{zéro}) s'il est à la fois initial et terminal.
\end{enumerate}
\end{definition}

\begin{proposition}\label{prop:init-term-unique}
Si $I_1$ et $I_2$ sont deux objets initiaux de $\mathcal{C}$, alors $I_1 \cong I_2$, et l'isomorphisme est unique. De même pour les objets terminaux.
\end{proposition}

\begin{proof}
Par la propriété universelle, il existe des morphismes uniques $f : I_1 \to I_2$ et $g : I_2 \to I_1$. Alors $g \circ f : I_1 \to I_1$ et $\id_{I_1} : I_1 \to I_1$ sont deux morphismes de $I_1$ vers lui-même ; par unicité (car $I_1$ est initial), $g \circ f = \id_{I_1}$. De même $f \circ g = \id_{I_2}$.
\end{proof}

\begin{exemple}[Objets initiaux et terminaux]\label{ex:init-term}
\index{objet!initial!exemples}\index{objet!terminal!exemples}
\begin{enumerate}[label=(\alph*)]
  \item Dans $\Set$ : $\varnothing$ est initial, tout singleton est terminal.
  \item Dans $\Grp$ et $\Ab$ : le groupe trivial $\{e\}$ est objet nul.
  \item Dans $\Ring$ : $\Z$ est initial, l'anneau nul $\{0\}$ est terminal.
  \item Dans $\Top$ : $\varnothing$ est initial, tout singleton (avec la topologie triviale) est terminal.
  \item Dans $\Vect_\K$ : l'espace nul $\{0\}$ est objet nul.
\end{enumerate}
\end{exemple}

% ============================================================
\section{Foncteurs}

\subsection{Définition et premiers exemples}

\begin{definition}[Foncteur (covariant)]\label{def:foncteur}
\index{foncteur!covariant}\index{foncteur!définition}
Soient $\mathcal{C}$ et $\mathcal{D}$ deux catégories. Un \textbf{foncteur} (covariant) $F : \mathcal{C} \to \mathcal{D}$ est la donnée de :
\begin{enumerate}[label=(\roman*)]
  \item une application $F : \Ob(\mathcal{C}) \to \Ob(\mathcal{D})$ ;
  \item pour tout couple $(A, B)$ d'objets de $\mathcal{C}$, une application
  \[
    F_{A,B} : \Hom_{\mathcal{C}}(A, B) \longrightarrow \Hom_{\mathcal{D}}(F(A), F(B))
  \]
  (on note simplement $F(f)$ pour $F_{A,B}(f)$) ;
\end{enumerate}
satisfaisant :
\begin{enumerate}[label=(F\arabic*)]
  \item $F(\id_A) = \id_{F(A)}$ pour tout objet $A$ de $\mathcal{C}$ ;
  \item $F(g \circ f) = F(g) \circ F(f)$ pour tous morphismes composables $f, g$ dans $\mathcal{C}$.
\end{enumerate}
\end{definition}

\begin{definition}[Foncteur contravariant]\label{def:foncteur-contra}
\index{foncteur!contravariant}
Un \textbf{foncteur contravariant} $F : \mathcal{C} \to \mathcal{D}$ est un foncteur $F : \mathcal{C}^{\op} \to \mathcal{D}$, c'est-à-dire qu'il \emph{renverse les flèches} :
\[
  F(g \circ f) = F(f) \circ F(g).
\]
\end{definition}

\begin{exemple}[Foncteurs classiques]\label{ex:foncteurs}
\index{foncteur!exemples}
\begin{enumerate}[label=(\alph*)]
  \item \textbf{Foncteur d'oubli}\index{foncteur!d'oubli} $U : \Grp \to \Set$ : à un groupe $G$ associe l'ensemble sous-jacent, à un homomorphisme l'application sous-jacente.
  \item \textbf{Foncteur libre}\index{foncteur!libre} $F : \Set \to \Grp$ : à un ensemble $X$ associe le groupe libre $F(X)$.
  \item \textbf{Foncteur abélianisation}\index{abélianisation} $\mathrm{ab} : \Grp \to \Ab$ : $G \mapsto G/[G,G]$.
  \item \textbf{Foncteur espace dual}\index{foncteur!dual} $(-)^* : \Vect_\K \to \Vect_\K^{\op}$ : $V \mapsto V^* = \Hom_\K(V, \K)$, contravariant.
  \item \textbf{Foncteur $\pi_1$}\index{foncteur!$\pi_1$} : $\Top_* \to \Grp$, le groupe fondamental.
  \item \textbf{Foncteur identité}\index{foncteur!identité} $\Id_\mathcal{C} : \mathcal{C} \to \mathcal{C}$ : $A \mapsto A$, $f \mapsto f$.
  \item \textbf{Foncteur constant}\index{foncteur!constant} $\Delta_D : \mathcal{C} \to \mathcal{D}$ : tout objet sur $D$, tout morphisme sur $\id_D$.
\end{enumerate}
\end{exemple}

\begin{proposition}\label{prop:foncteur-iso}
\index{foncteur!préserve les isomorphismes}
Tout foncteur préserve les isomorphismes : si $f$ est un isomorphisme dans $\mathcal{C}$, alors $F(f)$ est un isomorphisme dans $\mathcal{D}$, d'inverse $F(f^{-1})$.
\end{proposition}

\begin{proof}
$F(f) \circ F(f^{-1}) = F(f \circ f^{-1}) = F(\id_B) = \id_{F(B)}$ et de même pour l'autre composition.
\end{proof}

\subsection{Fidélité, plénitude, essentielle surjectivité}

\begin{definition}[Foncteur fidèle, plein, pleinement fidèle]\label{def:fidele-plein}
\index{foncteur!fidèle}\index{foncteur!plein}\index{foncteur!pleinement fidèle}
Un foncteur $F : \mathcal{C} \to \mathcal{D}$ est :
\begin{enumerate}[label=(\roman*)]
  \item \textbf{fidèle} si pour tous $A, B \in \Ob(\mathcal{C})$, l'application $F_{A,B} : \Hom_\mathcal{C}(A,B) \to \Hom_\mathcal{D}(F(A), F(B))$ est injective ;
  \item \textbf{plein} si chaque $F_{A,B}$ est surjective ;
  \item \textbf{pleinement fidèle} s'il est à la fois plein et fidèle (chaque $F_{A,B}$ est bijective).
\end{enumerate}
\end{definition}

\begin{definition}[Essentiellement surjectif]\label{def:ess-surj}
\index{foncteur!essentiellement surjectif}
Un foncteur $F : \mathcal{C} \to \mathcal{D}$ est \textbf{essentiellement surjectif} si pour tout objet $D \in \Ob(\mathcal{D})$, il existe $C \in \Ob(\mathcal{C})$ tel que $F(C) \cong D$.
\end{definition}

\begin{exemple}\label{ex:foncteur-fidele}
\begin{enumerate}[label=(\alph*)]
  \item Le foncteur d'oubli $U : \Grp \to \Set$ est fidèle mais pas plein.
  \item L'inclusion $\Ab \hookrightarrow \Grp$ est pleinement fidèle.
  \item Le foncteur $\mathrm{ab} : \Grp \to \Ab$ est plein et essentiellement surjectif (mais pas fidèle).
\end{enumerate}
\end{exemple}

\begin{proposition}\label{prop:plein-fidele-mono}
Un foncteur fidèle reflète les monomorphismes : si $F$ est fidèle et $F(f)$ est mono, alors $f$ est mono.
\end{proposition}

\begin{proof}
Si $f \circ g = f \circ h$, alors $F(f) \circ F(g) = F(f \circ g) = F(f \circ h) = F(f) \circ F(h)$, donc $F(g) = F(h)$ car $F(f)$ est mono. Par fidélité, $g = h$.
\end{proof}

\subsection{Composition de foncteurs}

\begin{proposition}\label{prop:comp-foncteurs}
\index{foncteur!composition}
Si $F : \mathcal{C} \to \mathcal{D}$ et $G : \mathcal{D} \to \mathcal{E}$ sont des foncteurs, alors $G \circ F : \mathcal{C} \to \mathcal{E}$ défini par $(G \circ F)(A) = G(F(A))$ et $(G \circ F)(f) = G(F(f))$ est un foncteur.
\end{proposition}

\begin{proof}
$(G \circ F)(\id_A) = G(F(\id_A)) = G(\id_{F(A)}) = \id_{G(F(A))}$ et $(G \circ F)(g \circ f) = G(F(g \circ f)) = G(F(g) \circ F(f)) = G(F(g)) \circ G(F(f))$.
\end{proof}

% ============================================================
\section{Transformations naturelles}

\subsection{Définition}

\begin{definition}[Transformation naturelle]\label{def:transfo-nat}
\index{transformation naturelle!définition}
Soient $F, G : \mathcal{C} \to \mathcal{D}$ deux foncteurs. Une \textbf{transformation naturelle} $\alpha : F \Rightarrow G$ est la donnée, pour chaque objet $A$ de $\mathcal{C}$, d'un morphisme
\[
  \alpha_A : F(A) \longrightarrow G(A)
\]
dans $\mathcal{D}$, appelé \textbf{composante}\index{composante} de $\alpha$ en $A$, tel que pour tout morphisme $f : A \to B$ dans $\mathcal{C}$, le diagramme suivant commute :
\[
\begin{tikzcd}
  F(A) \arrow[r, "\alpha_A"] \arrow[d, "F(f)"'] & G(A) \arrow[d, "G(f)"] \\
  F(B) \arrow[r, "\alpha_B"'] & G(B)
\end{tikzcd}
\]
Autrement dit : $\alpha_B \circ F(f) = G(f) \circ \alpha_A$ pour tout $f : A \to B$.
\end{definition}

\begin{definition}[Isomorphisme naturel]\label{def:iso-nat}
\index{isomorphisme naturel}\index{transformation naturelle!isomorphisme}
Une transformation naturelle $\alpha : F \Rightarrow G$ est un \textbf{isomorphisme naturel} si chaque composante $\alpha_A$ est un isomorphisme dans $\mathcal{D}$. On écrit alors $F \cong G$.
\end{definition}

\begin{exemple}[Double dual]\label{ex:double-dual}
\index{double dual}
Soit $\K$ un corps. Le morphisme canonique $\eta_V : V \to V^{**}$, $v \mapsto (\varphi \mapsto \varphi(v))$, définit une transformation naturelle $\eta : \Id_{\Vect_\K} \Rightarrow (-)^{**}$.

Pour $f : V \to W$ linéaire, la naturalité revient à $f^{**} \circ \eta_V = \eta_W \circ f$, ce qui se vérifie : pour $v \in V$ et $\psi \in W^*$,
\[
  (f^{**}(\eta_V(v)))(\psi) = \eta_V(v)(\psi \circ f) = (\psi \circ f)(v) = \psi(f(v)) = \eta_W(f(v))(\psi).
\]
En dimension finie, $\eta$ est un isomorphisme naturel.
\end{exemple}

\begin{exemple}[Déterminant]\label{ex:determinant-nat}
\index{déterminant}
Le déterminant $\det : \mathrm{GL}_n \Rightarrow (-)^\times$ est une transformation naturelle entre foncteurs $\Ring \to \Grp$, où $\mathrm{GL}_n(R)$ est le groupe des matrices inversibles et $(-)^\times$ associe le groupe des unités. La naturalité exprime que pour tout homomorphisme d'anneaux $\varphi : R \to S$, on a $\det \circ \mathrm{GL}_n(\varphi) = \varphi^\times \circ \det$.
\end{exemple}

\begin{exemple}[Abélianisation]\label{ex:abel-nat}
\index{abélianisation!transformation naturelle}
La projection canonique $\pi_G : G \twoheadrightarrow G/[G,G]$ définit une transformation naturelle $\pi : U \Rightarrow U \circ \mathrm{ab}$ (ici $U : \Grp \to \Set$ est l'oubli), ou plus précisément $\pi : \Id_{\Grp} \Rightarrow \iota \circ \mathrm{ab}$ où $\iota : \Ab \hookrightarrow \Grp$.
\end{exemple}

\subsection{Composition verticale et horizontale}

\begin{definition}[Composition verticale]\label{def:comp-verticale}
\index{composition!verticale}
Soient $\alpha : F \Rightarrow G$ et $\beta : G \Rightarrow H$ des transformations naturelles entre foncteurs $\mathcal{C} \to \mathcal{D}$. La \textbf{composition verticale} $\beta \circ \alpha : F \Rightarrow H$ est définie par
\[
  (\beta \circ \alpha)_A = \beta_A \circ \alpha_A.
\]
\[
\begin{tikzcd}
  F(A) \arrow[r, "\alpha_A"] \arrow[rr, bend left=40, "(\beta\circ\alpha)_A"] & G(A) \arrow[r, "\beta_A"] & H(A)
\end{tikzcd}
\]
\end{definition}

\begin{proposition}\label{prop:comp-vert-naturelle}
La composition verticale $\beta \circ \alpha$ est bien une transformation naturelle.
\end{proposition}

\begin{proof}
Pour $f : A \to B$ dans $\mathcal{C}$ :
$(\beta \circ \alpha)_B \circ F(f) = \beta_B \circ \alpha_B \circ F(f) = \beta_B \circ G(f) \circ \alpha_A = H(f) \circ \beta_A \circ \alpha_A = H(f) \circ (\beta \circ \alpha)_A$.
\end{proof}

\begin{definition}[Composition horizontale]\label{def:comp-horizontale}
\index{composition!horizontale}
Soient $F, G : \mathcal{C} \to \mathcal{D}$ et $H, K : \mathcal{D} \to \mathcal{E}$ des foncteurs, avec $\alpha : F \Rightarrow G$ et $\beta : H \Rightarrow K$. La \textbf{composition horizontale} $\beta * \alpha : H \circ F \Rightarrow K \circ G$ est définie par
\[
  (\beta * \alpha)_A = \beta_{G(A)} \circ H(\alpha_A) = K(\alpha_A) \circ \beta_{F(A)}.
\]
\[
\begin{tikzcd}
  \mathcal{C}
    \arrow[r, bend left=30, "F", ""{name=U,below}]
    \arrow[r, bend right=30, "G"', ""{name=D}]
  & \mathcal{D}
    \arrow[r, bend left=30, "H", ""{name=V,below}]
    \arrow[r, bend right=30, "K"', ""{name=W}]
  & \mathcal{E}
  \arrow[Rightarrow, from=U, to=D, "\alpha"]
  \arrow[Rightarrow, from=V, to=W, "\beta"]
\end{tikzcd}
\]
\end{definition}

\begin{proposition}[Loi d'échange]\label{prop:loi-echange}
\index{loi d'échange}
Soient $\alpha : F \Rightarrow G$, $\alpha' : G \Rightarrow H$ (foncteurs $\mathcal{C} \to \mathcal{D}$) et $\beta : J \Rightarrow K$, $\beta' : K \Rightarrow L$ (foncteurs $\mathcal{D} \to \mathcal{E}$). Alors
\[
  (\beta' \circ \beta) * (\alpha' \circ \alpha) = (\beta' * \alpha') \circ (\beta * \alpha).
\]
\end{proposition}

\begin{proof}
Vérification composante par composante en utilisant la naturalité de $\beta$ et $\beta'$.
\end{proof}

\subsection{Catégorie de foncteurs}

\begin{definition}[Catégorie de foncteurs]\label{def:cat-foncteurs}
\index{catégorie!de foncteurs}\index{$[\mathcal{C},\mathcal{D}]$}
Soient $\mathcal{C}$ et $\mathcal{D}$ deux catégories (avec $\mathcal{C}$ petite). La \textbf{catégorie de foncteurs} $[\mathcal{C}, \mathcal{D}]$ (aussi notée $\Fun(\mathcal{C}, \mathcal{D})$ ou $\mathcal{D}^\mathcal{C}$) est définie par :
\begin{itemize}
  \item les objets sont les foncteurs $F : \mathcal{C} \to \mathcal{D}$ ;
  \item les morphismes sont les transformations naturelles $\alpha : F \Rightarrow G$ ;
  \item la composition est la composition verticale ;
  \item l'identité de $F$ est la transformation naturelle $\id_F$ définie par $(\id_F)_A = \id_{F(A)}$.
\end{itemize}
\end{definition}

\begin{remarque}
L'hypothèse « $\mathcal{C}$ petite » est nécessaire pour garantir que les transformations naturelles entre deux foncteurs forment un ensemble (et non une classe propre).
\end{remarque}

\begin{proposition}\label{prop:iso-cat-foncteurs}
Un morphisme $\alpha : F \to G$ dans $[\mathcal{C}, \mathcal{D}]$ est un isomorphisme si et seulement si $\alpha$ est un isomorphisme naturel (i.e.\ chaque $\alpha_A$ est un isomorphisme dans $\mathcal{D}$).
\end{proposition}

\begin{proof}
Si $\alpha$ est un isomorphisme dans $[\mathcal{C}, \mathcal{D}]$, il existe $\beta : G \Rightarrow F$ avec $\beta \circ \alpha = \id_F$ et $\alpha \circ \beta = \id_G$. En composantes : $\beta_A \circ \alpha_A = \id_{F(A)}$ et $\alpha_A \circ \beta_A = \id_{G(A)}$, donc chaque $\alpha_A$ est un isomorphisme.

Réciproquement, si chaque $\alpha_A$ est un isomorphisme, on pose $\beta_A = \alpha_A^{-1}$. La naturalité de $\beta$ découle de celle de $\alpha$ : de $\alpha_B \circ F(f) = G(f) \circ \alpha_A$, on tire $F(f) = \alpha_B^{-1} \circ G(f) \circ \alpha_A$, d'où $\beta_B \circ G(f) = F(f) \circ \beta_A$.
\end{proof}

\begin{exemple}[Préfaisceaux]\label{ex:prefaisceaux}
\index{préfaisceau}
Si $\mathcal{C}$ est une petite catégorie, un \textbf{préfaisceau} (d'ensembles) sur $\mathcal{C}$ est un foncteur $F : \mathcal{C}^{\op} \to \Set$. La catégorie des préfaisceaux est $[\mathcal{C}^{\op}, \Set]$, souvent notée $\widehat{\mathcal{C}}$.
\end{exemple}

% ============================================================
\section{Exercices}

\begin{exercice}[Catégorie des flèches]\label{exo:cat-fleches}
\index{catégorie!des flèches}
Soit $\mathcal{C}$ une catégorie. Construire la \textbf{catégorie des flèches} $\mathcal{C}^{\to}$ dont les objets sont les morphismes de $\mathcal{C}$ et dont un morphisme de $f : A \to B$ vers $g : C \to D$ est un couple $(h, k)$ tel que le diagramme suivant commute :
\[
\begin{tikzcd}
  A \arrow[r, "h"] \arrow[d, "f"'] & C \arrow[d, "g"] \\
  B \arrow[r, "k"'] & D
\end{tikzcd}
\]
Montrer que $\mathcal{C}^{\to} \cong [\mathbf{2}, \mathcal{C}]$ où $\mathbf{2}$ est la catégorie $\{0 \to 1\}$.
\end{exercice}

\begin{exercice}[Sections et rétractions]\label{exo:section-retraction}
\index{section}\index{rétraction}
Un morphisme $f : A \to B$ est une \textbf{section} (ou mono scindé) s'il existe $g : B \to A$ tel que $g \circ f = \id_A$. Il est une \textbf{rétraction} (ou épi scindé) s'il existe $h : B \to A$ tel que $f \circ h = \id_B$.
\begin{enumerate}[label=(\alph*)]
  \item Montrer que toute section est un monomorphisme et toute rétraction est un épimorphisme.
  \item Donner un monomorphisme dans $\Set$ qui n'est pas une section.
  \item Montrer que dans $\Set$, tout épimorphisme est une rétraction (utiliser l'axiome du choix).
\end{enumerate}
\end{exercice}

\begin{exercice}[Foncteur puissance]\label{exo:foncteur-puissance}
\index{foncteur!puissance}
Montrer que l'application $\mathcal{P} : \Set \to \Set$ qui à un ensemble $X$ associe $\mathcal{P}(X)$ (ensemble des parties) peut être munie d'une structure de foncteur de deux façons :
\begin{enumerate}[label=(\alph*)]
  \item Covariant : image directe $f_*(S) = f(S)$.
  \item Contravariant : image réciproque $f^*(S) = f^{-1}(S)$.
\end{enumerate}
Vérifier les axiomes de foncteur dans chaque cas.
\end{exercice}

\begin{exercice}[Naturalité du centre]\label{exo:centre}
\index{centre d'un groupe}
Le centre $Z(G) = \{g \in G \mid \forall h \in G,\, gh = hg\}$ définit-il un foncteur $Z : \Grp \to \Ab$ ? Si oui, l'inclusion $Z(G) \hookrightarrow G$ est-elle une transformation naturelle ? Justifier.
\end{exercice}

\begin{exercice}[Équivalence mono/injectif dans Set]\label{exo:mono-set}
Montrer que dans $\Set$, un morphisme est un monomorphisme si et seulement s'il est injectif, et un épimorphisme si et seulement s'il est surjectif.
\end{exercice}

\begin{exercice}[Catégories de tranches]\label{exo:tranche}
\index{catégorie!de tranches}
Soit $\mathcal{C}$ une catégorie et $A$ un objet. Construire la \textbf{catégorie de tranches} $\mathcal{C}/A$ (ou \emph{comma category over $A$}) dont les objets sont les morphismes $f : X \to A$ et les morphismes de $(f : X \to A)$ vers $(g : Y \to A)$ sont les morphismes $h : X \to Y$ tels que $g \circ h = f$. Vérifier les axiomes de catégorie. Décrire $\Set/\{0,1\}$.
\end{exercice}

\begin{exercice}[Endomorphismes de l'identité]\label{exo:endo-id}
\index{transformation naturelle!endomorphismes de l'identité}
Soit $\mathcal{C} = \Ab$. Montrer que $\End({\Id_{\Ab}}) \cong \Z$ (l'anneau des endomorphismes de la transformation naturelle identité est isomorphe à $\Z$).

\emph{Indication} : une transformation naturelle $\alpha : \Id_{\Ab} \Rightarrow \Id_{\Ab}$ est déterminée par $\alpha_\Z(1)$.
\end{exercice}

\begin{exercice}[Foncteur nerf]\label{exo:nerf}
\index{foncteur!nerf}
Soit $\mathcal{C}$ une petite catégorie. Le \textbf{nerf}\index{nerf} $N(\mathcal{C})$ est l'ensemble simplicial défini par $N(\mathcal{C})_n = \{$chaînes de $n$ morphismes composables $A_0 \xrightarrow{f_1} A_1 \xrightarrow{f_2} \cdots \xrightarrow{f_n} A_n\}$.
\begin{enumerate}[label=(\alph*)]
  \item Décrire $N(\mathcal{C})_0$, $N(\mathcal{C})_1$ et $N(\mathcal{C})_2$.
  \item Montrer que $\mathcal{C} \mapsto N(\mathcal{C})$ définit un foncteur $N : \Cat \to \mathbf{sSet}$.
\end{enumerate}
\end{exercice}


% ============================================================
% CHAPITRE 2
% ============================================================
\chapter{Dualité et Principe de Dualité}
\minitoc

\section{La catégorie opposée}

\begin{definition}[Catégorie opposée]\label{def:cat-opposee}
\index{catégorie!opposée}\index{$\mathcal{C}^{\op}$}
Soit $\mathcal{C}$ une catégorie. La \textbf{catégorie opposée} (ou \textbf{duale}) $\mathcal{C}^{\op}$ est définie par :
\begin{enumerate}[label=(\roman*)]
  \item $\Ob(\mathcal{C}^{\op}) = \Ob(\mathcal{C})$ ;
  \item $\Hom_{\mathcal{C}^{\op}}(A, B) = \Hom_{\mathcal{C}}(B, A)$ ;
  \item la composition dans $\mathcal{C}^{\op}$ : si $f^{\op} \in \Hom_{\mathcal{C}^{\op}}(A,B)$ et $g^{\op} \in \Hom_{\mathcal{C}^{\op}}(B,C)$ (correspondant à $f : B \to A$ et $g : C \to B$ dans $\mathcal{C}$), alors $g^{\op} \circ_{\op} f^{\op} = (f \circ g)^{\op}$ ;
  \item $\id_A^{\op} = (\id_A)^{\op}$.
\end{enumerate}
\end{definition}

\begin{proposition}\label{prop:cop-cop}
Pour toute catégorie $\mathcal{C}$, on a $(\mathcal{C}^{\op})^{\op} = \mathcal{C}$.
\end{proposition}

\begin{proof}
Les objets sont identiques. Les morphismes : $\Hom_{(\mathcal{C}^{\op})^{\op}}(A,B) = \Hom_{\mathcal{C}^{\op}}(B,A) = \Hom_\mathcal{C}(A,B)$. La composition est restaurée par double renversement.
\end{proof}

\begin{exemple}\label{ex:cop-exemples}
\begin{enumerate}[label=(\alph*)]
  \item Si $(P, \leq)$ est un ensemble ordonné, $\mathcal{C}_P^{\op} = \mathcal{C}_{(P, \geq)}$.
  \item Si $G$ est un groupe vu comme catégorie à un objet, $(\mathcal{B}G)^{\op} = \mathcal{B}G^{\op}$ où $G^{\op}$ est le groupe opposé (même ensemble, multiplication $a \cdot^{\op} b = b \cdot a$). Pour les groupes abéliens, $G^{\op} \cong G$.
  \item $\Vect_\K^{\op}$ n'est pas « concrètement » la catégorie des espaces vectoriels, mais elle lui est équivalente (via le dual).
\end{enumerate}
\end{exemple}

\section{Le principe de dualité}

\begin{theoreme}[Principe de dualité]\label{thm:dualite}
\index{dualité!principe de}
Soit $\Phi$ un énoncé formulé dans le langage de la théorie des catégories (en termes d'objets, morphismes, composition et identités). Soit $\Phi^{\op}$ l'énoncé obtenu en renversant toutes les flèches (en échangeant source et but, en inversant l'ordre de composition). Alors :
\begin{center}
  $\Phi$ est vrai dans toute catégorie $\iff$ $\Phi^{\op}$ est vrai dans toute catégorie.
\end{center}
Plus précisément, $\Phi$ est vrai dans $\mathcal{C}$ si et seulement si $\Phi^{\op}$ est vrai dans $\mathcal{C}^{\op}$.
\end{theoreme}

\begin{proof}
Si $\Phi$ est vrai dans toute catégorie, alors en particulier $\Phi$ est vrai dans $\mathcal{C}^{\op}$ pour toute catégorie $\mathcal{C}$. Or l'énoncé « $\Phi$ dans $\mathcal{C}^{\op}$ » est exactement « $\Phi^{\op}$ dans $\mathcal{C}$ » (car passer à l'opposée renverse les flèches). Donc $\Phi^{\op}$ est vrai dans toute catégorie $\mathcal{C}$.
\end{proof}

\begin{remarque}
Le principe de dualité est un \emph{méta-théorème} : il porte sur les énoncés de la théorie et non sur des objets mathématiques particuliers. Son utilité pratique est considérable : tout théorème catégorique donne \emph{gratuitement} un second théorème par dualisation.
\end{remarque}

\begin{corollaire}\label{cor:dual-concepts}
Voici un dictionnaire des concepts duaux :
\begin{center}
\begin{tabular}{ll}
\toprule
\textbf{Concept} & \textbf{Concept dual} \\
\midrule
monomorphisme & épimorphisme \\
objet initial & objet terminal \\
section (mono scindé) & rétraction (épi scindé) \\
noyau & conoyau \\
produit & coproduit \\
limite & colimite \\
égaliseur & coégaliseur \\
\bottomrule
\end{tabular}
\end{center}
\end{corollaire}

\section{Exemples de dualisation}

\begin{exemple}[Dualisation des propositions du chapitre 1]\label{ex:dual-chap1}
\index{dualité!exemples}
\begin{enumerate}[label=(\alph*)]
  \item « La composée de monomorphismes est un monomorphisme » dualise en « la composée d'épimorphismes est un épimorphisme ».
  \item « Si $g \circ f$ est mono alors $f$ est mono » dualise en « si $g \circ f$ est épi alors $g$ est épi ».
  \item « L'objet initial est unique à isomorphisme unique près » dualise en la même affirmation pour l'objet terminal.
\end{enumerate}
\end{exemple}

\begin{proposition}\label{prop:contra-cop}
\index{foncteur!contravariant!et catégorie opposée}
Un foncteur contravariant $F : \mathcal{C} \to \mathcal{D}$ est la même chose qu'un foncteur (covariant) $F : \mathcal{C}^{\op} \to \mathcal{D}$, ou de façon équivalente un foncteur $F : \mathcal{C} \to \mathcal{D}^{\op}$.
\end{proposition}

\begin{proof}
La donnée d'une application $F : \Hom_\mathcal{C}(A,B) \to \Hom_\mathcal{D}(F(B), F(A))$ vérifiant $F(g \circ f) = F(f) \circ F(g)$ est exactement la donnée d'une application $\Hom_{\mathcal{C}^{\op}}(B,A) \to \Hom_\mathcal{D}(F(B), F(A))$ vérifiant la fonctorialité covariante (en utilisant la composition dans $\mathcal{C}^{\op}$).
\end{proof}

\section{Dualités concrètes}

Les dualités abstraites de la théorie des catégories se manifestent dans des théorèmes classiques profonds.

\subsection{Dualité de Pontryagin}

\begin{theoreme}[Pontryagin]\label{thm:pontryagin}
\index{dualité!de Pontryagin}\index{Pontryagin}
Soit $\mathbf{LCA}$ la catégorie des groupes abéliens localement compacts et homomorphismes continus. Le foncteur \textbf{dual de Pontryagin}
\[
  \widehat{(-)} : \mathbf{LCA} \longrightarrow \mathbf{LCA}^{\op}, \quad G \longmapsto \widehat{G} = \Hom_{\mathrm{cont}}(G, \R/\Z)
\]
(muni de la topologie compacte-ouverte) définit une équivalence de catégories $\mathbf{LCA} \simeq \mathbf{LCA}^{\op}$.

En particulier, la transformation naturelle canonique $G \to \widehat{\widehat{G}}$ est un isomorphisme naturel.
\end{theoreme}

\begin{remarque}
La dualité de Pontryagin échange :
\begin{itemize}
  \item groupes compacts $\longleftrightarrow$ groupes discrets ;
  \item $\Z \longleftrightarrow \R/\Z$ (le cercle) ;
  \item $\Z/n\Z \longleftrightarrow \Z/n\Z$ (auto-dual).
\end{itemize}
\end{remarque}

\subsection{Dualité de Stone}

\begin{theoreme}[Stone]\label{thm:stone}
\index{dualité!de Stone}\index{Stone}
Il existe une équivalence de catégories
\[
  \mathbf{Bool}^{\op} \simeq \mathbf{Stone}
\]
où $\mathbf{Bool}$ est la catégorie des algèbres de Boole et $\mathbf{Stone}$ est la catégorie des espaces de Stone (espaces topologiques compacts totalement discontinus).

Le foncteur $\mathbf{Bool} \to \mathbf{Stone}^{\op}$ associe à une algèbre de Boole $B$ son espace de spectres $\mathrm{Spec}(B)$ (ensemble des ultrafiltres muni de la topologie de Zariski). Le foncteur inverse associe à un espace de Stone $X$ l'algèbre $\mathrm{Clop}(X)$ de ses parties clopen (ouvertes et fermées).
\end{theoreme}

\begin{remarque}
La dualité de Stone peut être vue comme un précurseur de la dualité de Gelfand (entre $C^*$-algèbres commutatives et espaces compacts) et plus généralement des dualités en géométrie algébrique (le foncteur $\mathrm{Spec}$).
\end{remarque}

\section{Le foncteur Hom}

\begin{definition}[Foncteur Hom covariant]\label{def:hom-covariant}
\index{foncteur!Hom!covariant}\index{Hom@$\Hom$}
Soit $\mathcal{C}$ une catégorie localement petite et $A \in \Ob(\mathcal{C})$. Le \textbf{foncteur Hom covariant} est
\[
  \Hom_\mathcal{C}(A, -) : \mathcal{C} \longrightarrow \Set
\]
défini par :
\begin{itemize}
  \item sur les objets : $B \mapsto \Hom_\mathcal{C}(A, B)$ ;
  \item sur les morphismes : à $f : B \to C$, on associe la post-composition
  \[
    f_* = \Hom_\mathcal{C}(A, f) : \Hom_\mathcal{C}(A, B) \to \Hom_\mathcal{C}(A, C), \quad g \mapsto f \circ g.
  \]
\end{itemize}
\end{definition}

\begin{proposition}\label{prop:hom-cov-foncteur}
$\Hom_\mathcal{C}(A, -)$ est bien un foncteur.
\end{proposition}

\begin{proof}
$(\id_B)_*(g) = \id_B \circ g = g$, donc $(\id_B)_* = \id_{\Hom(A,B)}$. Pour la composition : $(f' \circ f)_*(g) = (f' \circ f) \circ g = f' \circ (f \circ g) = f'_*(f_*(g))$.
\end{proof}

\begin{definition}[Foncteur Hom contravariant]\label{def:hom-contravariant}
\index{foncteur!Hom!contravariant}
Soit $B \in \Ob(\mathcal{C})$. Le \textbf{foncteur Hom contravariant} est
\[
  \Hom_\mathcal{C}(-, B) : \mathcal{C}^{\op} \longrightarrow \Set
\]
défini par :
\begin{itemize}
  \item sur les objets : $A \mapsto \Hom_\mathcal{C}(A, B)$ ;
  \item sur les morphismes : à $f : A' \to A$ dans $\mathcal{C}$, on associe la pré-composition
  \[
    f^* = \Hom_\mathcal{C}(f, B) : \Hom_\mathcal{C}(A, B) \to \Hom_\mathcal{C}(A', B), \quad g \mapsto g \circ f.
  \]
\end{itemize}
\end{definition}

\begin{definition}[Bifoncteur Hom]\label{def:hom-bifoncteur}
\index{foncteur!Hom!bifoncteur}\index{bifoncteur}
Le \textbf{bifoncteur Hom} est le foncteur
\[
  \Hom_\mathcal{C}(-, -) : \mathcal{C}^{\op} \times \mathcal{C} \longrightarrow \Set
\]
défini par $(A, B) \mapsto \Hom_\mathcal{C}(A, B)$ et $(f, g) \mapsto (h \mapsto g \circ h \circ f)$ pour $f : A' \to A$ et $g : B \to B'$.
\end{definition}

\begin{proposition}\label{prop:hom-bifoncteur}
$\Hom_\mathcal{C}(-, -)$ est bien un foncteur $\mathcal{C}^{\op} \times \mathcal{C} \to \Set$.
\end{proposition}

\begin{proof}
Il faut vérifier les deux axiomes. Pour l'identité : $\Hom(\id_A, \id_B)(h) = \id_B \circ h \circ \id_A = h$. Pour la composition : si $(f_1, g_1)$ et $(f_2, g_2)$ sont composables, alors
\begin{align*}
  \Hom(f_1 \circ f_2, g_2 \circ g_1)(h) &= (g_2 \circ g_1) \circ h \circ (f_1 \circ f_2) \\
  &= g_2 \circ (g_1 \circ h \circ f_1) \circ f_2 \\
  &= \Hom(f_2, g_2)(\Hom(f_1, g_1)(h)).
\end{align*}
(Attention à l'ordre dans $\mathcal{C}^{\op}$ : la composition de $(f_2^{\op}, g_2)$ après $(f_1^{\op}, g_1)$ en tant que morphisme de $\mathcal{C}^{\op} \times \mathcal{C}$ donne $((f_1 \circ f_2)^{\op}, g_2 \circ g_1)$.)
\end{proof}

\begin{proposition}\label{prop:hom-mono-epi}
\index{monomorphisme!caractérisation par Hom}\index{épimorphisme!caractérisation par Hom}
Soit $\mathcal{C}$ localement petite et $f : A \to B$ un morphisme.
\begin{enumerate}[label=(\roman*)]
  \item $f$ est mono $\iff$ $f_* : \Hom(C, A) \to \Hom(C, B)$ est injective pour tout objet $C$.
  \item $f$ est épi $\iff$ $f^* : \Hom(B, C) \to \Hom(A, C)$ est injective pour tout objet $C$.
\end{enumerate}
\end{proposition}

\begin{proof}
Ce sont exactement les reformulations des définitions de mono et épi en termes des applications induites.
\end{proof}

\begin{proposition}\label{prop:hom-transfo-nat}
\index{transformation naturelle!entre foncteurs Hom}
Un morphisme $f : A \to B$ dans $\mathcal{C}$ induit une transformation naturelle
\[
  f^* : \Hom_\mathcal{C}(B, -) \Longrightarrow \Hom_\mathcal{C}(A, -)
\]
définie par $f^*_C(g) = g \circ f$ pour $g : B \to C$.
\end{proposition}

\begin{proof}
Pour $h : C \to D$, on vérifie la naturalité : $h_* \circ f^*_C = f^*_D \circ h_*$. En effet, pour $g : B \to C$ :
$(h_* \circ f^*_C)(g) = h \circ (g \circ f) = (h \circ g) \circ f = f^*_D(h \circ g) = (f^*_D \circ h_*)(g)$.
\end{proof}

\section{Exercices}

\begin{exercice}[Catégorie opposée de la catégorie opposée]\label{exo:cop-cop}
Vérifier en détail que $(\mathcal{C}^{\op})^{\op} = \mathcal{C}$ (égalité et non simple isomorphisme).
\end{exercice}

\begin{exercice}[Dual d'un foncteur]\label{exo:dual-foncteur}
Soit $F : \mathcal{C} \to \mathcal{D}$ un foncteur. Montrer que $F$ induit un foncteur $F^{\op} : \mathcal{C}^{\op} \to \mathcal{D}^{\op}$ défini par $F^{\op}(A) = F(A)$ et $F^{\op}(f^{\op}) = (F(f))^{\op}$. Vérifier les axiomes.
\end{exercice}

\begin{exercice}[Pleinement fidèle et isomorphismes]\label{exo:pf-iso}
\index{foncteur!pleinement fidèle!et isomorphismes}
Soit $F : \mathcal{C} \to \mathcal{D}$ un foncteur pleinement fidèle. Montrer que $F$ reflète les isomorphismes : si $F(f)$ est un isomorphisme dans $\mathcal{D}$, alors $f$ est un isomorphisme dans $\mathcal{C}$.
\end{exercice}

\begin{exercice}[Caractérisation des mono dans $\Top$]\label{exo:mono-top}
\index{monomorphisme!dans $\Top$}
\begin{enumerate}[label=(\alph*)]
  \item Montrer que dans $\Top$, les monomorphismes sont exactement les applications continues injectives.
  \item Montrer que dans $\Top$, les épimorphismes sont exactement les applications continues surjectives.
  \item En déduire qu'il existe dans $\Top$ des morphismes qui sont à la fois mono et épi sans être des isomorphismes. Donner un exemple explicite.
\end{enumerate}
\end{exercice}

\begin{exercice}[Hom et produit]\label{exo:hom-produit}
Soit $\mathcal{C}$ une catégorie localement petite possédant des produits binaires. Montrer qu'il existe un isomorphisme naturel (en $C$)
\[
  \Hom_\mathcal{C}(C, A \times B) \cong \Hom_\mathcal{C}(C, A) \times \Hom_\mathcal{C}(C, B).
\]
\end{exercice}

\begin{exercice}[Dualité de Gelfand -- énoncé catégorique]\label{exo:gelfand}
\index{dualité!de Gelfand}
Soit $\mathbf{CHaus}$ la catégorie des espaces compacts de Hausdorff et $\mathbf{cC^*Alg}$ la catégorie des $C^*$-algèbres commutatives unitaires.
\begin{enumerate}[label=(\alph*)]
  \item Décrire le foncteur $C : \mathbf{CHaus} \to \mathbf{cC^*Alg}^{\op}$ qui à $X$ associe $C(X, \C)$.
  \item Énoncer le théorème de Gelfand--Naimark en termes d'équivalence de catégories.
  \item En quoi cette dualité généralise-t-elle la dualité de Stone ?
\end{enumerate}
\end{exercice}

\begin{exercice}[Foncteur Hom interne]\label{exo:hom-interne}
\index{Hom interne}
Dans $\Vect_\K$, montrer que $\Hom_\K(V, W)$ est lui-même un $\K$-espace vectoriel, de sorte que le bifoncteur $\Hom$ se factorise par $\Vect_\K$ plutôt que par $\Set$. On dit que $\Vect_\K$ est \textbf{enrichie sur elle-même}. Est-ce le cas pour $\Grp$ ?
\end{exercice}

\begin{exercice}[Dualité dans les ensembles ordonnés]\label{exo:dual-poset}
Soit $(P, \leq)$ un ensemble ordonné fini.
\begin{enumerate}[label=(\alph*)]
  \item Montrer que les éléments minimaux de $P$ correspondent aux objets initiaux de $\mathcal{C}_P$, et les éléments maximaux aux objets terminaux.
  \item En déduire que la catégorie $\mathcal{C}_P$ admet un objet initial si et seulement si $P$ a un plus petit élément.
  \item Appliquer le principe de dualité pour obtenir le résultat analogue pour les objets terminaux.
\end{enumerate}
\end{exercice}

\begin{exercice}[Plongement de Yoneda -- première rencontre]\label{exo:yoneda-intro}
\index{Yoneda!plongement}
Soit $\mathcal{C}$ une catégorie localement petite. On définit le \textbf{plongement de Yoneda}\index{plongement de Yoneda}
\[
  \mathsf{y} : \mathcal{C} \longrightarrow [\mathcal{C}^{\op}, \Set], \quad A \longmapsto \Hom_\mathcal{C}(-, A).
\]
\begin{enumerate}[label=(\alph*)]
  \item Montrer que $\mathsf{y}$ est un foncteur (à un morphisme $f : A \to B$, associer la transformation naturelle $f_* : \Hom(-, A) \Rightarrow \Hom(-, B)$).
  \item Montrer que $\mathsf{y}$ est pleinement fidèle. (C'est une partie du \textbf{lemme de Yoneda}, qui sera traité en détail au chapitre suivant.)
\end{enumerate}
\end{exercice}
%%%%%%%%%%%%%%%%%%%%%%%%%%%%%%%%%%%%%%%%%%%%%%%%%%%%%%%%%%%%%
%% THÉORIE DES CATÉGORIES — Chapitres 3–4
%% Limites et Colimites / Adjonctions
%%%%%%%%%%%%%%%%%%%%%%%%%%%%%%%%%%%%%%%%%%%%%%%%%%%%%%%%%%%%%

%%%%%%%%%%%%%%%%%%%%%%%%%%%%%%%%%%%%%%%%%%%%%%%%%%%%%%%%%%%%
%%           CHAPITRE 3 : LIMITES ET COLIMITES            %%
%%%%%%%%%%%%%%%%%%%%%%%%%%%%%%%%%%%%%%%%%%%%%%%%%%%%%%%%%%%%

\chapter{Limites et Colimites}\label{ch:limites-colimites}
\minitoc

\vspace{1em}

Les limites et colimites sont les généralisations catégoriques de
constructions omniprésentes en mathématiques~: produits, intersections,
noyaux, images réciproques, limites projectives et inductives.
Dans ce chapitre, nous développons la théorie de façon systématique,
en partant du langage des diagrammes et des cônes, puis en énonçant
la propriété universelle qui définit une limite. Nous démontrons
le théorème fondamental d'existence --- les produits et égaliseurs
suffisent pour construire toutes les limites finies --- et nous
établissons la relation cruciale entre limites et foncteurs
représentables.

\index{limite|(}
\index{colimite|(}

%% ============================================================
\section{Diagrammes et catégories d'indices}\label{sec:diagrammes}
%% ============================================================

\begin{definition}[Diagramme]\label{def:diagramme}
\index{diagramme}
\index{catégorie d'indices}
\index{foncteur!diagramme}
Soit $\mathcal{C}$ une catégorie et $\mathcal{J}$ une petite catégorie
(la \emph{catégorie d'indices} ou \emph{catégorie de forme}).
Un \textbf{diagramme de forme~$\mathcal{J}$ dans~$\mathcal{C}$} est un foncteur
\[
    F \colon \mathcal{J} \longrightarrow \mathcal{C}.
\]
On note $F_j$ l'image d'un objet $j \in \Ob(\mathcal{J})$
et $F_\alpha \colon F_j \to F_k$ l'image d'un morphisme
$\alpha \colon j \to k$ dans~$\mathcal{J}$.
\end{definition}

\begin{exemple}\label{ex:formes-diagrammes}
\index{catégorie discrète}
\index{paire parallèle}
\index{span}
\index{cospan}
Les petites catégories suivantes apparaissent constamment comme catégories d'indices.
\begin{enumerate}[label=(\roman*)]
\item \textbf{Catégories discrètes.}
    Soit $\mathcal{J}$ une catégorie discrète sur un ensemble~$I$
    (objets~$I$, seuls morphismes~: les identités).
    Un diagramme $F\colon \mathcal{J}\to\mathcal{C}$ est simplement une famille
    d'objets $\{F_i\}_{i\in I}$.

\item \textbf{La paire parallèle.}
    $\mathcal{J} = (\bullet \rightrightarrows \bullet)$,
    i.e.\ deux objets $0,1$ avec deux morphismes non triviaux
    $\alpha,\beta\colon 0\to 1$.
    Un diagramme de cette forme est une paire de morphismes parallèles
    $f,g\colon A\to B$ dans~$\mathcal{C}$.

\item \textbf{Le span.}
    $\mathcal{J} = (\bullet \leftarrow \bullet \rightarrow \bullet)$.
    Un diagramme est un span $A \xleftarrow{f} C \xrightarrow{g} B$.

\item \textbf{Le cospan.}
    $\mathcal{J} = (\bullet \rightarrow \bullet \leftarrow \bullet)$.
    Un diagramme est un cospan $A \xrightarrow{f} C \xleftarrow{g} B$.
\end{enumerate}
\end{exemple}

\begin{remarque}\label{rem:diagramme-notation}
\index{diagramme!comme foncteur}
Un diagramme n'est rien d'autre qu'un foncteur ; le terme
\og{}diagramme\fg{} signale simplement que l'on s'apprête à en
prendre la limite ou la colimite.
\end{remarque}


%% ============================================================
\section{Cônes et cocônes}\label{sec:cones}
%% ============================================================

\begin{definition}[Cône]\label{def:cone}
\index{cône}
\index{cône!sommet}
\index{cône!jambe}
Soit $F\colon \mathcal{J}\to\mathcal{C}$ un diagramme.
Un \textbf{cône sur~$F$} est une paire $(N,\,\{\pi_j\}_{j\in\mathcal{J}})$
formée d'un objet $N\in\mathcal{C}$ (le \emph{sommet} du cône)
et d'une famille de morphismes $\pi_j\colon N\to F_j$ (les \emph{jambes}),
un pour chaque objet $j\in\mathcal{J}$,
tels que pour tout morphisme $\alpha\colon j\to k$ dans~$\mathcal{J}$
le triangle suivant commute~:
\[
\begin{tikzcd}
& N \arrow[dl, "\pi_j"'] \arrow[dr, "\pi_k"] & \\
F_j \arrow[rr, "F_\alpha"'] & & F_k
\end{tikzcd}
\]
Autrement dit, $F_\alpha \circ \pi_j = \pi_k$ pour tout $\alpha\colon j\to k$.
\end{definition}

\begin{definition}[Morphisme de cônes]\label{def:morphisme-cone}
\index{cône!morphisme}
Soient $(N,\pi_j)$ et $(N',\pi'_j)$ deux cônes sur~$F$.
Un \textbf{morphisme de cônes} de $(N,\pi_j)$ vers $(N',\pi'_j)$
est un morphisme $u\colon N\to N'$ dans~$\mathcal{C}$ tel que
$\pi'_j \circ u = \pi_j$ pour tout $j\in\mathcal{J}$~:
\[
\begin{tikzcd}
N \arrow[rr, "u"] \arrow[dr, "\pi_j"'] && N' \arrow[dl, "\pi'_j"] \\
& F_j &
\end{tikzcd}
\]
Les cônes sur~$F$ et leurs morphismes forment une catégorie,
notée $\mathsf{Cone}(F)$.
\index{cône!catégorie des cônes}
\end{definition}

\begin{definition}[Cocône]\label{def:cocone}
\index{cocône}
\index{cocône!sommet}
\index{cocône!jambe}
Dualement, un \textbf{cocône sous~$F$} est une paire
$(N,\,\{\iota_j\}_{j\in\mathcal{J}})$ où $\iota_j\colon F_j\to N$
et $\iota_k \circ F_\alpha = \iota_j$ pour tout $\alpha\colon j\to k$~:
\[
\begin{tikzcd}
F_j \arrow[rr, "F_\alpha"] \arrow[dr, "\iota_j"'] & & F_k \arrow[dl, "\iota_k"] \\
& N &
\end{tikzcd}
\]
Les cocônes sous~$F$ forment une catégorie $\mathsf{Cocone}(F)$.
\end{definition}

\begin{remarque}\label{rem:cone-transfo-nat}
\index{foncteur constant}
\index{transformation naturelle!et cônes}
Soit $\Delta_N \colon \mathcal{J}\to\mathcal{C}$ le foncteur constant
envoyant tout objet sur~$N$ et tout morphisme sur~$\id_N$.
Alors un cône $(N,\pi_j)$ sur~$F$ est exactement une transformation naturelle
$\pi\colon \Delta_N \Rightarrow F$.
Dualement, un cocône est une transformation naturelle $\iota\colon F\Rightarrow\Delta_N$.
\end{remarque}


%% ============================================================
\section{Limites}\label{sec:limites}
%% ============================================================

\begin{definition}[Limite]\label{def:limite}
\index{limite!définition}
\index{cône universel}
\index{limite!propriété universelle}
Soit $F\colon\mathcal{J}\to\mathcal{C}$ un diagramme.
Une \textbf{limite} de~$F$ est un cône $(\varprojlim F,\,\{\pi_j\}_{j})$
qui est \emph{terminal} dans $\mathsf{Cone}(F)$.
Explicitement~:
\begin{enumerate}[label=(\alph*)]
\item $(\varprojlim F,\,\pi_j)$ est un cône sur~$F$ ;
\item pour tout cône $(N,\,\psi_j)$ sur~$F$, il existe un
      \emph{unique} morphisme $u\colon N\to\varprojlim F$
      tel que $\pi_j\circ u = \psi_j$ pour tout $j\in\mathcal{J}$~:
\end{enumerate}
\[
\begin{tikzcd}
N \arrow[dr, "\psi_j"'] \arrow[rr, dashed, "\exists!\, u"] & & \varprojlim F \arrow[dl, "\pi_j"] \\
& F_j &
\end{tikzcd}
\]
On appelle les $\pi_j$ les \emph{projections canoniques}\index{projection canonique}.
\end{definition}

\begin{proposition}\label{prop:limite-unique}
\index{limite!unicité}
Si elle existe, la limite est unique à unique isomorphisme près.
\end{proposition}

\begin{proof}
Soient $(L,\pi_j)$ et $(L',\pi'_j)$ deux limites de~$F$.
Par la propriété universelle de~$L$, il existe un unique $u\colon L'\to L$
avec $\pi_j\circ u = \pi'_j$. Par celle de~$L'$, il existe un unique
$v\colon L\to L'$ avec $\pi'_j\circ v = \pi_j$.
Alors $u\circ v\colon L\to L$ satisfait $\pi_j\circ(u\circ v)=\pi_j$,
et par unicité $u\circ v = \id_L$. De même $v\circ u = \id_{L'}$.
\end{proof}


%% ============================================================
\section{Colimites}\label{sec:colimites}
%% ============================================================

\begin{definition}[Colimite]\label{def:colimite}
\index{colimite!définition}
\index{cocône universel}
Soit $F\colon\mathcal{J}\to\mathcal{C}$ un diagramme.
Une \textbf{colimite} de~$F$ est un cocône $(\varinjlim F,\,\{\iota_j\}_{j})$
qui est \emph{initial} dans $\mathsf{Cocone}(F)$.
Pour tout cocône $(N,\,\psi_j)$ sous~$F$, il existe un unique
$u\colon \varinjlim F\to N$ tel que $u\circ\iota_j = \psi_j$
pour tout $j$~:
\[
\begin{tikzcd}
& \varinjlim F \arrow[dd, dashed, "\exists!\, u"] & \\
F_j \arrow[ur, "\iota_j"] \arrow[dr, "\psi_j"'] & & F_k \arrow[ul, "\iota_k"'] \arrow[dl, "\psi_k"] \\
& N &
\end{tikzcd}
\]
Les $\iota_j$ sont les \emph{injections canoniques}\index{injection canonique}.
\end{definition}

\begin{remarque}\label{rem:colimite-dualite}
\index{dualité!limite-colimite}
La colimite de $F\colon\mathcal{J}\to\mathcal{C}$ est la limite de
$F^{\op}\colon\mathcal{J}^{\op}\to\mathcal{C}^{\op}$.
Tout énoncé sur les limites se dualise en un énoncé sur les colimites.
\end{remarque}


%% ============================================================
\section{Produits et coproduits}\label{sec:produits-coproduits}
%% ============================================================

\begin{definition}[Produit]\label{def:produit}
\index{produit}
\index{produit!projection}
Soit $\{A_i\}_{i\in I}$ une famille d'objets de~$\mathcal{C}$
(i.e.\ un diagramme de forme discrète~$I$).
Le \textbf{produit} $\prod_{i\in I}A_i$ est la limite de ce diagramme.
Concrètement, c'est un objet $P$ muni de projections
$\pi_i\colon P\to A_i$ tel que, pour tout objet~$X$ et toute famille
de morphismes $f_i\colon X\to A_i$, il existe un unique
$\langle f_i\rangle\colon X\to P$ avec $\pi_i\circ\langle f_i\rangle = f_i$~:
\[
\begin{tikzcd}
& X \arrow[dl, "f_i"'] \arrow[d, dashed, "{\langle f_i\rangle}"] \arrow[dr, "f_j"] & \\
A_i & \prod_{i\in I}A_i \arrow[l, "\pi_i"] \arrow[r, "\pi_j"'] & A_j
\end{tikzcd}
\]
\end{definition}

\begin{definition}[Coproduit]\label{def:coproduit}
\index{coproduit}
\index{coproduit!injection}
Dualement, le \textbf{coproduit} $\coprod_{i\in I}A_i$ est la colimite
du diagramme discret $\{A_i\}_{i\in I}$. C'est un objet~$S$ muni
d'injections $\iota_i\colon A_i\to S$ tel que, pour tout objet~$X$ et
toute famille $f_i\colon A_i\to X$, il existe un unique
$[f_i]\colon S\to X$ avec $[f_i]\circ\iota_i = f_i$.
\end{definition}

\begin{exemple}\label{ex:produits-concrets}
\index{produit!dans $\Set$}
\index{coproduit!dans $\Set$}
\index{produit!dans $\Grp$}
\index{coproduit!dans $\Grp$}
\begin{enumerate}[label=(\roman*)]
\item Dans $\Set$, le produit est le produit cartésien et le coproduit est
      l'union disjointe.
\item Dans $\Grp$, le produit est le produit direct et le coproduit est
      le produit libre.
\item Dans $\Ab$, le produit est le produit direct et le coproduit
      (fini) est la somme directe $\bigoplus_{i\in I}A_i$.
\item Dans $\Top$, le produit est le produit muni de la topologie produit
      et le coproduit est l'union disjointe avec la topologie somme.
\item Dans $\Vect_\K$, produit et coproduit finis coïncident avec
      la somme directe.
\end{enumerate}
\end{exemple}


%% ============================================================
\section{Égaliseurs et coégaliseurs}\label{sec:egaliseurs}
%% ============================================================

\begin{definition}[Égaliseur]\label{def:egaliseur}
\index{égaliseur}
Soient $f,g\colon A\rightrightarrows B$ deux morphismes parallèles.
L'\textbf{égaliseur} de~$f$ et~$g$ est la limite du diagramme
$\mathcal{J}=(\bullet\rightrightarrows\bullet)$, i.e.\ un objet~$E$
muni d'un morphisme $e\colon E\to A$ tel que $f\circ e = g\circ e$,
universel pour cette propriété~:
\[
\begin{tikzcd}
E \arrow[r, "e"] & A \arrow[r, shift left, "f"] \arrow[r, shift right, "g"'] & B
\end{tikzcd}
\]
Pour tout $h\colon X\to A$ avec $f\circ h = g\circ h$,
il existe un unique $\bar{h}\colon X\to E$ avec $e\circ\bar{h}=h$.
\end{definition}

\begin{definition}[Coégaliseur]\label{def:coegaliseur}
\index{coégaliseur}
Le \textbf{coégaliseur} de $f,g\colon A\rightrightarrows B$ est la colimite
du même diagramme~: un objet~$Q$ muni de $q\colon B\to Q$
tel que $q\circ f = q\circ g$, universel pour cette propriété~:
\[
\begin{tikzcd}
A \arrow[r, shift left, "f"] \arrow[r, shift right, "g"'] & B \arrow[r, "q"] & Q
\end{tikzcd}
\]
\end{definition}

\begin{exemple}\label{ex:egaliseurs-concrets}
\index{égaliseur!dans $\Set$}
\index{coégaliseur!dans $\Set$}
\begin{enumerate}[label=(\roman*)]
\item Dans $\Set$, l'égaliseur de $f,g\colon A\rightrightarrows B$ est
      le sous-ensemble $\{a\in A : f(a)=g(a)\}$ avec l'inclusion.
\item Dans $\Set$, le coégaliseur est le quotient $B/{\sim}$ où $\sim$
      est la plus petite relation d'équivalence telle que $f(a)\sim g(a)$
      pour tout $a\in A$.
\item Dans $\Ab$, l'égaliseur de $f,g$ est le noyau de $f-g$.
\end{enumerate}
\end{exemple}


%% ============================================================
\section{Noyaux et conoyaux}\label{sec:noyaux-conoyaux}
%% ============================================================

\begin{definition}[Noyau]\label{def:noyau}
\index{noyau!catégorique}
Dans une catégorie avec objet nul~$0$\index{objet nul}, le \textbf{noyau}
d'un morphisme $f\colon A\to B$ est l'égaliseur de~$f$ et du morphisme
nul $0_{AB}\colon A\to 0\to B$~:
\[
\begin{tikzcd}
\Ker(f) \arrow[r, hookrightarrow, "k"] & A \arrow[r, shift left, "f"] \arrow[r, shift right, "0"'] & B
\end{tikzcd}
\]
\end{definition}

\begin{definition}[Conoyau]\label{def:conoyau}
\index{conoyau}
Le \textbf{conoyau} de $f\colon A\to B$ est le coégaliseur de~$f$ et de~$0_{AB}$~:
\[
\begin{tikzcd}
A \arrow[r, shift left, "f"] \arrow[r, shift right, "0"'] & B \arrow[r, twoheadrightarrow, "c"] & \coker(f)
\end{tikzcd}
\]
\end{definition}

\begin{exemple}\label{ex:noyaux-concrets}
\index{noyau!dans $\Ab$}
\index{noyau!dans $\Mod$}
Dans $\Ab$ ou $\Mod_R$, le noyau et le conoyau catégoriques coïncident
avec le noyau et le conoyau algébriques habituels.
\end{exemple}


%% ============================================================
\section{Produits fibrés et sommes amalgamées}\label{sec:pullbacks-pushouts}
%% ============================================================

\begin{definition}[Produit fibré]\label{def:produit-fibre}
\index{produit fibré}
\index{pullback|see{produit fibré}}
Soit un cospan $A \xrightarrow{f} C \xleftarrow{g} B$.
Le \textbf{produit fibré} (ou \emph{pullback}) est la limite de ce diagramme.
C'est un objet $A\times_C B$ muni de $p\colon A\times_C B\to A$ et
$q\colon A\times_C B\to B$ tels que $f\circ p = g\circ q$,
universel pour cette propriété~:
\[
\begin{tikzcd}
A\times_C B \arrow[r, "q"] \arrow[d, "p"'] \arrow[dr, phantom, "\lrcorner", very near start]
    & B \arrow[d, "g"] \\
A \arrow[r, "f"'] & C
\end{tikzcd}
\]
Le carré ci-dessus est appelé \emph{carré cartésien}\index{carré cartésien}.
\end{definition}

\begin{definition}[Somme amalgamée]\label{def:somme-amalgamee}
\index{somme amalgamée}
\index{pushout|see{somme amalgamée}}
Soit un span $A \xleftarrow{f} C \xrightarrow{g} B$.
La \textbf{somme amalgamée} (ou \emph{pushout}) est la colimite de ce diagramme~:
un objet $A\sqcup_C B$ muni de $i\colon A\to A\sqcup_C B$ et
$j\colon B\to A\sqcup_C B$ tels que $i\circ f = j\circ g$~:
\[
\begin{tikzcd}
C \arrow[r, "g"] \arrow[d, "f"'] & B \arrow[d, "j"] \\
A \arrow[r, "i"'] & A\sqcup_C B \arrow[ul, phantom, "\ulcorner", very near start]
\end{tikzcd}
\]
\end{definition}

\begin{exemple}\label{ex:pullback-concrets}
\index{produit fibré!dans $\Set$}
\index{somme amalgamée!dans $\Set$}
\index{produit fibré!dans $\Top$}
\begin{enumerate}[label=(\roman*)]
\item Dans $\Set$, $A\times_C B = \{(a,b)\in A\times B : f(a)=g(b)\}$.
\item Dans $\Set$, $A\sqcup_C B = (A\sqcup B)/{\sim}$ où $f(c)\sim g(c)$
      pour tout $c\in C$.
\item Dans $\Top$, le produit fibré est le sous-espace correspondant du produit.
\item Dans $\Grp$, la somme amalgamée est le produit amalgamé de groupes.
\end{enumerate}
\end{exemple}

\begin{proposition}\label{prop:pullback-egaliseur-produit}
\index{produit fibré!via produit et égaliseur}
Si $\mathcal{C}$ admet des produits binaires et des égaliseurs, alors
$\mathcal{C}$ admet des produits fibrés. Plus précisément, le produit fibré
de $A\xrightarrow{f}C\xleftarrow{g}B$ est l'égaliseur de
\[
\begin{tikzcd}
A\times B \arrow[r, shift left, "f\circ\pi_A"] \arrow[r, shift right, "g\circ\pi_B"'] & C
\end{tikzcd}
\]
\end{proposition}

\begin{proof}
Soit $e\colon E\to A\times B$ l'égaliseur de $f\circ\pi_A$ et $g\circ\pi_B$.
Posons $p=\pi_A\circ e$ et $q=\pi_B\circ e$.
Alors $f\circ p = f\circ\pi_A\circ e = g\circ\pi_B\circ e = g\circ q$,
donc $(E,p,q)$ est un cône sur le cospan.
Soit $(X,p',q')$ un autre cône, i.e.\ $f\circ p'=g\circ q'$.
Par la propriété universelle du produit, il existe un unique
$h=\langle p',q'\rangle\colon X\to A\times B$ avec
$\pi_A\circ h=p'$ et $\pi_B\circ h=q'$.
Alors $f\circ\pi_A\circ h = f\circ p' = g\circ q' = g\circ\pi_B\circ h$,
donc par la propriété universelle de l'égaliseur, $h$ se factorise
uniquement par~$e$.
\end{proof}


%% ============================================================
\section{Théorème d'existence des limites finies}\label{sec:existence-limites}
%% ============================================================

\begin{theoreme}\label{thm:existence-limites-finies}
\index{limite!théorème d'existence}
\index{catégorie finiment complète}
Soit $\mathcal{C}$ une catégorie.
Les conditions suivantes sont équivalentes~:
\begin{enumerate}[label=(\roman*)]
\item $\mathcal{C}$ admet toutes les limites finies (i.e.\ les limites
      de diagrammes $F\colon\mathcal{J}\to\mathcal{C}$ avec $\mathcal{J}$ finie).
\item $\mathcal{C}$ admet les produits finis et les égaliseurs.
\item $\mathcal{C}$ admet un objet terminal, les produits binaires et
      les égaliseurs.
\end{enumerate}
On dit alors que $\mathcal{C}$ est \textbf{finiment complète}\index{finiment complète}.
\end{theoreme}

\begin{proof}
L'implication (i)$\Rightarrow$(ii) est claire puisque les produits finis
et les égaliseurs sont des limites finies.

(ii)$\Rightarrow$(iii)~: le produit vide est un objet terminal.

(iii)$\Rightarrow$(i)~: soit $F\colon\mathcal{J}\to\mathcal{C}$ un diagramme
avec $\mathcal{J}$ finie, disons $\Ob(\mathcal{J})=\{j_1,\dots,j_n\}$
et $\Mor(\mathcal{J})$ (non identités) $= \{\alpha_1,\dots,\alpha_m\}$
avec $\alpha_k\colon s(k)\to t(k)$.
Formons le produit $P = \prod_{i=1}^n F_{j_i}$ avec projections $\pi_i$,
et le produit $Q = \prod_{k=1}^m F_{t(k)}$ avec projections $\rho_k$.

Définissons deux morphismes $\varphi,\psi\colon P\to Q$ par~:
\[
    \rho_k\circ\varphi = F_{\alpha_k}\circ\pi_{s(k)},
    \qquad
    \rho_k\circ\psi = \pi_{t(k)}.
\]
Ces morphismes existent par la propriété universelle de~$Q$.
Soit $e\colon E\to P$ l'égaliseur de $\varphi$ et~$\psi$.

Montrons que $(E, \{\pi_i\circ e\}_i)$ est une limite de~$F$.
Pour tout $\alpha_k\colon s(k)\to t(k)$~:
\[
    F_{\alpha_k}\circ(\pi_{s(k)}\circ e)
    = \rho_k\circ\varphi\circ e
    = \rho_k\circ\psi\circ e
    = \pi_{t(k)}\circ e,
\]
donc $(E, \pi_i\circ e)$ est bien un cône.
Soit $(N,\psi_i)$ un autre cône. La famille $\psi_i$ induit un unique
$h\colon N\to P$ avec $\pi_i\circ h = \psi_i$.
La condition de cône donne $\varphi\circ h = \psi\circ h$, donc
$h$ se factorise uniquement par~$e$, ce qui conclut.
\end{proof}

\begin{corollaire}\label{cor:existence-colimites-finies}
\index{catégorie finiment cocomplète}
Dualement, $\mathcal{C}$ admet toutes les colimites finies
si et seulement si elle admet les coproduits finis et les coégaliseurs.
On dit alors que $\mathcal{C}$ est \textbf{finiment cocomplète}.
\end{corollaire}


%% ============================================================
\section{$\Hom$ préserve les limites}\label{sec:hom-preserve-limites}
%% ============================================================

\begin{definition}[Foncteur qui préserve les limites]\label{def:preservation-limites}
\index{foncteur!qui préserve les limites}
\index{foncteur continu}
Un foncteur $G\colon\mathcal{C}\to\mathcal{D}$ \textbf{préserve les limites}
de forme~$\mathcal{J}$ si, pour tout diagramme $F\colon\mathcal{J}\to\mathcal{C}$
dont la limite $(L,\pi_j)$ existe, le cône $(GL, G\pi_j)$ est une limite de
$GF\colon\mathcal{J}\to\mathcal{D}$.
Un foncteur qui préserve toutes les petites limites est dit \textbf{continu}\index{continu}.
\end{definition}

\begin{theoreme}\label{thm:hom-preserve-limites}
\index{Hom!préserve les limites}
\index{foncteur représentable!préserve les limites}
Pour tout objet $X\in\mathcal{C}$, le foncteur covariant
\[
    \Hom_{\mathcal{C}}(X,-)\colon\mathcal{C}\longrightarrow\Set
\]
préserve toutes les limites qui existent dans~$\mathcal{C}$.
Autrement dit, si $(L,\pi_j)=\varprojlim F$, alors
\[
    \Hom(X,\varprojlim F) \cong \varprojlim_j\, \Hom(X,F_j)
\]
naturellement en~$X$.
\end{theoreme}

\begin{proof}
Soit $F\colon\mathcal{J}\to\mathcal{C}$ un diagramme avec limite $(L,\pi_j)$.
La limite $\varprojlim_j\Hom(X,F_j)$ dans $\Set$ est
\[
    \bigl\{(h_j)_{j\in\mathcal{J}}\in\textstyle\prod_j\Hom(X,F_j)
    : F_\alpha\circ h_j = h_k \text{ pour tout } \alpha\colon j\to k\bigr\},
\]
c'est-à-dire l'ensemble des cônes de sommet~$X$ sur~$F$.

Définissons $\Phi\colon\Hom(X,L)\to\varprojlim_j\Hom(X,F_j)$ par
$\Phi(u)=(\pi_j\circ u)_j$. Pour tout $u$, la famille $(\pi_j\circ u)_j$
est un cône (car $(\pi_j)$ l'est), donc $\Phi$ est bien définie.

\emph{Injectivité}~: si $\Phi(u)=\Phi(v)$, alors $\pi_j\circ u=\pi_j\circ v$
pour tout~$j$, et par unicité dans la propriété universelle de~$L$, $u=v$.

\emph{Surjectivité}~: soit $(h_j)_j$ un cône de sommet~$X$.
Par la propriété universelle de~$L$, il existe un unique $u\colon X\to L$
avec $\pi_j\circ u = h_j$, i.e.\ $\Phi(u)=(h_j)_j$.

La naturalité en~$X$ est immédiate.
\end{proof}

\begin{corollaire}\label{cor:hom-contravariant-colimites}
\index{Hom!contravariant et colimites}
Le foncteur contravariant $\Hom_{\mathcal{C}}(-,Y)\colon\mathcal{C}^{\op}\to\Set$
transforme les colimites de~$\mathcal{C}$ en limites de~$\Set$~:
\[
    \Hom(\varinjlim F_j,\, Y) \cong \varprojlim_j\, \Hom(F_j,Y).
\]
\end{corollaire}


%% ============================================================
\section{Limites filtrantes}\label{sec:limites-filtrantes}
%% ============================================================

\begin{definition}[Catégorie filtrante]\label{def:categorie-filtrante}
\index{catégorie filtrante}
\index{filtrant}
Une petite catégorie $\mathcal{J}$ est \textbf{filtrante} si~:
\begin{enumerate}[label=(\alph*)]
\item $\mathcal{J}$ est non vide ;
\item pour tous $j,k\in\mathcal{J}$, il existe un objet~$\ell$ et des
      morphismes $j\to\ell$ et $k\to\ell$ ;
\item pour tous morphismes parallèles $f,g\colon j\rightrightarrows k$,
      il existe $h\colon k\to\ell$ tel que $h\circ f = h\circ g$.
\end{enumerate}
\end{definition}

\begin{definition}[Colimite filtrante]\label{def:colimite-filtrante}
\index{colimite!filtrante}
Une \textbf{colimite filtrante} est une colimite de diagramme
$F\colon\mathcal{J}\to\mathcal{C}$ où $\mathcal{J}$ est filtrante.
Dualement, une \textbf{limite cofiltrante}\index{limite!cofiltrante}
est une limite indexée par une catégorie cofiltrante
($\mathcal{J}^{\op}$ filtrante).
\end{definition}

\begin{exemple}\label{ex:colimites-filtrantes}
\index{système inductif}
\begin{enumerate}[label=(\roman*)]
\item Tout ensemble ordonné filtrant $(I,\leq)$ (vu comme catégorie)
      est filtrant. Les colimites indexées par un tel ensemble sont
      les \emph{limites inductives} classiques.
\item $\varinjlim_{n\in\N}\Z/p^n\Z \cong \Z[1/p]/\Z$ dans $\Ab$.
\end{enumerate}
\end{exemple}

\begin{theoreme}\label{thm:colimites-filtrantes-set}
\index{colimite!filtrante dans $\Set$}
Dans $\Set$, les colimites filtrantes commutent avec les limites finies.
Plus précisément, si $\mathcal{J}$ est filtrante et $\mathcal{K}$ est finie, alors
pour $F\colon\mathcal{J}\times\mathcal{K}\to\Set$~:
\[
    \varinjlim_j\, \varprojlim_k\, F(j,k)
    \cong \varprojlim_k\, \varinjlim_j\, F(j,k).
\]
\end{theoreme}

\begin{proof}[Esquisse de démonstration]
On vérifie directement que la description ensembliste de la colimite
filtrante (classes d'équivalence de paires $(j,x_j)$) permet de
commuter les deux constructions, en utilisant les propriétés (a)--(c)
de la catégorie filtrante pour traiter les cas des produits finis
et des égaliseurs.
\end{proof}


%% ============================================================
\section{Exercices}\label{sec:exo-limites}
%% ============================================================

\begin{exercice}\label{exo:produit-fibre-mono}
\index{monomorphisme!et produit fibré}
Soient $f\colon A\to C$ et $g\colon B\to C$ dans une catégorie
avec produits fibrés. Montrer que si $g$ est un monomorphisme,
alors la projection $p\colon A\times_C B\to A$ est un monomorphisme.
\end{exercice}

\begin{exercice}\label{exo:egaliseur-mono}
\index{égaliseur!est un monomorphisme}
Montrer que tout égaliseur est un monomorphisme.
\end{exercice}

\begin{exercice}\label{exo:limite-fonctorielle}
Soient $F,G\colon\mathcal{J}\to\mathcal{C}$ deux diagrammes et
$\alpha\colon F\Rightarrow G$ une transformation naturelle.
Si $\varprojlim F$ et $\varprojlim G$ existent, montrer qu'il
existe un unique morphisme
$\varprojlim\alpha\colon\varprojlim F\to\varprojlim G$
compatible avec les projections.
\end{exercice}

\begin{exercice}\label{exo:produit-dans-cat}
\index{produit!dans $\Cat$}
Décrire explicitement le produit de deux petites catégories
$\mathcal{C}\times\mathcal{D}$ dans $\Cat$.
\end{exercice}

\begin{exercice}\label{exo:limites-dans-foncteurs}
\index{limite!dans une catégorie de foncteurs}
Soit $\mathcal{C}$ une catégorie et $\mathcal{D}$ une catégorie
admettant toutes les limites de forme~$\mathcal{J}$.
Montrer que $\Fun(\mathcal{C},\mathcal{D})$ admet aussi les limites
de forme~$\mathcal{J}$, calculées \og{}objectivement\fg{}~:
$(\varprojlim F_i)(C) = \varprojlim F_i(C)$.
\end{exercice}

\begin{exercice}\label{exo:complete-cocomplete}
\index{catégorie complète}
Montrer que $\Set$, $\Grp$, $\Ab$, $\Top$ et $\Mod_R$
sont des catégories complètes et cocomplètes
(admettent toutes les petites limites et colimites).
\end{exercice}

\begin{exercice}\label{exo:colimite-filtrante-exacte}
\index{colimite!filtrante!exactitude}
Montrer que dans $\Ab$, les colimites filtrantes sont exactes,
i.e.\ le foncteur $\varinjlim\colon\Fun(\mathcal{J},\Ab)\to\Ab$
(pour $\mathcal{J}$ filtrante) est un foncteur exact.
\end{exercice}

\begin{exercice}\label{exo:pushout-dans-grp}
\index{somme amalgamée!dans $\Grp$}
Soient $H\xrightarrow{f}G_1$ et $H\xrightarrow{g}G_2$ des
homomorphismes de groupes. Montrer que la somme amalgamée
$G_1\sqcup_H G_2$ dans $\Grp$ est le produit amalgamé
$G_1 *_H G_2$.
\end{exercice}

\index{limite|)}
\index{colimite|)}


%%%%%%%%%%%%%%%%%%%%%%%%%%%%%%%%%%%%%%%%%%%%%%%%%%%%%%%%%%%%
%%           CHAPITRE 4 : ADJONCTIONS                     %%
%%%%%%%%%%%%%%%%%%%%%%%%%%%%%%%%%%%%%%%%%%%%%%%%%%%%%%%%%%%%

\chapter{Adjonctions}\label{ch:adjonctions}
\minitoc

\vspace{1em}

Les adjonctions sont l'un des concepts les plus fondamentaux de la théorie
des catégories, reliant deux foncteurs de directions opposées par une
bijection naturelle entre ensembles de morphismes.
Introduites par Daniel Kan en 1958\index{Kan, Daniel},
elles unifient une multitude de constructions mathématiques~:
les foncteurs libres, la tensorisation, l'image directe et réciproque
de faisceaux, les suspensions et espaces de lacets.
Comme le disait Saunders Mac~Lane, \og{}les adjonctions
apparaissent partout\fg{}.

\index{adjonction|(}

%% ============================================================
\section{Définition par bijection naturelle}\label{sec:adjonction-def}
%% ============================================================

\begin{definition}[Adjonction]\label{def:adjonction}
\index{adjonction!définition}
\index{adjoint à gauche}
\index{adjoint à droite}
\index{$F\dashv G$}
Soient $\mathcal{C}$ et $\mathcal{D}$ deux catégories et
\[
\begin{tikzcd}
\mathcal{C} \arrow[r, bend left=20, "F", ""{name=A,below}]
& \mathcal{D} \arrow[l, bend left=20, "G", ""{name=B,above}]
\arrow[from=A, to=B, phantom, "\dashv"]
\end{tikzcd}
\]
deux foncteurs. On dit que $F$ est \textbf{adjoint à gauche} de~$G$
(et $G$ \textbf{adjoint à droite} de~$F$), noté $F\dashv G$,
s'il existe une bijection
\[
    \Phi_{A,B}\colon \Hom_{\mathcal{D}}(FA,\, B)
    \xrightarrow{\;\sim\;}
    \Hom_{\mathcal{C}}(A,\, GB)
\]
naturelle en $A\in\mathcal{C}$ et $B\in\mathcal{D}$.
La naturalité signifie que pour tous $f\colon A'\to A$ dans $\mathcal{C}$
et $g\colon B\to B'$ dans $\mathcal{D}$, le diagramme suivant commute~:
\[
\begin{tikzcd}
\Hom(FA,B) \arrow[r, "\Phi_{A,B}"] \arrow[d, "{(Ff)^*\circ g_*}"']
& \Hom(A,GB) \arrow[d, "{f^*\circ (Gg)_*}"] \\
\Hom(FA',B') \arrow[r, "\Phi_{A',B'}"']
& \Hom(A',GB')
\end{tikzcd}
\]
\end{definition}

\begin{notation}\label{not:adjonction-transposee}
\index{transposée}
\index{morphisme adjoint}
Si $h\colon FA\to B$, on appelle $\tilde{h}=\Phi(h)\colon A\to GB$
la \emph{transposée} (ou \emph{adjointe}) de~$h$. Réciproquement,
si $k\colon A\to GB$, on note $\hat{k}=\Phi^{-1}(k)\colon FA\to B$.
\end{notation}


%% ============================================================
\section{Unité et counité}\label{sec:unite-counite}
%% ============================================================

\begin{definition}[Unité et counité]\label{def:unite-counite}
\index{unité!d'une adjonction}
\index{counité!d'une adjonction}
\index{$\eta$ (unité)}
\index{$\varepsilon$ (counité)}
Soit $F\dashv G$ une adjonction avec bijection~$\Phi$.
\begin{enumerate}[label=(\roman*)]
\item L'\textbf{unité} est la transformation naturelle
$\eta\colon\Id_{\mathcal{C}}\Rightarrow GF$ définie par
\[
    \eta_A = \Phi_{A,FA}(\id_{FA})\colon A\to GFA.
\]
\item La \textbf{counité} est la transformation naturelle
$\varepsilon\colon FG\Rightarrow\Id_{\mathcal{D}}$ définie par
\[
    \varepsilon_B = \Phi^{-1}_{GB,B}(\id_{GB})\colon FGB\to B.
\]
\end{enumerate}
\end{definition}

\begin{proposition}\label{prop:transposee-via-unite}
\index{transposée!via unité et counité}
Pour tout $h\colon FA\to B$ et tout $k\colon A\to GB$~:
\[
    \tilde{h} = Gh\circ\eta_A, \qquad \hat{k} = \varepsilon_B\circ Fk.
\]
\end{proposition}

\begin{proof}
Par naturalité de~$\Phi$ en~$B$, pour $h\colon FA\to B$~:
\[
    \Phi_{A,B}(h)
    = \Phi_{A,B}(h\circ\id_{FA})
    = Gh \circ \Phi_{A,FA}(\id_{FA})
    = Gh\circ\eta_A.
\]
De même, par naturalité en~$A$, pour $k\colon A\to GB$~:
\[
    \Phi^{-1}_{A,B}(k)
    = \Phi^{-1}_{A,B}(\id_{GB}\circ k)
    = \Phi^{-1}_{GB,B}(\id_{GB})\circ Fk
    = \varepsilon_B\circ Fk. \qedhere
\]
\end{proof}


%% ============================================================
\section{Identités triangulaires}\label{sec:identites-triangulaires}
%% ============================================================

\begin{theoreme}[Identités triangulaires]\label{thm:identites-triangulaires}
\index{identités triangulaires}
\index{identités zigzag|see{identités triangulaires}}
Soit $F\dashv G$ avec unité~$\eta$ et counité~$\varepsilon$.
Les compositions suivantes sont des identités~:
\begin{enumerate}[label=(\roman*)]
\item $\varepsilon F\circ F\eta = \id_F$, i.e.\ pour tout $A$~:
\[
\begin{tikzcd}
FA \arrow[r, "F\eta_A"] \arrow[dr, "\id_{FA}"'] & FGFA \arrow[d, "\varepsilon_{FA}"] \\
& FA
\end{tikzcd}
\]
\item $G\varepsilon\circ\eta G = \id_G$, i.e.\ pour tout $B$~:
\[
\begin{tikzcd}
GB \arrow[r, "\eta_{GB}"] \arrow[dr, "\id_{GB}"'] & GFGB \arrow[d, "G\varepsilon_B"] \\
& GB
\end{tikzcd}
\]
\end{enumerate}
\end{theoreme}

\begin{proof}
(i) Soit $A\in\mathcal{C}$. En utilisant la \autoref{prop:transposee-via-unite}
avec $h=\id_{FA}\colon FA\to FA$~:
\[
    \tilde{\id}_{FA} = G(\id_{FA})\circ\eta_A = \eta_A.
\]
Maintenant, appliquons $\Phi^{-1}$ à $\eta_A = \tilde{\id}_{FA}$~:
\[
    \hat{\eta}_A = \varepsilon_{FA}\circ F\eta_A.
\]
Mais $\hat{\eta}_A = \Phi^{-1}(\Phi(\id_{FA})) = \id_{FA}$,
donc $\varepsilon_{FA}\circ F\eta_A = \id_{FA}$.

(ii) Soit $B\in\mathcal{D}$. Avec $k=\id_{GB}\colon GB\to GB$~:
\[
    \hat{\id}_{GB} = \varepsilon_B\circ F(\id_{GB}) = \varepsilon_B.
\]
Appliquons $\Phi$ à $\varepsilon_B = \hat{\id}_{GB}$~:
\[
    \tilde{\varepsilon}_B = G\varepsilon_B\circ\eta_{GB}.
\]
Mais $\tilde{\varepsilon}_B = \Phi(\Phi^{-1}(\id_{GB})) = \id_{GB}$,
donc $G\varepsilon_B\circ\eta_{GB} = \id_{GB}$.
\end{proof}


%% ============================================================
\section{Équivalence des définitions}\label{sec:equivalence-def-adjonctions}
%% ============================================================

\begin{theoreme}\label{thm:equivalence-adjonctions}
\index{adjonction!équivalence des définitions}
Les données suivantes sont équivalentes~:
\begin{enumerate}[label=(\alph*)]
\item Une adjonction $F\dashv G$ donnée par une bijection naturelle
      $\Phi\colon\Hom(FA,B)\xrightarrow{\sim}\Hom(A,GB)$.
\item Un quadruplet $(F,G,\eta,\varepsilon)$ où
      $\eta\colon\Id_\mathcal{C}\Rightarrow GF$ et
      $\varepsilon\colon FG\Rightarrow\Id_\mathcal{D}$
      sont des transformations naturelles satisfaisant les
      identités triangulaires.
\end{enumerate}
\end{theoreme}

\begin{proof}
\emph{(a)$\Rightarrow$(b)}~: c'est ce que nous avons construit dans les
\autoref{def:unite-counite} et \autoref{thm:identites-triangulaires}.

\emph{(b)$\Rightarrow$(a)}~: étant donné $(F,G,\eta,\varepsilon)$
satisfaisant les identités triangulaires, définissons
\[
    \Phi_{A,B}(h) = Gh\circ\eta_A
    \quad\text{pour}\quad h\colon FA\to B,
\]
\[
    \Psi_{A,B}(k) = \varepsilon_B\circ Fk
    \quad\text{pour}\quad k\colon A\to GB.
\]

Vérifions que $\Psi\circ\Phi = \id$.
Soit $h\colon FA\to B$~:
\[
    \Psi(\Phi(h))
    = \varepsilon_B\circ F(Gh\circ\eta_A)
    = \varepsilon_B\circ FGh\circ F\eta_A.
\]
Par naturalité de~$\varepsilon$ (appliquée à $h\colon FA\to B$)~:
$\varepsilon_B\circ FGh = h\circ\varepsilon_{FA}$, donc
\[
    \Psi(\Phi(h))
    = h\circ\varepsilon_{FA}\circ F\eta_A
    = h\circ\id_{FA} = h,
\]
en utilisant la première identité triangulaire.

Vérifions que $\Phi\circ\Psi=\id$.
Soit $k\colon A\to GB$~:
\[
    \Phi(\Psi(k))
    = G(\varepsilon_B\circ Fk)\circ\eta_A
    = G\varepsilon_B\circ GFk\circ\eta_A.
\]
Par naturalité de~$\eta$ (appliquée à $k\colon A\to GB$)~:
$GFk\circ\eta_A = \eta_{GB}\circ k$, donc
\[
    \Phi(\Psi(k))
    = G\varepsilon_B\circ\eta_{GB}\circ k
    = \id_{GB}\circ k = k,
\]
en utilisant la seconde identité triangulaire.

La naturalité de~$\Phi$ en~$A$ et~$B$ suit de la fonctorialité de~$G$
et de la naturalité de~$\eta$.
\end{proof}


%% ============================================================
\section{Exemples fondamentaux}\label{sec:exemples-adjonctions}
%% ============================================================

\begin{exemple}[Libre $\dashv$ Oubli]\label{ex:libre-oubli}
\index{adjonction!libre-oubli}
\index{foncteur libre}
\index{foncteur d'oubli}
Le foncteur d'oubli $U\colon\Grp\to\Set$ admet un adjoint à gauche $F$,
le foncteur \og{}groupe libre\fg{}~:
\[
\begin{tikzcd}
\Set \arrow[r, bend left=20, "F", ""{name=A,below}]
& \Grp \arrow[l, bend left=20, "U", ""{name=B,above}]
\arrow[from=A, to=B, phantom, "\dashv"]
\end{tikzcd}
\qquad
\Hom_{\Grp}(FX,\,G) \cong \Hom_{\Set}(X,\,UG).
\]
L'unité $\eta_X\colon X\to UFX$ est l'inclusion de~$X$ comme
générateurs dans le groupe libre $FX$. La counité
$\varepsilon_G\colon FUG\to G$ envoie chaque mot formel dans
le groupe libre sur les éléments de~$G$ sur son évaluation dans~$G$.

De même pour $\Ab$, $\Ring$, $\Mod_R$, $\Vect_\K$, etc.
\end{exemple}

\begin{exemple}[Produit $\dashv$ Hom interne]\label{ex:produit-hom}
\index{adjonction!produit-hom}
\index{catégorie cartésienne fermée}
\index{exponentiation}
Dans $\Set$ (ou plus généralement dans toute catégorie cartésienne fermée),
pour tout ensemble~$A$~:
\[
    (-)\times A \;\dashv\; (-)^A,
\]
i.e.\ $\Hom(X\times A,\, B)\cong\Hom(X,\, B^A)$ où $B^A=\Hom(A,B)$.
C'est la \emph{curryfication}\index{curryfication}.
\end{exemple}

\begin{exemple}[Image réciproque $\dashv$ image directe]\label{ex:faisceaux}
\index{adjonction!$f^*\dashv f_*$}
\index{image directe}
\index{image réciproque}
Soit $f\colon X\to Y$ une application continue entre espaces topologiques.
Les foncteurs image réciproque $f^*$ et image directe $f_*$ sur les
faisceaux satisfont $f^*\dashv f_*$~:
\[
    \Hom_{\mathsf{Sh}(X)}(f^*\mathcal{G},\,\mathcal{F})
    \cong \Hom_{\mathsf{Sh}(Y)}(\mathcal{G},\,f_*\mathcal{F}).
\]
\end{exemple}

\begin{exemple}[Suspension $\dashv$ Lacets]\label{ex:suspension-lacets}
\index{adjonction!$\Sigma\dashv\Omega$}
\index{suspension}
\index{espace de lacets}
Dans la catégorie homotopique pointée $\Top_*$, la suspension réduite
$\Sigma$ et le foncteur espace de lacets~$\Omega$ satisfont
\[
    \Sigma \dashv \Omega,\qquad
    [\Sigma X,\,Y]_* \cong [X,\,\Omega Y]_*
\]
où $[-,-]_*$ désigne les classes d'homotopie pointée.
\end{exemple}

\begin{exemple}[Foncteur diagonal $\dashv$ Produit]\label{ex:diagonal-produit}
\index{adjonction!diagonal-produit}
\index{foncteur diagonal}
Soit $\Delta\colon\mathcal{C}\to\mathcal{C}\times\mathcal{C}$ le foncteur
diagonal $\Delta(A)=(A,A)$. Si $\mathcal{C}$ admet des produits binaires,
alors $\Delta\dashv\times$~:
\[
    \Hom_{\mathcal{C}\times\mathcal{C}}(\Delta A,\,(B,C))
    = \Hom(A,B)\times\Hom(A,C)
    \cong \Hom(A,\,B\times C).
\]
Dualement, le coproduit est adjoint à gauche du diagonal~:
$\sqcup\dashv\Delta$.
\end{exemple}

\begin{exemple}[Extension et restriction des scalaires]\label{ex:scalaires}
\index{adjonction!extension-restriction des scalaires}
\index{extension des scalaires}
Soit $\varphi\colon R\to S$ un homomorphisme d'anneaux.
Le foncteur de restriction des scalaires $\varphi^*\colon\Mod_S\to\Mod_R$
admet un adjoint à gauche, l'extension des scalaires
$S\otimes_R(-)\colon\Mod_R\to\Mod_S$~:
\[
    \Hom_S(S\otimes_R M,\, N) \cong \Hom_R(M,\,\varphi^* N).
\]
\end{exemple}


%% ============================================================
\section{Adjoints et limites}\label{sec:adjoints-limites}
%% ============================================================

\begin{theoreme}[Les adjoints à droite préservent les limites]\label{thm:rapl}
\index{adjoint à droite!préserve les limites}
\index{RAPL}
Soit $F\dashv G$ avec $F\colon\mathcal{C}\to\mathcal{D}$ et
$G\colon\mathcal{D}\to\mathcal{C}$.
Alors $G$ préserve toutes les limites qui existent dans~$\mathcal{D}$.
\end{theoreme}

\begin{proof}
Soit $D\colon\mathcal{J}\to\mathcal{D}$ un diagramme avec limite
$(L,\pi_j)$. Montrons que $(GL,G\pi_j)$ est une limite de
$GD\colon\mathcal{J}\to\mathcal{C}$.

Soit $(X,\psi_j)$ un cône sur~$GD$ dans~$\mathcal{C}$,
i.e.\ $\psi_j\colon X\to GD_j$ avec $GD_\alpha\circ\psi_j=\psi_k$
pour tout $\alpha\colon j\to k$.

Par l'adjonction, chaque $\psi_j\colon X\to GD_j$ correspond à un unique
$\hat{\psi}_j = \varepsilon_{D_j}\circ F\psi_j\colon FX\to D_j$.
Vérifions que $(\hat{\psi}_j)_j$ est un cône sur~$D$.
Pour $\alpha\colon j\to k$~:
\begin{align*}
    D_\alpha\circ\hat{\psi}_j
    &= D_\alpha\circ\varepsilon_{D_j}\circ F\psi_j
    = \varepsilon_{D_k}\circ FGD_\alpha\circ F\psi_j \\
    &= \varepsilon_{D_k}\circ F(GD_\alpha\circ\psi_j)
    = \varepsilon_{D_k}\circ F\psi_k = \hat{\psi}_k,
\end{align*}
où l'on a utilisé la naturalité de~$\varepsilon$.

Par la propriété universelle de~$(L,\pi_j)$, il existe un unique
$v\colon FX\to L$ avec $\pi_j\circ v = \hat{\psi}_j$.
En transposant, on obtient $u = \tilde{v} = Gv\circ\eta_X\colon X\to GL$.
Alors~:
\[
    G\pi_j\circ u = G\pi_j\circ Gv\circ\eta_X
    = G(\pi_j\circ v)\circ\eta_X
    = G\hat{\psi}_j\circ\eta_X
    = \widetilde{\hat{\psi}}_j = \psi_j.
\]
L'unicité de~$u$ découle de celle de~$v$ et de la bijectivité de la transposition.
\end{proof}

\begin{theoreme}[Les adjoints à gauche préservent les colimites]\label{thm:lapc}
\index{adjoint à gauche!préserve les colimites}
\index{LAPC}
Soit $F\dashv G$ avec $F\colon\mathcal{C}\to\mathcal{D}$ et
$G\colon\mathcal{D}\to\mathcal{C}$.
Alors $F$ préserve toutes les colimites qui existent dans~$\mathcal{C}$.
\end{theoreme}

\begin{proof}
Soit $D\colon\mathcal{J}\to\mathcal{C}$ un diagramme avec colimite
$(C,\iota_j)$. Montrons que $(FC,F\iota_j)$ est une colimite de
$FD\colon\mathcal{J}\to\mathcal{D}$.

Soit $(Y,\psi_j)$ un cocône sous~$FD$, i.e.\ $\psi_j\colon FD_j\to Y$.
Par l'adjonction, chaque $\psi_j$ correspond à
$\tilde{\psi}_j\colon D_j\to GY$.
La famille $(\tilde{\psi}_j)_j$ est un cocône sous~$D$ dans~$\mathcal{C}$
(par naturalité de la bijection).

Par la propriété universelle de~$(C,\iota_j)$, il existe un unique
$w\colon C\to GY$ avec $w\circ\iota_j = \tilde{\psi}_j$.
En transposant, $\hat{w} = \varepsilon_Y\circ Fw\colon FC\to Y$ satisfait~:
\[
    \hat{w}\circ F\iota_j
    = \varepsilon_Y\circ Fw\circ F\iota_j
    = \varepsilon_Y\circ F(w\circ\iota_j)
    = \varepsilon_Y\circ F\tilde{\psi}_j
    = \hat{\tilde{\psi}}_j = \psi_j.
\]
L'unicité suit comme dans la preuve précédente.
\end{proof}

\begin{corollaire}\label{cor:adj-hom-lim}
\index{adjonction!et Hom}
Si $F\dashv G$, alors pour tout diagramme $D\colon\mathcal{J}\to\mathcal{D}$
dont la limite existe~:
\[
    G(\varprojlim D_j) \cong \varprojlim G(D_j),
    \qquad
    F(\varinjlim D_j) \cong \varinjlim F(D_j).
\]
\end{corollaire}


%% ============================================================
\section{Le théorème de Freyd (GAFT)}\label{sec:freyd}
%% ============================================================

Nous établissons maintenant un critère puissant garantissant l'existence
d'un foncteur adjoint.

\begin{definition}[Condition d'ensemble de solutions]\label{def:solution-set}
\index{condition d'ensemble de solutions}
\index{solution set condition}
Un foncteur $G\colon\mathcal{D}\to\mathcal{C}$ satisfait la
\textbf{condition d'ensemble de solutions} si, pour tout $A\in\mathcal{C}$,
il existe un ensemble (et non une classe propre)
$\{f_i\colon A\to GB_i\}_{i\in I_A}$ de morphismes tel que tout morphisme
$f\colon A\to GB$ se factorise comme $f = Gh\circ f_i$ pour un certain
$i\in I_A$ et un certain $h\colon B_i\to B$.
\[
\begin{tikzcd}
A \arrow[r, "f_i"] \arrow[dr, "f"'] & GB_i \arrow[d, "Gh"] \\
& GB
\end{tikzcd}
\]
\end{definition}

\begin{theoreme}[Théorème de l'adjoint de Freyd (GAFT)]\label{thm:gaft}
\index{théorème de Freyd}
\index{GAFT}
\index{théorème de l'adjoint}
Soit $G\colon\mathcal{D}\to\mathcal{C}$ un foncteur où $\mathcal{D}$
est complète (admet toutes les petites limites).
Alors $G$ admet un adjoint à gauche si et seulement si~:
\begin{enumerate}[label=(\roman*)]
\item $G$ préserve toutes les petites limites (est continu) ;
\item $G$ satisfait la condition d'ensemble de solutions.
\end{enumerate}
\end{theoreme}

\begin{proof}
\emph{Nécessité}~: si $F\dashv G$, alors $G$ préserve les limites par
le \autoref{thm:rapl}. Pour la condition d'ensemble de solutions,
l'ensemble $\{\eta_A\colon A\to GFA\}$ (réduit à un seul élément)
convient, car tout $f\colon A\to GB$ a pour transposée $\hat{f}\colon FA\to B$,
et $f = G\hat{f}\circ\eta_A$.

\emph{Suffisance}~: fixons $A\in\mathcal{C}$.
Nous construisons un objet initial dans la catégorie comma $(A\downarrow G)$\index{catégorie comma},
dont les objets sont les paires $(B,f\colon A\to GB)$ et les morphismes
$(B,f)\to(B',f')$ sont les $h\colon B\to B'$ avec $Gh\circ f = f'$.

\emph{Étape 1.} Par la condition d'ensemble de solutions, soit
$\{(B_i,f_i)\}_{i\in I}$ un ensemble de solutions.
Posons $B_0 = \prod_{i\in I}B_i$ (produit dans~$\mathcal{D}$,
qui existe car $\mathcal{D}$ est complète).
Comme $G$ préserve les limites, $GB_0\cong\prod_i GB_i$,
et la famille $(f_i)_i$ induit $f_0\colon A\to GB_0$.
L'objet $(B_0,f_0)$ est \og{}faiblement initial\fg{}~: tout objet
de $(A\downarrow G)$ reçoit au moins un morphisme depuis $(B_0,f_0)$.

\emph{Étape 2.} Pour rendre l'objet initial, nous égalisons.
Considérons l'ensemble de tous les endomorphismes
$h\colon B_0\to B_0$ dans $(A\downarrow G)$
(c'est un ensemble car $\Hom(B_0,B_0)$ est un ensemble).
Formons l'égaliseur $e\colon E\to B_0$ de tous ces endomorphismes
dans~$\mathcal{D}$.
(Plus précisément, l'égaliseur de la famille de paires
$(\id_{B_0}, h)$ pour tout endomorphisme~$h$.)
Comme $G$ préserve les égaliseurs, $Ge\colon GE\to GB_0$ est
l'égaliseur correspondant dans~$\mathcal{C}$.
Puisque $\id_{B_0}$ est un endomorphisme de $(B_0,f_0)$,
on a $Gh\circ f_0 = f_0$ pour tout~$h$, donc
$f_0$ se factorise par $Ge$ en un unique $f_E\colon A\to GE$.

\emph{Étape 3.} Montrons que $(E,f_E)$ est initial dans $(A\downarrow G)$.
Par l'étape~1 (via la factorisation $e$), tout objet $(B,f)$ reçoit un
morphisme depuis $(E,f_E)$~: en effet, le morphisme $(B_0,f_0)\to(B,f)$
se factorise par~$E$.
Pour l'unicité, soient $u,v\colon E\to B$ deux morphismes avec
$Gu\circ f_E = Gv\circ f_E = f$.
Formons l'égaliseur $w\colon W\to E$ de $u$ et~$v$.
Par faible initialité de~$(B_0,f_0)$, il existe $s\colon B_0\to W$.
La composée $e' = e\circ w\circ s\colon B_0\to B_0$
est un endomorphisme de $(B_0,f_0)$, donc $e\circ w\circ s$ est égalisé
par~$e$, ce qui implique $w\circ s\circ e_0 = \id$ (par la propriété
de l'égaliseur en étape~2), et finalement $w$ est un isomorphisme,
donc $u=v$.

Posons alors $FA = E$. La correspondance $A\mapsto FA$ s'étend
en un foncteur adjoint à gauche de~$G$, avec unité $\eta_A = f_E$.
\end{proof}

\begin{corollaire}\label{cor:saft}
\index{SAFT}
\index{théorème de l'adjoint spécial}
\textbf{(Théorème de l'adjoint spécial, SAFT.)}
Soit $G\colon\mathcal{D}\to\mathcal{C}$ un foncteur continu
où $\mathcal{D}$ est complète et localement petite.
Si $\mathcal{D}$ admet une famille cogénératrice\index{famille cogénératrice},
alors $G$ admet un adjoint à gauche.
\end{corollaire}

\begin{proof}[Esquisse]
La famille cogénératrice fournit automatiquement la condition d'ensemble
de solutions, et on applique le \autoref{thm:gaft}.
\end{proof}

\begin{exemple}\label{ex:applications-gaft}
\begin{enumerate}[label=(\roman*)]
\item \index{complétion!de groupe abélien}
      Le foncteur d'oubli $U\colon\Ab\to\Set$ préserve les limites
      et $\Ab$ est complète. L'ensemble $\{f\colon X\to UB : |B|\leq|X|+\aleph_0\}$
      fournit un ensemble de solutions, d'où $F\dashv U$ (groupes abéliens libres).
\item \index{théorème de Stone-Čech}
      Le foncteur d'oubli $U\colon\mathsf{CompHaus}\to\Top$
      satisfait les conditions du SAFT, ce qui donne la
      compactification de Stone--Čech $\beta\dashv U$.
\end{enumerate}
\end{exemple}


%% ============================================================
\section{Adjonctions et catégories comma}\label{sec:adj-comma}
%% ============================================================

\begin{proposition}\label{prop:adj-objet-initial-comma}
\index{adjonction!et catégorie comma}
\index{catégorie comma}
Un foncteur $G\colon\mathcal{D}\to\mathcal{C}$ admet un adjoint à gauche
si et seulement si, pour tout $A\in\mathcal{C}$, la catégorie comma
$(A\downarrow G)$ admet un objet initial.
\end{proposition}

\begin{proof}
Si $F\dashv G$ avec unité~$\eta$, alors $(FA,\eta_A)$ est initial
dans $(A\downarrow G)$~: pour tout $(B,f\colon A\to GB)$, le morphisme
$\hat{f}=\varepsilon_B\circ Ff\colon FA\to B$ est l'unique morphisme
avec $G\hat{f}\circ\eta_A = f$.

Réciproquement, si $(L_A,\eta_A\colon A\to GL_A)$ est initial dans
$(A\downarrow G)$ pour tout~$A$, on pose $FA=L_A$ et on vérifie
la fonctorialité et la bijection naturelle.
\end{proof}


%% ============================================================
\section{Propriétés supplémentaires}\label{sec:adj-proprietes}
%% ============================================================

\begin{proposition}\label{prop:adj-pleinement-fidele}
\index{adjonction!et foncteur pleinement fidèle}
\index{foncteur pleinement fidèle!et adjonction}
Soit $F\dashv G$.
\begin{enumerate}[label=(\roman*)]
\item $G$ est fidèle si et seulement si chaque composante de la
      counité~$\varepsilon_B$ est un épimorphisme.
\item $G$ est pleinement fidèle si et seulement si $\varepsilon$
      est un isomorphisme naturel.
\end{enumerate}
\end{proposition}

\begin{proof}
(i) $G$ fidèle $\Leftrightarrow$ $\Phi$ injective sur les fibres
$\Leftrightarrow$ la composée $h\mapsto\varepsilon_B\circ Fk\mapsto h$
est injective
$\Leftrightarrow$ $\varepsilon_B$ est un épimorphisme (car $\hat{k}=\varepsilon_B\circ Fk$).

(ii) $G$ pleinement fidèle $\Leftrightarrow$ $\Phi$ bijectif pour $A=GB$
$\Leftrightarrow$ $\Hom(FGB,B)\cong\Hom(GB,GB)$ par $h\mapsto Gh\circ\eta_{GB}$,
et $\varepsilon_B$ est l'antécédent de $\id_{GB}$, d'où le résultat.
\end{proof}

\begin{proposition}[Unicité des adjoints]\label{prop:unicite-adjoints}
\index{adjonction!unicité}
Si $F\dashv G$ et $F'\dashv G$, alors $F\cong F'$ (naturellement).
De même, si $F\dashv G$ et $F\dashv G'$, alors $G\cong G'$.
\end{proposition}

\begin{proof}
Par le lemme de Yoneda, $F$ est déterminé par le foncteur
$\Hom(F(-),-)=\Hom(-,G(-))$, qui ne dépend que de~$G$.
\end{proof}

\begin{proposition}[Composition d'adjonctions]\label{prop:composition-adj}
\index{adjonction!composition}
Si $F\dashv G\colon\mathcal{C}\rightleftarrows\mathcal{D}$ et
$F'\dashv G'\colon\mathcal{D}\rightleftarrows\mathcal{E}$, alors
$F'F\dashv GG'$.
\end{proposition}

\begin{proof}
$\Hom_\mathcal{E}(F'FA,B)
\cong\Hom_\mathcal{D}(FA,G'B)
\cong\Hom_\mathcal{C}(A,GG'B)$,
naturellement en $A$ et~$B$.
\end{proof}


%% ============================================================
\section{Monades issues d'adjonctions}\label{sec:monades-apercu}
%% ============================================================

\begin{remarque}\label{rem:monade-apercu}
\index{monade}
Toute adjonction $F\dashv G$ engendre une \emph{monade}
$T = GF\colon\mathcal{C}\to\mathcal{C}$ munie de~:
\begin{itemize}
\item $\eta\colon\Id\Rightarrow T$ (l'unité de l'adjonction),
\item $\mu = G\varepsilon F\colon T^2 = GFGF\Rightarrow GF = T$
      (la multiplication).
\end{itemize}
Les identités triangulaires entraînent les axiomes de monade~:
$\mu\circ T\eta = \mu\circ\eta T = \id_T$ et
$\mu\circ T\mu = \mu\circ\mu T$ (associativité).
Nous développerons cette théorie au chapitre suivant.
\end{remarque}


%% ============================================================
\section{Exercices}\label{sec:exo-adjonctions}
%% ============================================================

\begin{exercice}\label{exo:adj-libre-ab}
\index{adjonction!groupe abélien libre}
Soit $F\colon\Set\to\Ab$ le foncteur \og{}groupe abélien libre\fg{}
($FX = \bigoplus_{x\in X}\Z$) et $U\colon\Ab\to\Set$ l'oubli.
Montrer que $F\dashv U$ en exhibant l'unité, la counité,
et en vérifiant les identités triangulaires.
\end{exercice}

\begin{exercice}\label{exo:adj-tensor-hom}
\index{adjonction!tenseur-hom}
Soit $R$ un anneau commutatif et $M$ un $R$-module.
Montrer que $M\otimes_R(-)\dashv\Hom_R(M,-)$ en établissant
la bijection naturelle
$\Hom_R(M\otimes_R N,\,P)\cong\Hom_R(N,\,\Hom_R(M,P))$.
\end{exercice}

\begin{exercice}\label{exo:adj-preserve-epi}
\index{adjoint à gauche!préserve les épimorphismes}
Montrer que tout adjoint à gauche préserve les épimorphismes
et que tout adjoint à droite préserve les monomorphismes.
\end{exercice}

\begin{exercice}\label{exo:adj-equivalence}
\index{équivalence de catégories!et adjonction}
Montrer qu'une adjonction $F\dashv G$ est une équivalence de catégories
si et seulement si $\eta$ et $\varepsilon$ sont des isomorphismes naturels.
\end{exercice}

\begin{exercice}\label{exo:adj-reflexive}
\index{sous-catégorie réflexive}
Soit $\mathcal{D}\hookrightarrow\mathcal{C}$ une sous-catégorie pleine.
On dit que $\mathcal{D}$ est \emph{réflexive} si le foncteur d'inclusion
$\iota\colon\mathcal{D}\hookrightarrow\mathcal{C}$ admet un adjoint à gauche
$L\colon\mathcal{C}\to\mathcal{D}$ (le \emph{réflecteur}).
\begin{enumerate}[label=(\alph*)]
\item Montrer que $\varepsilon$ est un isomorphisme naturel.
\item Montrer que $L$ est essentiellement surjectif.
\item Montrer que $\Ab$ est une sous-catégorie réflexive de $\Grp$
      (le réflecteur est l'abélianisation $G\mapsto G/[G,G]$).
\end{enumerate}
\end{exercice}

\begin{exercice}\label{exo:gaft-application}
\index{théorème de Freyd!application}
En utilisant le GAFT, montrer que le foncteur d'oubli
$U\colon\mathsf{CompHaus}\to\Top$ admet un adjoint à gauche
(la compactification de Stone--Čech).
\emph{Indication}~: vérifier que $\mathsf{CompHaus}$ est complète,
que $U$ préserve les limites, et exhiber un ensemble de solutions.
\end{exercice}

\begin{exercice}\label{exo:adj-kan}
\index{extension de Kan!à gauche}
Soient $K\colon\mathcal{A}\to\mathcal{B}$ et $F\colon\mathcal{A}\to\mathcal{C}$
des foncteurs avec $\mathcal{A}$ petite et $\mathcal{C}$ cocomplète.
Montrer que l'extension de Kan à gauche $\mathrm{Lan}_K F$
existe et que $\mathrm{Lan}_K\dashv K^*$ où
$K^*\colon\Fun(\mathcal{B},\mathcal{C})\to\Fun(\mathcal{A},\mathcal{C})$
est la précomposition par~$K$.
\end{exercice}

\begin{exercice}\label{exo:adj-monade}
\index{monade!axiomes}
Soit $F\dashv G$ avec $T=GF$, $\eta$ l'unité et $\mu=G\varepsilon F$.
Vérifier en détail les axiomes de monade~:
\begin{enumerate}[label=(\alph*)]
\item $\mu\circ T\eta = \id_T$ et $\mu\circ\eta T = \id_T$
      (axiomes d'unité).
\item $\mu\circ T\mu = \mu\circ\mu T$ (associativité).
\end{enumerate}
\emph{Indication}~: utiliser les identités triangulaires et la
naturalité de $\eta$ et $\varepsilon$.
\end{exercice}

\begin{exercice}\label{exo:adj-galois}
\index{correspondance de Galois}
\index{adjonction!de Galois}
Soient $(P,\leq)$ et $(Q,\leq)$ deux ensembles partiellement ordonnés,
vus comme catégories. Montrer qu'une adjonction $F\dashv G$
entre ces catégories est exactement une \emph{connexion de Galois}~:
$F(p)\leq q \Leftrightarrow p\leq G(q)$.
Donner un exemple en théorie des treillis.
\end{exercice}

\index{adjonction|)}
%%%%%%%%%%%%%%%%%%%%%%%%%%%%%%%%%%%%%%%%%%%%%%%%%%%%%%%%%%%%%
%% THÉORIE DES CATÉGORIES — Chapitres 5–6
%% Lemme de Yoneda et Représentabilité / Monades et Algèbres
%%%%%%%%%%%%%%%%%%%%%%%%%%%%%%%%%%%%%%%%%%%%%%%%%%%%%%%%%%%%%

%%%%%%%%%%%%%%%%%%%%%%%%%%%%%%%%%%%%%%%%%%%%%%%%%%%%%%%%%%%%
%%     CHAPITRE 5 : LEMME DE YONEDA ET REPRÉSENTABILITÉ   %%
%%%%%%%%%%%%%%%%%%%%%%%%%%%%%%%%%%%%%%%%%%%%%%%%%%%%%%%%%%%%

\chapter{Lemme de Yoneda et représentabilité}\label{ch:yoneda}
\minitoc

\vspace{1em}

Le lemme de Yoneda\index{Yoneda, lemme de|(} est, selon beaucoup de catégoriciens,
le résultat le plus fondamental de la théorie des catégories.
Il affirme qu'un objet est entièrement déterminé
par l'ensemble des morphismes qui y arrivent (ou qui en partent),
et que cette détermination est fonctorielle.
Plus précisément, le lemme établit une bijection naturelle entre les
transformations naturelles issues d'un foncteur représentable et les
éléments de l'objet cible, ce qui permet de \emph{plonger} toute catégorie
dans une catégorie de foncteurs.

Ce chapitre développe la théorie des préfaisceaux, démontre le lemme de Yoneda
dans sa forme complète, en déduit le plongement de Yoneda, puis étudie
les foncteurs représentables et leurs propriétés.

\index{préfaisceau|(}
\index{foncteur représentable|(}

%% ============================================================
\section{Préfaisceaux et foncteurs \texorpdfstring{$\Hom$}{Hom}}\label{sec:prefaisceaux}
%% ============================================================

\begin{definition}[Préfaisceau]\label{def:prefaisceau}
\index{préfaisceau!définition}
\index{catégorie!des préfaisceaux}
Soit $\mathcal{C}$ une catégorie localement petite.
Un \textbf{préfaisceau} sur $\mathcal{C}$ est un foncteur
\[
    F \colon \mathcal{C}^{\op} \longrightarrow \Set.
\]
La catégorie des préfaisceaux sur $\mathcal{C}$ est la catégorie de foncteurs
$\Fun(\mathcal{C}^{\op}, \Set)$,
que l'on note $\widehat{\mathcal{C}}$ ou $\mathrm{PSh}(\mathcal{C})$.
\index{$\widehat{\mathcal{C}}$}
\index{PSh@$\mathrm{PSh}(\mathcal{C})$}
Les morphismes dans $\widehat{\mathcal{C}}$ sont les transformations naturelles.
\end{definition}

\begin{remarque}\label{rem:prefaisceau-taille}
\index{préfaisceau!questions de taille}
La catégorie $\widehat{\mathcal{C}}$ n'est en général pas petite, même si
$\mathcal{C}$ est petite : les foncteurs de $\mathcal{C}^{\op}$ vers $\Set$
forment une classe propre dès que $\mathcal{C}$ a suffisamment d'objets.
Toutefois, si $\mathcal{C}$ est petite, alors $\widehat{\mathcal{C}}$ est
localement petite et admet toutes les limites et colimites petites
(calculées argument par argument).
\end{remarque}

\begin{definition}[Foncteur $\Hom$ covariant]\label{def:hom-covariant}
\index{foncteur Hom@foncteur $\Hom$!covariant}
\index{h@$h^A$}
Soit $\mathcal{C}$ une catégorie localement petite et $A \in \Ob(\mathcal{C})$.
Le \textbf{foncteur $\Hom$ covariant} représenté par $A$ est le foncteur
\[
    h^A = \Hom(A, -) \colon \mathcal{C} \longrightarrow \Set
\]
défini sur les objets par $h^A(B) = \Hom(A, B)$ et sur les morphismes
$f \colon B \to C$ par la post-composition :
\[
    h^A(f) = f_* \colon \Hom(A, B) \longrightarrow \Hom(A, C), \quad
    g \longmapsto f \circ g.
\]
\end{definition}

\begin{definition}[Foncteur $\Hom$ contravariant]\label{def:hom-contravariant}
\index{foncteur Hom@foncteur $\Hom$!contravariant}
\index{h@$h_A$}
Soit $A \in \Ob(\mathcal{C})$.
Le \textbf{foncteur $\Hom$ contravariant} représenté par $A$ est le foncteur
\[
    h_A = \Hom(-, A) \colon \mathcal{C}^{\op} \longrightarrow \Set
\]
défini sur les objets par $h_A(B) = \Hom(B, A)$ et sur les morphismes
$f \colon B \to C$ (dans $\mathcal{C}$, donc $f \colon C \to B$ dans
$\mathcal{C}^{\op}$) par la pré-composition :
\[
    h_A(f) = f^* \colon \Hom(C, A) \longrightarrow \Hom(B, A), \quad
    g \longmapsto g \circ f.
\]
Ainsi $h_A$ est un préfaisceau sur $\mathcal{C}$.
\end{definition}

\begin{exemple}\label{ex:hom-foncteurs}
\index{foncteur Hom@foncteur $\Hom$!exemples}
\begin{enumerate}[label=(\roman*)]
\item Dans $\Set$, le foncteur $h^{\{*\}} = \Hom(\{*\}, -)$ est
    (naturellement isomorphe à) le foncteur identité $\Id_{\Set}$,
    car un élément de $\Hom(\{*\}, X)$ correspond exactement
    à un élément de $X$.

\item Dans $\Ab$, le foncteur $h^{\Z} = \Hom(\Z, -)$ est le
    foncteur oubli\index{foncteur oubli} $U \colon \Ab \to \Set$,
    car un morphisme de groupes $\Z \to G$ est déterminé par l'image de $1$.

\item Dans $\Top$, le foncteur $h^{\{*\}}$ est le foncteur
    ensemble sous-jacent. Le foncteur $h^{[0,1]}$ est le foncteur
    des chemins\index{foncteur!des chemins}, dont l'étude mène à la
    théorie de l'homotopie.

\item Dans $\mathsf{Vect}_k$, le foncteur $h_V = \Hom(-, V)$
    envoie $W$ sur l'espace des applications linéaires $W \to V$.
    En particulier, $h_{k} = \Hom(-, k)$ est le foncteur dual\index{foncteur dual}.
\end{enumerate}
\end{exemple}

\begin{definition}[Foncteur $\Hom$ bifunctoriel]\label{def:hom-bifuncteur}
\index{foncteur Hom@foncteur $\Hom$!bifuncteur}
\index{bifuncteur}
Le \textbf{foncteur $\Hom$ bifunctoriel} est le foncteur
\[
    \Hom(-,-) \colon \mathcal{C}^{\op} \times \mathcal{C} \longrightarrow \Set
\]
défini sur les paires d'objets par $\Hom(A, B)$ et sur les paires de morphismes
$(f \colon A' \to A,\, g \colon B \to B')$ par
\[
    \Hom(f, g) \colon \Hom(A, B) \longrightarrow \Hom(A', B'), \quad
    h \longmapsto g \circ h \circ f.
\]
\end{definition}

\begin{proposition}\label{prop:hom-bifuncteur-fonctorialite}
\index{foncteur Hom@foncteur $\Hom$!fonctorialité}
Le foncteur $\Hom(-,-)$ est bien un bifuncteur : pour tous morphismes
composables $(f_1, g_1)$ et $(f_2, g_2)$, on a
\[
    \Hom(f_1 \circ f_2,\, g_2 \circ g_1) = \Hom(f_2, g_2) \circ \Hom(f_1, g_1),
\]
et $\Hom(\id_A, \id_B) = \id_{\Hom(A,B)}$.
\end{proposition}

\begin{proof}
Pour tout $h \in \Hom(A, B)$, on a
\begin{align*}
    \Hom(f_1 \circ f_2,\, g_2 \circ g_1)(h)
    &= (g_2 \circ g_1) \circ h \circ (f_1 \circ f_2) \\
    &= g_2 \circ (g_1 \circ h \circ f_1) \circ f_2 \\
    &= \Hom(f_2, g_2)\bigl(\Hom(f_1, g_1)(h)\bigr).
\end{align*}
L'assertion sur les identités est immédiate.
\end{proof}


%% ============================================================
\section{Le plongement de Yoneda}\label{sec:plongement-yoneda}
%% ============================================================

\begin{definition}[Plongement de Yoneda covariant]\label{def:plongement-yoneda-cov}
\index{plongement de Yoneda!covariant}
\index{Yoneda!plongement}
Le \textbf{plongement de Yoneda covariant} est le foncteur
\[
    \mathbf{y} \colon \mathcal{C} \longrightarrow \Fun(\mathcal{C}^{\op}, \Set)
    = \widehat{\mathcal{C}}, \qquad
    A \longmapsto h_A = \Hom(-, A).
\]
Sur les morphismes $f \colon A \to B$, il est défini par la transformation
naturelle $\mathbf{y}(f) = f_* \colon h_A \Rightarrow h_B$ dont la composante
en $C \in \mathcal{C}$ est
\[
    (f_*)_C \colon \Hom(C, A) \longrightarrow \Hom(C, B), \qquad
    g \longmapsto f \circ g.
\]
\end{definition}

\begin{definition}[Plongement de Yoneda contravariant]\label{def:plongement-yoneda-contrav}
\index{plongement de Yoneda!contravariant}
De même, le \textbf{plongement de Yoneda contravariant} est le foncteur
\[
    \mathbf{y}^{\op} \colon \mathcal{C}^{\op} \longrightarrow
    \Fun(\mathcal{C}, \Set), \qquad
    A \longmapsto h^A = \Hom(A, -).
\]
\end{definition}

\begin{remarque}\label{rem:yoneda-notations}
\index{plongement de Yoneda!notations}
Suivant les auteurs, le plongement de Yoneda covariant est parfois noté
$\mathcal{C} \hookrightarrow \widehat{\mathcal{C}}$,
$A \mapsto \mathcal{C}(-, A)$, ou encore $A \mapsto \mathrm{Yo}(A)$.
Nous utiliserons principalement la notation $h_A$ pour le foncteur
contravariant représenté par $A$ et $h^A$ pour le covariant.
\end{remarque}

\begin{lemme}\label{lem:y-bien-defini}
\index{plongement de Yoneda!bonne définition}
Le plongement de Yoneda $\mathbf{y}$ est un foncteur bien défini :
pour tout morphisme $f \colon A \to B$ dans $\mathcal{C}$,
la famille $(f_*)_C = f \circ (-)$ est une transformation naturelle
$h_A \Rightarrow h_B$.
\end{lemme}

\begin{proof}
Il faut vérifier que pour tout morphisme $u \colon C' \to C$ dans $\mathcal{C}$,
le diagramme suivant commute :
\[
\begin{tikzcd}
\Hom(C, A) \arrow[r, "(f_*)_C"] \arrow[d, "u^*"'] &
\Hom(C, B) \arrow[d, "u^*"] \\
\Hom(C', A) \arrow[r, "(f_*)_{C'}"'] &
\Hom(C', B)
\end{tikzcd}
\]
Pour tout $g \in \Hom(C, A)$, on a
\[
    u^* \circ (f_*)_C(g) = u^*(f \circ g) = f \circ g \circ u
    = (f_*)_{C'}(g \circ u) = (f_*)_{C'} \circ u^*(g).
\]
Donc le diagramme commute.
\end{proof}


%% ============================================================
\section{Le lemme de Yoneda}\label{sec:lemme-yoneda}
%% ============================================================

Nous énonçons et démontrons maintenant le résultat central de ce chapitre.
Nous traitons d'abord la version contravariante (préfaisceaux), puis
indiquons la version covariante.

\begin{theoreme}[Lemme de Yoneda]\label{thm:lemme-yoneda}
\index{Yoneda, lemme de!énoncé}
\index{transformation naturelle!et lemme de Yoneda}
Soit $\mathcal{C}$ une catégorie localement petite,
$A \in \Ob(\mathcal{C})$, et $F \colon \mathcal{C}^{\op} \to \Set$
un préfaisceau. Il existe une bijection
\[
    \Phi_A \colon \Nat(h_A, F) \xrightarrow{\;\sim\;} F(A)
\]
qui est naturelle en $A$ et en $F$.
Explicitement, $\Phi_A$ est définie par
\[
    \Phi_A(\alpha) = \alpha_A(\id_A)
\]
pour toute transformation naturelle $\alpha \colon h_A \Rightarrow F$.
\end{theoreme}

\begin{proof}
\index{Yoneda, lemme de!preuve}
La preuve se décompose en quatre étapes : construction de la bijection,
construction de l'inverse, vérification que ces constructions sont
mutuellement inverses, et preuve de la naturalité.

\medskip
\textbf{Étape 1 : Définition de $\Phi_A$.}
Soit $\alpha \colon h_A \Rightarrow F$ une transformation naturelle.
Pour chaque objet $B$ de $\mathcal{C}$, $\alpha$ possède une composante
$\alpha_B \colon \Hom(B, A) \to F(B)$.
En particulier, pour $B = A$, on obtient
$\alpha_A \colon \Hom(A, A) \to F(A)$.
On pose
\[
    \Phi_A(\alpha) \coloneqq \alpha_A(\id_A) \in F(A).
\]

\medskip
\textbf{Étape 2 : Construction de l'inverse $\Psi_A$.}
Soit $x \in F(A)$. On définit une transformation naturelle
$\Psi_A(x) = \beta^x \colon h_A \Rightarrow F$ en posant,
pour chaque objet $B$ de $\mathcal{C}$ et chaque morphisme
$f \in \Hom(B, A) = h_A(B)$,
\[
    \beta^x_B(f) \coloneqq F(f)(x) \in F(B).
\]
Ici, $F(f) \colon F(A) \to F(B)$ est l'application induite par $F$
(rappelons que $F$ est un foncteur de $\mathcal{C}^{\op}$ vers $\Set$,
donc un morphisme $f \colon B \to A$ dans $\mathcal{C}$ induit
$F(f) \colon F(A) \to F(B)$).

Vérifions que $\beta^x$ est bien une transformation naturelle.
Pour un morphisme $u \colon C \to B$ dans $\mathcal{C}$,
il faut montrer que le diagramme suivant commute :
\[
\begin{tikzcd}
\Hom(B, A) \arrow[r, "\beta^x_B"] \arrow[d, "u^*"'] &
F(B) \arrow[d, "F(u)"] \\
\Hom(C, A) \arrow[r, "\beta^x_C"'] &
F(C)
\end{tikzcd}
\]
Pour tout $f \in \Hom(B, A)$, on a
\begin{align*}
    F(u) \circ \beta^x_B(f)
    &= F(u)\bigl(F(f)(x)\bigr) \\
    &= \bigl(F(u) \circ F(f)\bigr)(x) \\
    &= F(f \circ u)(x) \quad \text{(fonctorialité de $F$ sur $\mathcal{C}^{\op}$)} \\
    &= \beta^x_C(f \circ u) \\
    &= \beta^x_C(u^*(f)).
\end{align*}
Donc $\beta^x$ est bien naturelle.

\medskip
\textbf{Étape 3 : $\Phi_A$ et $\Psi_A$ sont mutuellement inverses.}

\emph{Vérification $\Phi_A \circ \Psi_A = \id_{F(A)}$.}
Pour tout $x \in F(A)$,
\[
    \Phi_A(\Psi_A(x))
    = \Phi_A(\beta^x)
    = \beta^x_A(\id_A)
    = F(\id_A)(x)
    = \id_{F(A)}(x)
    = x.
\]

\emph{Vérification $\Psi_A \circ \Phi_A = \id_{\Nat(h_A, F)}$.}
Pour toute transformation naturelle $\alpha \colon h_A \Rightarrow F$,
posons $x = \Phi_A(\alpha) = \alpha_A(\id_A)$.
Il faut montrer que $\beta^x = \alpha$,
c'est-à-dire que pour tout objet $B$ et tout $f \in \Hom(B, A)$,
\[
    \beta^x_B(f) = \alpha_B(f).
\]
Par définition, $\beta^x_B(f) = F(f)(x) = F(f)(\alpha_A(\id_A))$.
Or la naturalité de $\alpha$ appliquée au morphisme $f \colon B \to A$
donne le diagramme commutatif :
\[
\begin{tikzcd}
\Hom(A, A) \arrow[r, "\alpha_A"] \arrow[d, "f^*"'] &
F(A) \arrow[d, "F(f)"] \\
\Hom(B, A) \arrow[r, "\alpha_B"'] &
F(B)
\end{tikzcd}
\]
En évaluant sur $\id_A \in \Hom(A, A)$, on obtient
\[
    F(f)(\alpha_A(\id_A)) = \alpha_B(f^*(\id_A)) = \alpha_B(\id_A \circ f)
    = \alpha_B(f).
\]
Donc $\beta^x_B(f) = \alpha_B(f)$ pour tout $B$ et tout $f$,
ce qui montre $\Psi_A(\Phi_A(\alpha)) = \alpha$.

\medskip
\textbf{Étape 4 : Naturalité.}

\emph{Naturalité en $A$.}
Soit $u \colon A' \to A$ un morphisme dans $\mathcal{C}$.
Il faut vérifier que le diagramme suivant commute :
\[
\begin{tikzcd}
\Nat(h_A, F) \arrow[r, "\Phi_A"] \arrow[d, "u^*"'] &
F(A) \arrow[d, "F(u)"] \\
\Nat(h_{A'}, F) \arrow[r, "\Phi_{A'}"'] &
F(A')
\end{tikzcd}
\]
où $u^*(\alpha) = \alpha \circ \mathbf{y}(u) = \alpha \circ u_*$,
c'est-à-dire la transformation naturelle dont la composante en $B$ est
$\alpha_B \circ (u \circ -)$.
Pour toute $\alpha \colon h_A \Rightarrow F$,
\begin{align*}
    F(u)(\Phi_A(\alpha))
    &= F(u)(\alpha_A(\id_A)), \\
    \Phi_{A'}(u^*(\alpha))
    &= (u^*(\alpha))_{A'}(\id_{A'})
    = \alpha_{A'}(u \circ \id_{A'})
    = \alpha_{A'}(u).
\end{align*}
Par naturalité de $\alpha$ en le morphisme $u \colon A' \to A$,
$F(u)(\alpha_A(\id_A)) = \alpha_{A'}(u)$. Le diagramme commute.

\emph{Naturalité en $F$.}
Soit $\eta \colon F \Rightarrow G$ une transformation naturelle
entre préfaisceaux. On vérifie que
\[
\begin{tikzcd}
\Nat(h_A, F) \arrow[r, "\Phi^F_A"] \arrow[d, "\eta_*"'] &
F(A) \arrow[d, "\eta_A"] \\
\Nat(h_A, G) \arrow[r, "\Phi^G_A"'] &
G(A)
\end{tikzcd}
\]
commute, où $\eta_*(\alpha) = \eta \circ \alpha$. Pour toute $\alpha$,
\[
    \eta_A(\Phi^F_A(\alpha)) = \eta_A(\alpha_A(\id_A))
    = (\eta \circ \alpha)_A(\id_A)
    = \Phi^G_A(\eta_*(\alpha)). \qedhere
\]
\end{proof}

\begin{remarque}\label{rem:yoneda-covariant}
\index{Yoneda, lemme de!version covariante}
La version \emph{covariante} du lemme de Yoneda s'énonce comme suit.
Pour tout foncteur $G \colon \mathcal{C} \to \Set$ et tout
$A \in \Ob(\mathcal{C})$, on a une bijection naturelle
\[
    \Nat(h^A, G) \cong G(A),
    \qquad \alpha \mapsto \alpha_A(\id_A).
\]
La preuve est entièrement analogue.
\end{remarque}


%% ============================================================
\section{Corollaires du lemme de Yoneda}\label{sec:yoneda-corollaires}
%% ============================================================

\begin{corollaire}[Le plongement de Yoneda est pleinement fidèle]%
\label{cor:yoneda-pleinement-fidele}
\index{plongement de Yoneda!pleinement fidèle}
\index{foncteur!pleinement fidèle}
Le foncteur $\mathbf{y} \colon \mathcal{C} \to \widehat{\mathcal{C}}$,
$A \mapsto h_A$, est pleinement fidèle.
Autrement dit, pour tous objets $A, B$ de $\mathcal{C}$,
\[
    \mathbf{y}_{A,B} \colon \Hom_{\mathcal{C}}(A, B)
    \xrightarrow{\;\sim\;} \Nat(h_A, h_B)
\]
est une bijection.
\end{corollaire}

\begin{proof}
Appliquons le lemme de Yoneda~\ref{thm:lemme-yoneda} au préfaisceau
$F = h_B = \Hom(-, B)$. On obtient la bijection naturelle
\[
    \Phi_A \colon \Nat(h_A, h_B) \xrightarrow{\;\sim\;} h_B(A) = \Hom(A, B).
\]
Par construction, $\Phi_A(\alpha) = \alpha_A(\id_A)$ pour toute
transformation naturelle $\alpha \colon h_A \Rightarrow h_B$.

Or, par définition du plongement de Yoneda, pour tout morphisme
$f \colon A \to B$, on a $\mathbf{y}(f) = f_*$ avec
$(f_*)_C(g) = f \circ g$ pour tout $g \in \Hom(C, A)$.
Donc
\[
    \Phi_A(\mathbf{y}(f)) = (f_*)_A(\id_A) = f \circ \id_A = f.
\]
Cela montre que $\Phi_A \circ \mathbf{y}_{A,B} = \id_{\Hom(A,B)}$,
donc $\mathbf{y}_{A,B}$ est l'inverse de $\Phi_A$, et en particulier
c'est une bijection.
\end{proof}

\begin{corollaire}\label{cor:h-iso-implique-iso}
\index{Yoneda, lemme de!conséquence isomorphisme}
\index{isomorphisme!et foncteurs représentables}
Soient $A, B$ des objets d'une catégorie localement petite $\mathcal{C}$.
Alors
\[
    h_A \cong h_B \quad \text{dans } \widehat{\mathcal{C}}
    \qquad \Longleftrightarrow \qquad
    A \cong B \quad \text{dans } \mathcal{C}.
\]
De même, $h^A \cong h^B$ dans $\Fun(\mathcal{C}, \Set)$
si et seulement si $A \cong B$ dans $\mathcal{C}$.
\end{corollaire}

\begin{proof}
$(\Leftarrow)$ Si $f \colon A \xrightarrow{\sim} B$ est un isomorphisme,
alors $\mathbf{y}(f) = f_* \colon h_A \Rightarrow h_B$ est un isomorphisme
dans $\widehat{\mathcal{C}}$ car $\mathbf{y}$ est un foncteur.

$(\Rightarrow)$ Supposons $\alpha \colon h_A \xrightarrow{\sim} h_B$
un isomorphisme de préfaisceaux.
Puisque $\mathbf{y}$ est pleinement fidèle
(corollaire~\ref{cor:yoneda-pleinement-fidele}), il existe un unique
$f \colon A \to B$ tel que $\mathbf{y}(f) = \alpha$,
et un unique $g \colon B \to A$ tel que $\mathbf{y}(g) = \alpha^{-1}$.
Alors
\[
    \mathbf{y}(g \circ f) = \mathbf{y}(g) \circ \mathbf{y}(f)
    = \alpha^{-1} \circ \alpha = \id_{h_A} = \mathbf{y}(\id_A),
\]
et la fidélité de $\mathbf{y}$ donne $g \circ f = \id_A$.
De même, $f \circ g = \id_B$, donc $f$ est un isomorphisme.
\end{proof}

\begin{corollaire}\label{cor:yoneda-elements}
\index{Yoneda, lemme de!et éléments}
Pour tout préfaisceau $F \in \widehat{\mathcal{C}}$ et tout objet
$A \in \mathcal{C}$, les éléments de $F(A)$ correspondent
bijectivement aux morphismes $h_A \to F$ dans $\widehat{\mathcal{C}}$.
En particulier, les «~éléments généralisés de type $A$~» de $F$
sont exactement les transformations naturelles $h_A \Rightarrow F$.
\end{corollaire}

\begin{remarque}\label{rem:yoneda-philosophie}
\index{Yoneda, lemme de!interprétation philosophique}
Le lemme de Yoneda exprime le principe suivant : \emph{un objet est
entièrement déterminé par ses relations avec tous les autres objets}.
Plus précisément, connaître le préfaisceau $h_A = \Hom(-, A)$ revient
à connaître $A$ à isomorphisme unique près. Ce point de vue
est central en géométrie algébrique, où les «~points à valeurs dans $S$~»
d'un schéma $X$ sont les morphismes $S \to X$, c'est-à-dire les éléments
de $h_X(S)$.
\end{remarque}


%% ============================================================
\section{Foncteurs représentables}\label{sec:foncteurs-representables}
%% ============================================================

\begin{definition}[Foncteur représentable]\label{def:foncteur-representable}
\index{foncteur représentable!définition}
\index{objet représentant}
Un préfaisceau $F \colon \mathcal{C}^{\op} \to \Set$ est dit
\textbf{représentable} s'il existe un objet $A \in \mathcal{C}$
et un isomorphisme naturel $\alpha \colon h_A \xrightarrow{\sim} F$.
L'objet $A$ est appelé l'\textbf{objet représentant} de $F$,
et le couple $(A, \alpha)$ est une \textbf{représentation} de $F$.

Par le lemme de Yoneda, donner un tel isomorphisme revient
à donner un élément $\xi = \alpha_A(\id_A) \in F(A)$
tel que pour tout objet $B$ et tout $x \in F(B)$,
il existe un unique $f \colon B \to A$ avec $F(f)(\xi) = x$.
L'élément $\xi$ est l'\textbf{élément universel}\index{élément universel}.
\end{definition}

\begin{exemple}\label{ex:representables}
\index{foncteur représentable!exemples}
\begin{enumerate}[label=(\roman*)]
\item \textbf{Foncteur oubli des groupes.}\index{foncteur oubli!des groupes}
    Le foncteur oubli $U \colon \Grp \to \Set$ est représenté par $\Z$,
    car $\Hom_{\Grp}(\Z, G) \cong U(G)$ via $\varphi \mapsto \varphi(1)$.
    L'élément universel est $1 \in U(\Z) = \Z$.

\item \textbf{Foncteur ensemble de puissance.}\index{foncteur!ensemble de puissance}
    Le foncteur $\mathcal{P} \colon \Set^{\op} \to \Set$ qui envoie
    $X$ sur $\mathcal{P}(X) = \{A \subseteq X\}$ est représenté
    par l'ensemble $\{0, 1\}$ : un sous-ensemble $A \subseteq X$
    correspond à sa fonction caractéristique $\chi_A \colon X \to \{0,1\}$.

\item \textbf{Produits.}\index{produit!comme foncteur représentable}
    Dans une catégorie $\mathcal{C}$ avec produits binaires, le foncteur
    $\mathcal{C}^{\op} \to \Set$ défini par
    $X \mapsto \Hom(X, A) \times \Hom(X, B)$
    est représenté par $A \times B$.

\item \textbf{Grassmanniennes.}\index{Grassmannienne}
    En géométrie algébrique, le foncteur de Grassmann
    $\mathrm{Grass}(r, n) \colon \mathsf{Sch}^{\op} \to \Set$
    qui envoie un schéma $S$ sur l'ensemble des sous-fibrés localement
    libres de rang $r$ de $\mathcal{O}_S^n$ est représentable par
    la Grassmannienne $\mathrm{Gr}(r, n)$.
\end{enumerate}
\end{exemple}

\begin{proposition}\label{prop:representant-unique}
\index{foncteur représentable!unicité du représentant}
Si $F$ est un préfaisceau représentable, son objet représentant
est unique à isomorphisme unique près.
\end{proposition}

\begin{proof}
Si $(A, \alpha)$ et $(B, \beta)$ sont deux représentations de $F$,
alors $\beta^{-1} \circ \alpha \colon h_A \xrightarrow{\sim} h_B$
est un isomorphisme de préfaisceaux. Par le
corollaire~\ref{cor:h-iso-implique-iso}, $A \cong B$ dans $\mathcal{C}$.
L'unicité de l'isomorphisme découle de la fidélité de $\mathbf{y}$.
\end{proof}


%% ============================================================
\section{Tout préfaisceau est colimite de représentables}%
\label{sec:prefaisceau-colimite}
%% ============================================================

\begin{definition}[Catégorie des éléments]\label{def:categorie-elements}
\index{catégorie des éléments}
\index{$\int_{\mathcal{C}} F$}
Soit $F \colon \mathcal{C}^{\op} \to \Set$ un préfaisceau.
La \textbf{catégorie des éléments} de $F$, notée
$\int_{\mathcal{C}} F$ ou $\mathrm{el}(F)$, est la catégorie dont :
\begin{itemize}
    \item les objets sont les paires $(C, x)$ avec $C \in \mathcal{C}$
    et $x \in F(C)$ ;
    \item un morphisme $(C, x) \to (C', x')$ est un morphisme
    $f \colon C \to C'$ dans $\mathcal{C}$ tel que $F(f)(x') = x$.
\end{itemize}
\end{definition}

\begin{proposition}\label{prop:el-projection}
\index{catégorie des éléments!foncteur projection}
Le foncteur de projection $\pi \colon \int_{\mathcal{C}} F \to \mathcal{C}$
défini par $\pi(C, x) = C$ et $\pi(f) = f$ est un foncteur fidèle.
De plus, la catégorie $\int_{\mathcal{C}} F$ est petite lorsque
$\mathcal{C}$ est petite.
\end{proposition}

\begin{proof}
La fidélité est immédiate puisque $\pi$ est injectif sur les morphismes.
Si $\mathcal{C}$ est petite, alors $\Ob(\int_{\mathcal{C}} F)
= \coprod_{C \in \mathcal{C}} F(C)$ est un ensemble, et de même pour
les morphismes.
\end{proof}

\begin{theoreme}[Tout préfaisceau est colimite de représentables]%
\label{thm:prefaisceau-colimite-representables}
\index{préfaisceau!comme colimite de représentables}
\index{colimite!de représentables}
Soit $\mathcal{C}$ une catégorie petite et
$F \colon \mathcal{C}^{\op} \to \Set$ un préfaisceau. Alors
\[
    F \cong \colim_{(C, x) \in \int_{\mathcal{C}} F} h_C
\]
dans $\widehat{\mathcal{C}}$, où la colimite est prise le long du
foncteur $\int_{\mathcal{C}} F \to \widehat{\mathcal{C}}$,
$(C, x) \mapsto h_C$.
\end{theoreme}

\begin{proof}
Notons $G = \colim_{(C,x) \in \int_{\mathcal{C}} F} h_C$.
Par le lemme de Yoneda, pour tout préfaisceau $P$,
\begin{align*}
    \Nat(G, P)
    &\cong \lim_{(C,x) \in (\int_{\mathcal{C}} F)^{\op}} \Nat(h_C, P) \\
    &\cong \lim_{(C,x) \in (\int_{\mathcal{C}} F)^{\op}} P(C).
\end{align*}
Un élément de cette limite est une famille compatible
$(s_{(C,x)})_{(C,x)} \in \prod_{(C,x)} P(C)$ telle que
pour tout morphisme $f \colon (C,x) \to (C', x')$ dans $\int F$,
on a $P(f)(s_{(C',x')}) = s_{(C,x)}$.
Or une transformation naturelle $\eta \colon F \Rightarrow P$ est
une famille $(\eta_C)_{C \in \mathcal{C}}$ avec
$\eta_C \colon F(C) \to P(C)$ naturelle en $C$.
La correspondance est $\eta \mapsto (\eta_C(x))_{(C,x)}$,
et on vérifie que c'est une bijection. Ainsi
$\Nat(G, P) \cong \Nat(F, P)$ naturellement en $P$,
et le lemme de Yoneda (appliqué dans $\widehat{\mathcal{C}}$)
donne $G \cong F$.
\end{proof}

\begin{remarque}\label{rem:colimite-representables-densite}
\index{Yoneda!densité}
\index{théorème de densité}
Le théorème~\ref{thm:prefaisceau-colimite-representables} est souvent
appelé \emph{théorème de densité} de Yoneda : les objets représentables
forment une famille «~dense~» dans la catégorie des préfaisceaux, au sens
où tout préfaisceau s'obtient par colimites à partir des représentables.
\end{remarque}


%% ============================================================
\section{Exercices}\label{sec:exo-yoneda}
%% ============================================================

\begin{exercice}\label{exo:yoneda-groupes}
\index{Yoneda, lemme de!exercice}
Soit $G$ un groupe vu comme une catégorie $\mathbf{B}G$ à un seul objet
$*$ avec $\Hom(*, *) = G$.
\begin{enumerate}[label=(\alph*)]
\item Montrer qu'un préfaisceau sur $\mathbf{B}G$ est la même chose
    qu'un $G$-ensemble\index{G-ensemble@$G$-ensemble} (à droite).
\item Identifier l'image du plongement de Yoneda dans ce cas.
\item Énoncer et vérifier le lemme de Yoneda dans ce contexte :
    pour tout $G$-ensemble $X$,
    $\Nat(h_*, X) \cong X$ (en tant qu'ensembles).
\end{enumerate}
\end{exercice}

\begin{exercice}\label{exo:yoneda-preordre}
\index{préordre!et préfaisceaux}
Soit $(P, \leq)$ un ensemble préordonné vu comme une catégorie.
\begin{enumerate}[label=(\alph*)]
\item Montrer qu'un préfaisceau sur $P$ est un foncteur
    $P^{\op} \to \Set$, c'est-à-dire un ensemble $F(p)$ pour chaque
    $p \in P$ avec des applications de restriction
    $F(q) \to F(p)$ pour $p \leq q$, compatibles avec la composition.
\item Décrire le préfaisceau représentable $h_p$ pour $p \in P$.
\item Interpréter le lemme de Yoneda dans ce cas.
\end{enumerate}
\end{exercice}

\begin{exercice}\label{exo:representable-limites}
\index{foncteur représentable!et limites}
Montrer que tout foncteur représentable $h_A \colon \mathcal{C}^{\op} \to \Set$
préserve les limites. (C'est-à-dire, $h_A$ envoie les colimites de
$\mathcal{C}$ en limites de $\Set$.)
\end{exercice}

\begin{exercice}\label{exo:elements-comma}
\index{catégorie des éléments!et catégorie comma}
Montrer que la catégorie des éléments $\int_{\mathcal{C}} F$ est isomorphe
à la catégorie comma $(\mathbf{y} \downarrow F)$, où
$\mathbf{y} \colon \mathcal{C} \to \widehat{\mathcal{C}}$ est le plongement
de Yoneda.
\end{exercice}

\begin{exercice}\label{exo:coyoneda}
\index{Yoneda, lemme de!co-Yoneda}
\index{co-Yoneda, lemme de}
(\emph{Lemme de co-Yoneda.})
Soit $F \colon \mathcal{C} \to \Set$ un foncteur covariant.
Montrer que pour tout $A \in \mathcal{C}$,
\[
    F(A) \cong \int^{C \in \mathcal{C}} \Hom(C, A) \times F(C),
\]
où $\int^C$ désigne la coégalisatrice (coend)\index{coend}.
\end{exercice}

\begin{exercice}\label{exo:presheaf-complete}
\index{préfaisceau!complétude}
Soit $\mathcal{C}$ une catégorie petite. Montrer que
$\widehat{\mathcal{C}} = \Fun(\mathcal{C}^{\op}, \Set)$
admet toutes les petites limites et colimites, calculées
argument par argument.
\end{exercice}

\begin{exercice}\label{exo:yoneda-extension}
\index{extension de Kan!et Yoneda}
(\emph{Extension de Kan à gauche le long de Yoneda.})
Soit $\mathcal{C}$ une catégorie petite, $\mathcal{D}$ une catégorie
cocomplète, et $G \colon \mathcal{C} \to \mathcal{D}$ un foncteur.
Montrer que l'extension de Kan à gauche $\mathrm{Lan}_{\mathbf{y}} G
\colon \widehat{\mathcal{C}} \to \mathcal{D}$ existe et est donnée par
\[
    (\mathrm{Lan}_{\mathbf{y}} G)(F)
    = \colim_{(C, x) \in \int_{\mathcal{C}} F} G(C).
\]
\end{exercice}

\index{Yoneda, lemme de|)}
\index{préfaisceau|)}
\index{foncteur représentable|)}


%%%%%%%%%%%%%%%%%%%%%%%%%%%%%%%%%%%%%%%%%%%%%%%%%%%%%%%%%%%%
%%         CHAPITRE 6 : MONADES ET ALGÈBRES               %%
%%%%%%%%%%%%%%%%%%%%%%%%%%%%%%%%%%%%%%%%%%%%%%%%%%%%%%%%%%%%

\chapter{Monades et algèbres}\label{ch:monades}
\minitoc

\vspace{1em}

Les monades\index{monade|(} constituent un concept unificateur remarquable
qui relie la théorie des adjonctions, l'algèbre universelle et
l'informatique théorique.
Toute adjonction engendre une monade ; réciproquement, toute monade
provient d'une adjonction, et ce de deux manières canoniques
(Eilenberg--Moore et Kleisli).
Le théorème de Beck, résultat culminant de ce chapitre, caractérise
les catégories qui peuvent être décrites comme des «~catégories d'algèbres~»
sur une monade.

Ce chapitre suppose connue la théorie des adjonctions développée
aux chapitres précédents.

%% ============================================================
\section{Définition et premiers exemples}\label{sec:monades-def}
%% ============================================================

\begin{definition}[Monade]\label{def:monade}
\index{monade!définition}
\index{endofoncteur}
\index{multiplication!d'une monade}
\index{unité!d'une monade}
Une \textbf{monade} sur une catégorie $\mathcal{C}$ est un triplet
$(T, \eta, \mu)$ où :
\begin{itemize}
    \item $T \colon \mathcal{C} \to \mathcal{C}$ est un endofoncteur ;
    \item $\eta \colon \Id_{\mathcal{C}} \Rightarrow T$ est une
    transformation naturelle appelée \textbf{unité} ;
    \item $\mu \colon T^2 \Rightarrow T$ est une transformation naturelle
    appelée \textbf{multiplication},
\end{itemize}
satisfaisant les axiomes suivants :
\begin{enumerate}[label=(\roman*)]
\item \textbf{Associativité}\index{monade!associativité} :
    $\mu \circ T\mu = \mu \circ \mu T$, c'est-à-dire le diagramme
    suivant commute :
\[
\begin{tikzcd}
T^3 \arrow[r, "T\mu"] \arrow[d, "\mu T"'] & T^2 \arrow[d, "\mu"] \\
T^2 \arrow[r, "\mu"'] & T
\end{tikzcd}
\]

\item \textbf{Unités}\index{monade!axiomes d'unité} :
    $\mu \circ T\eta = \id_T = \mu \circ \eta T$, c'est-à-dire :
\[
\begin{tikzcd}
T \arrow[r, "T\eta"] \arrow[dr, equal] & T^2 \arrow[d, "\mu"] &
T \arrow[l, "\eta T"'] \arrow[dl, equal] \\
& T &
\end{tikzcd}
\]
\end{enumerate}
\end{definition}

\begin{remarque}\label{rem:monade-monoide}
\index{monade!comme monoïde}
\index{monoïde!dans une catégorie monoïdale}
Une monade sur $\mathcal{C}$ est exactement un monoïde dans la catégorie
monoïdale $(\mathrm{End}(\mathcal{C}), \circ, \Id_{\mathcal{C}})$
des endofoncteurs de $\mathcal{C}$, munie de la composition comme produit
tensoriel.
\end{remarque}

\begin{notation}\label{not:monade}
\index{monade!notation}
Par abus de langage, on note souvent la monade simplement $T$,
les transformations $\eta$ et $\mu$ étant sous-entendues.
On écrit parfois $\eta^T$ et $\mu^T$ lorsqu'il y a ambiguïté.
\end{notation}

\begin{definition}[Morphisme de monades]\label{def:morphisme-monades}
\index{monade!morphisme de monades}
Soient $(S, \eta^S, \mu^S)$ et $(T, \eta^T, \mu^T)$ deux monades
sur $\mathcal{C}$.
Un \textbf{morphisme de monades} $\sigma \colon S \to T$ est une
transformation naturelle $\sigma \colon S \Rightarrow T$ telle que :
\[
    \sigma \circ \eta^S = \eta^T, \qquad
    \sigma \circ \mu^S = \mu^T \circ (\sigma \star \sigma),
\]
où $\sigma \star \sigma = \sigma T \circ S\sigma = T\sigma \circ \sigma S
\colon S^2 \Rightarrow T^2$.
\end{definition}


%% ============================================================
\section{Monades issues d'adjonctions}\label{sec:monade-adjonction}
%% ============================================================

\begin{theoreme}\label{thm:adjonction-monade}
\index{adjonction!et monade}
\index{monade!issue d'une adjonction}
Toute adjonction $(F \dashv G, \eta, \varepsilon) \colon
\mathcal{C} \rightleftarrows \mathcal{D}$ engendre une monade
$(T, \eta, \mu)$ sur $\mathcal{C}$, où
\[
    T = G \circ F, \qquad
    \mu = G \varepsilon F \colon T^2 = GFGF \Rightarrow GF = T.
\]
\end{theoreme}

\begin{proof}
L'unité $\eta \colon \Id_{\mathcal{C}} \Rightarrow GF = T$ est
l'unité de l'adjonction. Il faut vérifier les axiomes de monade.

\emph{Associativité.}
On a $\mu \circ T\mu = G\varepsilon F \circ GF \cdot G\varepsilon F
= G\varepsilon F \circ G(F G\varepsilon F)
= G(\varepsilon F \circ FG\varepsilon F)
= G(\varepsilon \cdot \varepsilon FGF)$... Plus précisément,
$\mu \circ T\mu$ a pour composante en $C$ :
\[
    G(\varepsilon_{FC}) \circ GF(G(\varepsilon_{FC}))
    = G\bigl(\varepsilon_{FC} \circ F G(\varepsilon_{FC})\bigr).
\]
Et $\mu \circ \mu T$ a pour composante en $C$ :
\[
    G(\varepsilon_{FC}) \circ G(\varepsilon_{GFFC})
    = G\bigl(\varepsilon_{FC} \circ \varepsilon_{GFFC}\bigr).
\]
Or la naturalité de $\varepsilon$ appliquée au morphisme
$\varepsilon_{FC} \colon FGFC \to FC$ donne
$\varepsilon_{FC} \circ FG(\varepsilon_{FC})
= \varepsilon_{FC} \circ \varepsilon_{FGFC}$,
ce qui est exactement l'identité de gauche.

\emph{Unité à gauche.}
$\mu \circ \eta T$ a pour composante en $C$ :
$G(\varepsilon_{FC}) \circ \eta_{GFC}$.
Par l'identité triangulaire $G\varepsilon \circ \eta G = \id_G$,
en l'évaluant sur $FC$ on obtient
$G(\varepsilon_{FC}) \circ \eta_{GFC} = \id_{GFC} = \id_{TC}$.

\emph{Unité à droite.}
$\mu \circ T\eta$ a pour composante en $C$ :
$G(\varepsilon_{FC}) \circ GF(\eta_C)$.
Par l'identité triangulaire $\varepsilon F \circ F\eta = \id_F$,
en l'évaluant sur $C$ on obtient
$\varepsilon_{FC} \circ F(\eta_C) = \id_{FC}$,
d'où en appliquant $G$ :
$G(\varepsilon_{FC}) \circ GF(\eta_C) = G(\id_{FC}) = \id_{GFC}$.
\end{proof}

\begin{exemple}\label{ex:monades-adjonctions}
\index{monade!exemples}
\begin{enumerate}[label=(\roman*)]
\item \textbf{Monade libre sur les groupes.}%
\index{monade!libre!groupes}
    L'adjonction libre/oubli $F \dashv U \colon \Grp \rightleftarrows \Set$
    engendre la monade $T = UF$ sur $\Set$ : pour un ensemble $X$,
    $T(X)$ est l'ensemble sous-jacent au groupe libre $F(X)$.
    L'unité $\eta_X \colon X \to T(X)$ envoie chaque élément sur le
    mot de longueur $1$ correspondant. La multiplication
    $\mu_X \colon T^2(X) \to T(X)$ «~aplatit~» les mots de mots
    en les concaténant.

\item \textbf{Monade libre sur les modules.}%
\index{monade!libre!modules}
    Pour un anneau $R$, l'adjonction libre/oubli
    $F \dashv U \colon R\text{-}\mathsf{Mod} \rightleftarrows \Set$
    donne $T(X) = \bigoplus_{x \in X} R \cong R^{(X)}$.
    L'unité envoie $x$ sur le vecteur de base $e_x$
    et la multiplication est la $R$-linéarisation.

\item \textbf{Monade ensemble de puissance.}%
\index{monade!ensemble de puissance}
    Le foncteur $\mathcal{P} \colon \Set \to \Set$
    avec $\eta_X(x) = \{x\}$ et $\mu_X(\mathfrak{A}) = \bigcup \mathfrak{A}$
    pour $\mathfrak{A} \in \mathcal{P}(\mathcal{P}(X))$
    est une monade. Elle provient de l'auto-adjonction
    $\mathcal{P} \dashv \mathcal{P}^{\op}$ (à travers l'isomorphisme
    $\Set^{\op} \cong \mathsf{CABA}$).

\item \textbf{Monades en informatique.}%
\index{monade!en informatique}
    En programmation fonctionnelle (Haskell, Scala), les monades modélisent
    les effets de calcul. La monade \texttt{Maybe} correspond au foncteur
    $T(X) = X \sqcup \{*\}$ (valeur optionnelle). La monade \texttt{List}
    correspond au foncteur \emph{monoïde libre} $T(X) = \coprod_{n \geq 0} X^n$.
    La monade d'état $T(X) = (X \times S)^S$ pour un ensemble d'états $S$
    modélise les calculs avec état mutable.
\end{enumerate}
\end{exemple}

\begin{proposition}\label{prop:comonade}
\index{comonade}
\index{comonade!co-unité}
\index{comonade!comultiplication}
Dualement, toute adjonction $(F \dashv G)$ engendre une \textbf{comonade}
$(K, \varepsilon, \delta)$ sur $\mathcal{D}$, où
$K = F \circ G$, $\varepsilon$ est la co-unité de l'adjonction,
et $\delta = F\eta G \colon K \Rightarrow K^2$.
\end{proposition}

\begin{proof}
C'est le dual du théorème~\ref{thm:adjonction-monade}, en considérant
l'adjonction dans les catégories opposées.
\end{proof}


%% ============================================================
\section{Algèbres d'Eilenberg--Moore}\label{sec:eilenberg-moore}
%% ============================================================

La question réciproque se pose : toute monade provient-elle d'une adjonction ?
La réponse est oui, et il y a (au moins) deux constructions canoniques.
Nous commençons par celle d'Eilenberg--Moore\index{Eilenberg--Moore|(}.

\begin{definition}[$T$-algèbre]\label{def:T-algebre}
\index{T-algèbre@$T$-algèbre!définition}
\index{algèbre!d'une monade}
\index{application de structure}
Soit $(T, \eta, \mu)$ une monade sur $\mathcal{C}$.
Une \textbf{$T$-algèbre} (ou \textbf{algèbre d'Eilenberg--Moore})
est un couple $(A, a)$ où :
\begin{itemize}
    \item $A \in \Ob(\mathcal{C})$ est un objet (le \textbf{porteur}) ;
    \item $a \colon T(A) \to A$ est un morphisme
    (l'\textbf{application de structure}),
\end{itemize}
satisfaisant :
\begin{enumerate}[label=(\roman*)]
\item \textbf{Unité}\index{T-algèbre@$T$-algèbre!axiome d'unité} :
    $a \circ \eta_A = \id_A$ :
\[
\begin{tikzcd}
A \arrow[r, "\eta_A"] \arrow[dr, equal] & T(A) \arrow[d, "a"] \\
& A
\end{tikzcd}
\]

\item \textbf{Associativité}\index{T-algèbre@$T$-algèbre!axiome d'associativité} :
    $a \circ \mu_A = a \circ T(a)$ :
\[
\begin{tikzcd}
T^2(A) \arrow[r, "T(a)"] \arrow[d, "\mu_A"'] & T(A) \arrow[d, "a"] \\
T(A) \arrow[r, "a"'] & A
\end{tikzcd}
\]
\end{enumerate}
\end{definition}

\begin{definition}[Morphisme de $T$-algèbres]\label{def:morphisme-T-algebre}
\index{T-algèbre@$T$-algèbre!morphisme}
Un \textbf{morphisme de $T$-algèbres} de $(A, a)$ vers $(B, b)$ est
un morphisme $f \colon A \to B$ dans $\mathcal{C}$ tel que
$f \circ a = b \circ T(f)$ :
\[
\begin{tikzcd}
T(A) \arrow[r, "T(f)"] \arrow[d, "a"'] & T(B) \arrow[d, "b"] \\
A \arrow[r, "f"'] & B
\end{tikzcd}
\]
\end{definition}

\begin{definition}[Catégorie d'Eilenberg--Moore]\label{def:categorie-EM}
\index{catégorie!d'Eilenberg--Moore}
\index{Eilenberg--Moore!catégorie}
\index{$\mathcal{C}^T$}
La \textbf{catégorie d'Eilenberg--Moore} de la monade $T$,
notée $\mathcal{C}^T$, est la catégorie dont les objets sont les
$T$-algèbres et les morphismes sont les morphismes de $T$-algèbres.
La composition et les identités sont héritées de $\mathcal{C}$.
\end{definition}

\begin{exemple}\label{ex:T-algebres}
\index{T-algèbre@$T$-algèbre!exemples}
\begin{enumerate}[label=(\roman*)]
\item \textbf{Monade libre sur les groupes.}
    La monade $T = UF$ sur $\Set$ (foncteur mot).
    Une $T$-algèbre est un ensemble $X$ muni d'une application
    $a \colon T(X) \to X$ qui est associative et unitale.
    Cela revient exactement à munir $X$ d'une structure de groupe.
    Ainsi $\Set^T \simeq \Grp$.\index{monade!libre!algèbres}

\item \textbf{Monade ensemble de puissance.}
    Les $\mathcal{P}$-algèbres sur $\Set$ sont les treillis complets
    avec morphismes préservant les suprema arbitraires.
    \index{treillis complet}

\item \textbf{Algèbre libre.}
    Pour toute monade $T$ et tout objet $C \in \mathcal{C}$,
    le couple $(T(C), \mu_C)$ est une $T$-algèbre, appelée
    la \textbf{$T$-algèbre libre}\index{T-algèbre@$T$-algèbre!libre}
    engendrée par $C$.
    En effet, $\mu_C \circ \eta_{TC} = \id_{TC}$ (unité de droite)
    et $\mu_C \circ \mu_{TC} = \mu_C \circ T\mu_C$ (associativité).
\end{enumerate}
\end{exemple}

\begin{theoreme}\label{thm:adjonction-EM}
\index{adjonction!d'Eilenberg--Moore}
\index{Eilenberg--Moore!adjonction}
Soit $(T, \eta, \mu)$ une monade sur $\mathcal{C}$.
Il existe une adjonction
\[
    F^T \dashv U^T \colon \mathcal{C}^T \rightleftarrows \mathcal{C}
\]
telle que la monade associée est $(T, \eta, \mu)$, où :
\begin{itemize}
    \item le foncteur d'oubli $U^T \colon \mathcal{C}^T \to \mathcal{C}$
    est défini par $U^T(A, a) = A$ et $U^T(f) = f$ ;
    \item le foncteur libre $F^T \colon \mathcal{C} \to \mathcal{C}^T$
    est défini par $F^T(C) = (TC, \mu_C)$ et
    $F^T(f) = T(f)$.
\end{itemize}
\end{theoreme}

\begin{proof}
\emph{$F^T$ est bien défini.}
Pour tout $C \in \mathcal{C}$, $(TC, \mu_C)$ est une $T$-algèbre
libre (exemple~\ref{ex:T-algebres}(iii)).
Pour un morphisme $f \colon C \to D$, $T(f) \colon TC \to TD$ est
un morphisme de $T$-algèbres car
$T(f) \circ \mu_C = \mu_D \circ T^2(f)$ par naturalité de $\mu$.

\emph{Adjonction.}
On définit l'unité $\eta^{\mathrm{adj}}_C = \eta_C \colon C \to U^T F^T(C) = TC$
et la co-unité $\varepsilon^{\mathrm{adj}}_{(A,a)} = a \colon F^T U^T(A,a)
= (TA, \mu_A) \to (A,a)$.
Vérifions que $a$ est un morphisme de $T$-algèbres :
$a \circ \mu_A = a \circ Ta$ par l'axiome d'associativité de $(A,a)$.

\emph{Identités triangulaires.}
\begin{itemize}
    \item $U^T \varepsilon \circ \eta_{U^T} \colon U^T \Rightarrow U^T$:
    pour $(A,a)$, la composante est $a \circ \eta_A = \id_A$
    par l'axiome d'unité. Donc c'est $\id_{U^T}$.

    \item $\varepsilon_{F^T} \circ F^T\eta \colon F^T \Rightarrow F^T$:
    pour $C$, la composante est $\mu_C \circ T\eta_C = \id_{TC}$
    par l'axiome d'unité de droite de la monade. Donc c'est $\id_{F^T}$.
\end{itemize}

\emph{La monade associée est $T$.}
$U^T F^T(C) = TC$, et la multiplication est
$U^T \varepsilon_{F^T C} = \mu_C$. L'unité est $\eta_C$.
\end{proof}


%% ============================================================
\section{Catégorie de Kleisli}\label{sec:kleisli}
%% ============================================================

\begin{definition}[Catégorie de Kleisli]\label{def:categorie-kleisli}
\index{Kleisli!catégorie de}
\index{catégorie!de Kleisli}
\index{$\mathcal{C}_T$}
Soit $(T, \eta, \mu)$ une monade sur $\mathcal{C}$.
La \textbf{catégorie de Kleisli} $\mathcal{C}_T$ est définie par :
\begin{itemize}
    \item $\Ob(\mathcal{C}_T) = \Ob(\mathcal{C})$ ;
    \item pour $A, B \in \Ob(\mathcal{C}_T)$, on pose
    $\Hom_{\mathcal{C}_T}(A, B) = \Hom_{\mathcal{C}}(A, TB)$ ;
    \item la composition de $f \colon A \to TB$ et $g \colon B \to TC$
    dans $\mathcal{C}_T$ est
    \[
        g \circ_T f \coloneqq \mu_C \circ Tg \circ f \colon A \to TC ;
    \]
    \item l'identité de $A$ dans $\mathcal{C}_T$ est $\eta_A \colon A \to TA$.
\end{itemize}
\end{definition}

\begin{proposition}\label{prop:kleisli-categorie}
\index{Kleisli!bonne définition}
La catégorie de Kleisli $\mathcal{C}_T$ est bien une catégorie.
\end{proposition}

\begin{proof}
\emph{Associativité.}
Pour $f \colon A \to TB$, $g \colon B \to TC$, $h \colon C \to TD$,
\begin{align*}
    h \circ_T (g \circ_T f)
    &= \mu_D \circ Th \circ (\mu_C \circ Tg \circ f) \\
    &= \mu_D \circ Th \circ \mu_C \circ Tg \circ f \\
    &= \mu_D \circ \mu_{TD} \circ T^2 h \circ Tg \circ f
    \quad \text{(naturalité de $\mu$)}.
\end{align*}
D'autre part,
\begin{align*}
    (h \circ_T g) \circ_T f
    &= \mu_D \circ T(\mu_D \circ Th \circ g) \circ f \\
    &= \mu_D \circ T\mu_D \circ T^2 h \circ Tg \circ f.
\end{align*}
Par l'associativité de la monade $\mu_D \circ \mu_{TD} = \mu_D \circ T\mu_D$,
les deux expressions sont égales.

\emph{Unité à gauche.}
$f \circ_T \eta_A = \mu_B \circ Tf \circ \eta_A
= \mu_B \circ \eta_{TB} \circ f = f$
par naturalité de $\eta$ et l'axiome d'unité.

\emph{Unité à droite.}
$\eta_B \circ_T f = \mu_B \circ T\eta_B \circ f = \id_{TB} \circ f = f$
par l'axiome d'unité de la monade.
\end{proof}

\begin{theoreme}\label{thm:adjonction-kleisli}
\index{Kleisli!adjonction}
\index{adjonction!de Kleisli}
Il existe une adjonction $F_T \dashv G_T \colon \mathcal{C}_T
\rightleftarrows \mathcal{C}$ telle que la monade associée est $T$, où :
\begin{itemize}
    \item $F_T \colon \mathcal{C} \to \mathcal{C}_T$ est défini par
    $F_T(A) = A$ sur les objets et $F_T(f) = \eta_B \circ f$ pour
    $f \colon A \to B$ ;
    \item $G_T \colon \mathcal{C}_T \to \mathcal{C}$ est défini par
    $G_T(A) = TA$ et $G_T(f) = \mu_B \circ Tf$ pour $f \colon A \to TB$.
\end{itemize}
\end{theoreme}

\begin{proof}
La vérification est directe. L'unité de l'adjonction est
$\eta_A \colon A \to G_T F_T(A) = TA$, et la co-unité est
$\varepsilon_A = \id_{TA} \colon F_T G_T(A) = TA \to TA$ dans
$\mathcal{C}_T$ (vue comme $\id_{TA} \in \Hom(TA, TA)$).
Les identités triangulaires se vérifient à l'aide des axiomes de monade.
\end{proof}

\begin{remarque}\label{rem:kleisli-EM-universalite}
\index{Kleisli!propriété universelle}
\index{Eilenberg--Moore!propriété universelle}
La catégorie de Kleisli est \emph{initiale} et celle d'Eilenberg--Moore
est \emph{terminale} parmi les adjonctions engendrant la monade $T$.
Plus précisément, si $(F \dashv G) \colon \mathcal{C} \rightleftarrows
\mathcal{D}$ est une adjonction engendrant $T$, il existe des foncteurs
uniques (à isomorphisme unique près) formant le diagramme commutatif :
\[
\begin{tikzcd}
& \mathcal{D} \arrow[dd, "K"] \\
\mathcal{C}_T \arrow[ur, "J"] \arrow[dr] & \\
& \mathcal{C}^T
\end{tikzcd}
\]
où $J$ est pleinement fidèle et $K$ est le foncteur de comparaison.
\end{remarque}


%% ============================================================
\section{Théorème de monadicité de Beck}\label{sec:beck}
%% ============================================================

Le théorème de Beck\index{Beck, théorème de|(} caractérise les adjonctions
pour lesquelles le foncteur de comparaison est une équivalence.

\begin{definition}[Paire coégalisatrice scindée]\label{def:paire-coeg-scindee}
\index{paire coégalisatrice scindée}
\index{coégalisateur!scindé}
Soit $\mathcal{C}$ une catégorie.
Une \textbf{paire coégalisatrice scindée} est un diagramme
\[
\begin{tikzcd}
A \arrow[r, shift left=0.5ex, "f"] \arrow[r, shift right=0.5ex, "g"'] &
B \arrow[r, "e"] \arrow[l, bend right=40, "s"'] &
C \arrow[l, bend right=40, "t"']
\end{tikzcd}
\]
tel que $e \circ f = e \circ g$, $e \circ t = \id_C$,
$f \circ s = \id_B$ et $g \circ s = t \circ e$.
\end{definition}

\begin{lemme}\label{lem:coeg-scindee-absolue}
\index{coégalisateur!absolu}
\index{paire coégalisatrice scindée!est coégalisateur}
Toute paire coégalisatrice scindée est un coégalisateur absolu,
c'est-à-dire que $e$ est le coégalisateur de $(f, g)$ et ce coégalisateur
est préservé par tout foncteur.
\end{lemme}

\begin{proof}
Montrons que $e$ est le coégalisateur de $(f, g)$.
On sait que $e \circ f = e \circ g$.
Soit $h \colon B \to D$ tel que $h \circ f = h \circ g$.
Posons $\bar{h} = h \circ t \colon C \to D$. Alors
$\bar{h} \circ e = h \circ t \circ e = h \circ g \circ s \cdot$...
Plus précisément :
\[
    \bar{h} \circ e = h \circ t \circ e.
\]
Or $g \circ s = t \circ e$, donc $t \circ e = g \circ s$ et
\[
    h \circ t \circ e = h \circ g \circ s = h \circ f \circ s = h \circ \id_B = h.
\]
Donc $\bar{h} \circ e = h$.

Pour l'unicité, si $\bar{h}' \circ e = h$, alors
$\bar{h}' = \bar{h}' \circ e \circ t = h \circ t = \bar{h}$.

Puisque les identités scindantes se transportent par tout foncteur $U$
(en appliquant $U$ aux morphismes $s, t$),
le coégalisateur est préservé par $U$.
\end{proof}

\begin{definition}[Foncteur de comparaison]\label{def:foncteur-comparaison}
\index{foncteur!de comparaison}
\index{comparaison, foncteur de}
Soit $(F \dashv G, \eta, \varepsilon)$ une adjonction
$\mathcal{C} \rightleftarrows \mathcal{D}$ engendrant la monade $T = GF$.
Le \textbf{foncteur de comparaison} est le foncteur
$K \colon \mathcal{D} \to \mathcal{C}^T$ défini par :
\begin{itemize}
    \item $K(D) = (GD,\, G\varepsilon_D)$ pour tout $D \in \mathcal{D}$ ;
    \item $K(f) = Gf$ pour tout morphisme $f$ dans $\mathcal{D}$.
\end{itemize}
\end{definition}

\begin{proposition}\label{prop:comparaison-bien-defini}
\index{foncteur!de comparaison!bonne définition}
Le foncteur de comparaison $K$ est bien défini :
$(GD, G\varepsilon_D)$ est une $T$-algèbre pour tout $D$.
\end{proposition}

\begin{proof}
\emph{Unité.}
$G\varepsilon_D \circ \eta_{GD} = \id_{GD}$ par l'identité triangulaire.

\emph{Associativité.}
$G\varepsilon_D \circ \mu_{GD} = G\varepsilon_D \circ G\varepsilon_{FGD}
= G(\varepsilon_D \circ \varepsilon_{FGD})$.
Par ailleurs, $G\varepsilon_D \circ GF(G\varepsilon_D)
= G(\varepsilon_D \circ FG\varepsilon_D)$.
La naturalité de $\varepsilon$ appliquée à $\varepsilon_D$ donne
$\varepsilon_D \circ FG\varepsilon_D = \varepsilon_D \circ \varepsilon_{FGD}$.
\end{proof}

\begin{definition}[Foncteur monadique]\label{def:foncteur-monadique}
\index{foncteur!monadique}
\index{monadique, foncteur}
Un foncteur $G \colon \mathcal{D} \to \mathcal{C}$ est dit
\textbf{monadique} s'il admet un adjoint à gauche $F$ et si le
foncteur de comparaison $K \colon \mathcal{D} \to \mathcal{C}^T$
(pour la monade $T = GF$) est une équivalence de catégories.
Il est \textbf{strictement monadique} si $K$ est un isomorphisme
de catégories.
\end{definition}

\begin{definition}[$G$-paire coégalisatrice scindée]\label{def:G-paire-coeg}
\index{paire coégalisatrice!$G$-scindée}
Soit $G \colon \mathcal{D} \to \mathcal{C}$ un foncteur.
Une paire de morphismes $f, g \colon A \rightrightarrows B$ dans $\mathcal{D}$
est une \textbf{$G$-paire coégalisatrice scindée} si la paire
$Gf, Gg \colon GA \rightrightarrows GB$ admet une paire coégalisatrice
scindée dans $\mathcal{C}$.
\end{definition}

\begin{theoreme}[Théorème de monadicité de Beck]\label{thm:beck}
\index{Beck, théorème de!énoncé}
\index{monadicité, théorème de}
Soit $G \colon \mathcal{D} \to \mathcal{C}$ un foncteur admettant
un adjoint à gauche $F$. Alors $G$ est monadique si et seulement si :
\begin{enumerate}[label=(\roman*)]
\item $G$ reflète les isomorphismes\index{foncteur!qui reflète les isomorphismes}
    (c'est-à-dire : si $Gf$ est un isomorphisme dans $\mathcal{C}$,
    alors $f$ est un isomorphisme dans $\mathcal{D}$) ;
\item $\mathcal{D}$ admet les coégalisateurs de toutes les
    $G$-paires coégalisatrices scindées, et $G$ les préserve.
\end{enumerate}
\end{theoreme}

\begin{proof}
\index{Beck, théorème de!preuve}
La preuve se décompose en deux implications.

\medskip
\textbf{($\Rightarrow$) Condition nécessaire.}

Supposons $G$ monadique, c'est-à-dire que $K \colon \mathcal{D} \to
\mathcal{C}^T$ est une équivalence.

\emph{(i) $G$ reflète les isomorphismes.}
Puisque $U^T \circ K = G$ et que $K$ est une équivalence,
il suffit de montrer que $U^T$ reflète les isomorphismes.
Soit $f \colon (A, a) \to (B, b)$ un morphisme de $T$-algèbres
tel que $U^T(f) = f \colon A \to B$ est un isomorphisme dans $\mathcal{C}$.
Alors $f^{-1}$ existe dans $\mathcal{C}$, et il faut montrer que
$f^{-1}$ est un morphisme de $T$-algèbres, c'est-à-dire
$f^{-1} \circ b = a \circ T(f^{-1})$.
Or $f \circ a = b \circ Tf$ donne
$a = f^{-1} \circ b \circ Tf$, donc
$a \circ T(f^{-1}) = f^{-1} \circ b \circ Tf \circ T(f^{-1})
= f^{-1} \circ b$.

\emph{(ii) Existence et préservation des coégalisateurs.}
Puisque $K$ est une équivalence, il suffit de montrer le résultat
pour $\mathcal{C}^T$. Si $(f, g) \colon (A,a) \rightrightarrows (B,b)$
est une paire dans $\mathcal{C}^T$ telle que $U^T f, U^T g$ admet
une paire coégalisatrice scindée dans $\mathcal{C}$, alors le coégalisateur
scindé est absolu (lemme~\ref{lem:coeg-scindee-absolue}) et on peut le relever
en une $T$-algèbre.

\medskip
\textbf{($\Leftarrow$) Condition suffisante.}

Supposons (i) et (ii). Montrons que $K$ est une équivalence
en vérifiant qu'il est pleinement fidèle et essentiellement surjectif.

\emph{$K$ est fidèle.}
$K$ est fidèle car $U^T \circ K = G$ et $U^T$ est fidèle
(comme foncteur d'oubli d'algèbres).

\emph{$K$ est plein.}
Soit $\varphi \colon K(D) \to K(D')$ un morphisme de $T$-algèbres,
c'est-à-dire $\varphi \colon GD \to GD'$ avec
$\varphi \circ G\varepsilon_D = G\varepsilon_{D'} \circ GF\varphi$.
Considérons le diagramme dans $\mathcal{D}$ :
\[
\begin{tikzcd}
FGFGD \arrow[r, shift left, "FG\varepsilon_D"] \arrow[r, shift right, "\varepsilon_{FGD}"'] &
FGD \arrow[r, "\varepsilon_D"] & D
\end{tikzcd}
\]
En appliquant $G$, on obtient une paire coégalisatrice scindée
(les sections proviennent de $\eta$). Par (ii), $\varepsilon_D$ est
le coégalisateur de cette paire dans $\mathcal{D}$.
De même pour $D'$. La condition de morphisme de $T$-algèbres sur
$\varphi$ assure que $F\varphi \colon FGD \to FGD'$ est compatible,
et par la propriété universelle du coégalisateur, il existe un unique
$f \colon D \to D'$ avec $Gf = \varphi$. Donc $K(f) = \varphi$.

\emph{$K$ est essentiellement surjectif.}
Soit $(A, a)$ une $T$-algèbre. Considérons la paire
\[
\begin{tikzcd}
FTA \arrow[r, shift left, "Fa"] \arrow[r, shift right, "\varepsilon_{FA} \circ FT(\eta_A)"'] &
FA
\end{tikzcd}
\]
Plus précisément, on considère les deux morphismes
$F(a), \varepsilon_{FA} \colon FGFA = FTA \rightrightarrows FA$.
En appliquant $G$, on obtient $GFa, G\varepsilon_{FA} \colon
GFGFA \rightrightarrows GFA$, qui est une $G$-paire coégalisatrice
scindée (les sections sont données par $\eta_{GFA}$ et $GF\eta_A$,
en utilisant les axiomes de la $T$-algèbre).

Par (ii), soit $e \colon FA \to D$ le coégalisateur de
$(F(a), \varepsilon_{FA})$ dans $\mathcal{D}$, préservé par $G$.
Alors $Ge \colon GFA = TA \to GD$ est le coégalisateur de
$(GFa, G\varepsilon_{FA}) = (Ta, \mu_A)$ dans $\mathcal{C}$.
Or $a$ coégalise cette paire ($a \circ Ta = a \circ \mu_A$
par l'axiome d'associativité), donc il existe un unique
$\bar{a} \colon GD \to A$ avec $\bar{a} \circ Ge = a$.

On vérifie que $\bar{a}$ est un isomorphisme.
En effet, $\bar{a} \circ Ge \circ \eta_A = a \circ \eta_A = \id_A$,
et on montre que $Ge \circ \eta_A \circ \bar{a} = \id_{GD}$
en utilisant que $Ge$ est un épimorphisme (étant un coégalisateur).
Donc $G(\bar{a})$... plus précisément, on montre que $\bar{a}$ induit
un isomorphisme $GD \cong A$ et par (i) ($G$ reflète les isomorphismes),
l'isomorphisme se relève.

Ainsi $(A, a) \cong K(D) = (GD, G\varepsilon_D)$ dans $\mathcal{C}^T$.
\end{proof}

\begin{remarque}\label{rem:beck-versions}
\index{Beck, théorème de!version stricte}
Il existe une version «~stricte~» du théorème de Beck :
$G$ est \emph{strictement} monadique si et seulement si
$G$ reflète les isomorphismes et $\mathcal{D}$ admet les
coégalisateurs de \emph{toutes} les paires réflexives
(c'est-à-dire $f, g \colon A \rightrightarrows B$ avec une section
commune $s \colon B \to A$, $fs = gs = \id_B$), et $G$ les préserve.

On utilise aussi fréquemment le \emph{théorème de monadicité par
coégalisateurs réflexifs}\index{coégalisateur!réflexif} :
si $G$ crée les coégalisateurs de $G$-paires coégalisatrices scindées
(au lieu de simplement les préserver), alors $G$ est strictement monadique.
\end{remarque}

\index{Beck, théorème de|)}


%% ============================================================
\section{Applications : catégories algébriques}\label{sec:categories-algebriques}
%% ============================================================

Le théorème de Beck permet de reconnaître les catégories
«~algébriques~», c'est-à-dire celles qui sont monadiques sur $\Set$.

\begin{definition}[Catégorie algébrique]\label{def:categorie-algebrique}
\index{catégorie!algébrique}
Une catégorie $\mathcal{D}$ est dite \textbf{algébrique} (au sens de
la monadicité) si le foncteur d'oubli $G \colon \mathcal{D} \to \Set$
est monadique.
\end{definition}

\begin{corollaire}\label{cor:grp-monadique}
\index{Grp@$\Grp$!monadicité}
\index{Ab@$\Ab$!monadicité}
\index{Ring@$\mathsf{Ring}$!monadicité}
Les catégories $\Grp$, $\Ab$, $R\text{-}\mathsf{Mod}$,
$\mathsf{Ring}$ sont monadiques sur $\Set$.
\end{corollaire}

\begin{proof}
Nous vérifions les conditions du théorème de Beck pour le foncteur
d'oubli $U \colon \Grp \to \Set$ (les autres cas sont analogues).

\emph{$U$ admet un adjoint à gauche.}
C'est le foncteur groupe libre $F \colon \Set \to \Grp$.

\emph{$U$ reflète les isomorphismes.}
Un morphisme de groupes $\varphi \colon G \to H$ tel que $U\varphi$
est bijectif est un isomorphisme de groupes (l'inverse ensembliste
est automatiquement un morphisme de groupes :
$\varphi^{-1}(xy) = \varphi^{-1}(x)\varphi^{-1}(y)$
car $\varphi$ est un morphisme bijectif).

\emph{$\Grp$ admet les coégalisateurs des $U$-paires coégalisatrices
scindées et $U$ les préserve.}
Plus généralement, $\Grp$ admet tous les coégalisateurs
(c'est la catégorie quotient par la relation d'équivalence engendrée).
Si la paire est $U$-scindée, le coégalisateur dans $\Grp$
est préservé par $U$ car les paires coégalisatrices scindées
donnent des coégalisateurs absolus
(lemme~\ref{lem:coeg-scindee-absolue}).
\end{proof}

\begin{exemple}\label{ex:non-monadique}
\index{catégorie!non monadique}
\begin{enumerate}[label=(\roman*)]
\item \textbf{$\Top$ n'est pas monadique sur $\Set$.}
    Le foncteur d'oubli $U \colon \Top \to \Set$ ne reflète pas
    les isomorphismes : une bijection continue n'est pas nécessairement
    un homéomorphisme. Donc $\Top$ n'est pas monadique sur $\Set$.

\item \textbf{Les espaces de Hausdorff compacts.}\index{espace de Hausdorff compact}
    Le foncteur d'oubli $U \colon \mathsf{CHaus} \to \Set$
    est monadique (théorème de Manes\index{Manes, théorème de}).
    La monade associée est la monade ultrafiltres
    $\beta \colon \Set \to \Set$\index{monade!ultrafiltres},
    qui envoie un ensemble sur l'ensemble de ses ultrafiltres.
\end{enumerate}
\end{exemple}

\begin{proposition}\label{prop:monadique-proprietes}
\index{catégorie!monadique!propriétés}
Soit $G \colon \mathcal{D} \to \mathcal{C}$ un foncteur monadique.
Alors :
\begin{enumerate}[label=(\roman*)]
\item $G$ reflète les isomorphismes ;
\item $G$ crée les limites que $\mathcal{C}$ possède
    (en particulier $\mathcal{D}$ est complète si $\mathcal{C}$ l'est) ;
\item $G$ est fidèle et conservatif.
\end{enumerate}
\end{proposition}

\begin{proof}
(i) est le théorème de Beck. Pour (ii), on montre que $U^T$ crée
les limites. Soit $\{(A_i, a_i)\}$ un diagramme de $T$-algèbres et
$L = \lim_i A_i$ dans $\mathcal{C}$ avec projections $\pi_i$.
Alors $TL = T(\lim_i A_i)$ et on définit $\ell \colon TL \to L$
comme l'unique morphisme tel que $\pi_i \circ \ell = a_i \circ T\pi_i$
pour tout $i$. On vérifie que $(L, \ell)$ est une $T$-algèbre et
que c'est la limite dans $\mathcal{C}^T$. Comme $K$ est une équivalence,
le résultat se transporte à $\mathcal{D}$.

(iii) découle de (i) et de l'existence de l'adjoint à gauche.
\end{proof}

\begin{theoreme}[Critère de monadicité de Beck simplifié]%
\label{thm:beck-simplifie}
\index{Beck, théorème de!critère simplifié}
\index{monadicité!critère simplifié}
Soit $G \colon \mathcal{D} \to \mathcal{C}$ un foncteur admettant un
adjoint à gauche. Si $G$ reflète les isomorphismes et si $\mathcal{D}$
admet tous les coégalisateurs réflexifs\index{coégalisateur!réflexif}
et $G$ les préserve, alors $G$ est monadique.
\end{theoreme}

\begin{proof}
Toute $G$-paire coégalisatrice scindée est en particulier réflexive
(la section $s$ fournit la section commune). Donc les hypothèses du
théorème de Beck~\ref{thm:beck} sont satisfaites.
\end{proof}


%% ============================================================
\section{Exercices}\label{sec:exo-monades}
%% ============================================================

\begin{exercice}\label{exo:monade-idempotente}
\index{monade!idempotente}
Une monade $(T, \eta, \mu)$ est dite \textbf{idempotente} si
$\mu \colon T^2 \Rightarrow T$ est un isomorphisme naturel.
\begin{enumerate}[label=(\alph*)]
\item Montrer que $T$ est idempotente si et seulement si
    $T\eta = \eta T \colon T \Rightarrow T^2$.
\item Montrer que $T$ est idempotente si et seulement si
    $U^T \colon \mathcal{C}^T \to \mathcal{C}$ est pleinement fidèle.
\item En déduire que les $T$-algèbres correspondent à une
    sous-catégorie pleine réflexive de $\mathcal{C}$.
\end{enumerate}
\end{exercice}

\begin{exercice}\label{exo:kleisli-informatique}
\index{Kleisli!en informatique}
Soit $T$ la monade \texttt{Maybe} sur $\Set$ :
$T(X) = X \sqcup \{\bot\}$, $\eta_X(x) = x$,
$\mu_X(x) = x$ si $x \in X$ et $\mu_X(\bot) = \bot$.
\begin{enumerate}[label=(\alph*)]
\item Vérifier les axiomes de monade.
\item Décrire la catégorie de Kleisli $\Set_T$ : quels sont les morphismes
    de $A$ vers $B$ ? Interpréter comme des fonctions partielles.
\item Décrire les $T$-algèbres (catégorie d'Eilenberg--Moore).
\end{enumerate}
\end{exercice}

\begin{exercice}\label{exo:monade-distribution}
\index{monade!de distribution}
Soit $T \colon \Set \to \Set$ le foncteur qui envoie $X$ sur l'ensemble
des distributions de probabilité finiment supportées sur $X$ :
\[
    T(X) = \Bigl\{ p \colon X \to [0,1] \;\Big|\;
    \mathrm{supp}(p) \text{ fini}, \sum_{x \in X} p(x) = 1 \Bigr\}.
\]
\begin{enumerate}[label=(\alph*)]
\item Définir $\eta$ et $\mu$ pour en faire une monade.
\item Identifier les $T$-algèbres comme des espaces convexes
    (ensembles munis de combinaisons convexes).
\end{enumerate}
\end{exercice}

\begin{exercice}\label{exo:adjonction-monade-aller-retour}
\index{adjonction!et monade}
Soit $(F \dashv G, \eta, \varepsilon)$ une adjonction engendrant
la monade $T$ et la comonade $K$.
\begin{enumerate}[label=(\alph*)]
\item Montrer que les identités triangulaires de l'adjonction sont
    exactement ce qu'il faut pour que $\eta$ et $\mu = G\varepsilon F$
    satisfassent les axiomes de monade.
\item Montrer que si $F$ est pleinement fidèle, alors
    $\eta$ est un isomorphisme naturel, et $T$ est idempotente.
\item Montrer que si $G$ est pleinement fidèle, alors
    $\varepsilon$ est un isomorphisme naturel, et $K$ est idempotente.
\end{enumerate}
\end{exercice}

\begin{exercice}\label{exo:algebres-libres-colimites}
\index{T-algèbre@$T$-algèbre!libre!colimites}
Soit $T$ une monade sur $\mathcal{C}$.
\begin{enumerate}[label=(\alph*)]
\item Montrer que le foncteur libre $F^T \colon \mathcal{C} \to \mathcal{C}^T$
    préserve les colimites.
\item Montrer que toute $T$-algèbre $(A, a)$ est le coégalisateur dans
    $\mathcal{C}^T$ de la paire
    \[
    \begin{tikzcd}
    F^T T(A) \arrow[r, shift left, "F^T(a)"]
    \arrow[r, shift right, "\varepsilon_{F^T(A)}"'] & F^T(A)
    \end{tikzcd}
    \]
    c'est-à-dire que toute algèbre est un coégalisateur d'algèbres libres.
\end{enumerate}
\end{exercice}

\begin{exercice}\label{exo:beck-top}
\index{Top@$\Top$!non monadicité}
Donner un exemple explicite montrant que le foncteur d'oubli
$U \colon \Top \to \Set$ ne reflète pas les isomorphismes.
En déduire que $\Top$ n'est pas monadique sur $\Set$.
\end{exercice}

\begin{exercice}\label{exo:distributive-law}
\index{loi distributive}
Soient $S$ et $T$ deux monades sur $\mathcal{C}$.
Une \textbf{loi distributive} de $S$ sur $T$ est une
transformation naturelle $\lambda \colon TS \Rightarrow ST$ compatible
avec les unités et multiplications des deux monades.
\begin{enumerate}[label=(\alph*)]
\item Écrire les quatre diagrammes de compatibilité que $\lambda$ doit
    satisfaire.
\item Montrer qu'une loi distributive $\lambda$ définit une monade
    $ST$ sur $\mathcal{C}$ avec unité $\eta^S \eta^T$ et multiplication
    $\mu^S \mu^T \circ S\lambda T$.
\item Donner un exemple avec $S = (-) \times M$ (monade écrivain pour
    un monoïde $M$) et $T = (-)^R$ (monade lecteur pour un ensemble $R$).
\end{enumerate}
\end{exercice}

\begin{exercice}\label{exo:monade-codensity}
\index{monade!de codensité}
(\emph{Monade de codensité.})
Soit $G \colon \mathcal{D} \to \mathcal{C}$ un foncteur qui
n'admet pas nécessairement d'adjoint à gauche.
On définit l'extension de Kan à droite
$\mathrm{Ran}_G G \colon \mathcal{C} \to \mathcal{C}$ (lorsqu'elle existe).
\begin{enumerate}[label=(\alph*)]
\item Montrer que $T = \mathrm{Ran}_G G$ est naturellement muni d'une
    structure de monade, appelée la \textbf{monade de codensité} de $G$.
\item Montrer que lorsque $G$ admet un adjoint à gauche $F$,
    la monade de codensité coïncide avec $GF$.
\item Pour $G \colon \mathsf{FinSet} \hookrightarrow \Set$ l'inclusion,
    montrer que la monade de codensité est la monade des ultrafiltres $\beta$.
\end{enumerate}
\end{exercice}

\index{Eilenberg--Moore|)}
\index{monade|)}
%%%%%%%%%%%%%%%%%%%%%%%%%%%%%%%%%%%%%%%%%%%%%%%%%%%%%%%%%%%%%
%% THÉORIE DES CATÉGORIES — Chapitres 7–8
%% Catégories Abéliennes et Suites Exactes /
%% Catégories Dérivées — Introduction
%%%%%%%%%%%%%%%%%%%%%%%%%%%%%%%%%%%%%%%%%%%%%%%%%%%%%%%%%%%%%

%%%%%%%%%%%%%%%%%%%%%%%%%%%%%%%%%%%%%%%%%%%%%%%%%%%%%%%%%%%%
%%           CHAPITRE 7 : CATÉGORIES ABÉLIENNES           %%
%%                  ET SUITES EXACTES                      %%
%%%%%%%%%%%%%%%%%%%%%%%%%%%%%%%%%%%%%%%%%%%%%%%%%%%%%%%%%%%%

\chapter{Catégories abéliennes et suites exactes}\label{ch:abeliennes}
\minitoc

\vspace{1em}

Les catégories abéliennes\index{catégorie!abélienne|(} constituent le cadre
algébrique dans lequel se développe l'algèbre homologique. Introduites par
Grothendieck\index{Grothendieck, Alexander} dans son célèbre article
\emph{Tôhoku}\index{Tôhoku (article)} (1957), elles axiomatisent les
propriétés essentielles des catégories de modules et de faisceaux qui
permettent de définir les foncteurs dérivés et de manipuler les suites
exactes. Dans ce chapitre, nous construisons la théorie pas à pas : catégories
pré-additives, additives, puis abéliennes, avant d'énoncer et de démontrer
les lemmes fondamentaux (serpent, cinq) qui sont les outils de base
du calcul homologique.

\index{suite exacte|(}

%% ============================================================
\section{Catégories pré-additives}\label{sec:preadditive}
%% ============================================================

\begin{definition}[Catégorie pré-additive]\label{def:preadditive}
\index{catégorie!pré-additive}
\index{enrichissement!sur $\Ab$}
Une catégorie $\mathcal{A}$ est dite \textbf{pré-additive}
(ou \textbf{$\Ab$-catégorie}\index{Ab-catégorie@$\Ab$-catégorie})
si, pour tous objets $A, B \in \Ob(\mathcal{A})$,
l'ensemble $\Hom_{\mathcal{A}}(A, B)$ est muni d'une structure de
groupe abélien\index{groupe abélien} telle que la composition
\[
  \circ \colon \Hom_{\mathcal{A}}(B, C) \times \Hom_{\mathcal{A}}(A, B)
  \longrightarrow \Hom_{\mathcal{A}}(A, C)
\]
est bilinéaire\index{bilinéarité de la composition} :
pour tous $f, f' \colon A \to B$ et $g, g' \colon B \to C$,
\[
  g \circ (f + f') = g \circ f + g \circ f', \qquad
  (g + g') \circ f = g \circ f + g' \circ f.
\]
\end{definition}

\begin{exemple}\label{ex:preadditive}
\index{catégorie!des modules}
\begin{enumerate}[label=(\roman*)]
\item Pour tout anneau $R$, la catégorie $R\text{-}\Mod$\index{R-Mod@$R\text{-}\Mod$}
  est pré-additive : si $f, g \colon M \to N$ sont des $R$-morphismes,
  on pose $(f + g)(m) = f(m) + g(m)$.

\item La catégorie $\Ab$\index{Ab@$\Ab$} des groupes abéliens
  est le cas $R = \Z$.

\item La catégorie $\Set$\index{Set@$\Set$} n'est \emph{pas}
  pré-additive, car $\Hom_{\Set}(A, B)$ n'a pas de structure naturelle
  de groupe abélien (sauf cas dégénérés).
\end{enumerate}
\end{exemple}

\begin{remarque}\label{rem:zero-morphism}
\index{morphisme nul}
Dans une catégorie pré-additive, chaque $\Hom(A, B)$ possède un
élément neutre $0_{A,B}$ pour l'addition. On a
$0_{B,C} \circ f = 0_{A,C}$ et $g \circ 0_{A,B} = 0_{A,C}$
pour tous $f \colon A \to B$ et $g \colon B \to C$.
\end{remarque}

\begin{definition}[Foncteur additif]\label{def:additive-functor}
\index{foncteur!additif}
Un foncteur $F \colon \mathcal{A} \to \mathcal{B}$ entre catégories
pré-additives est \textbf{additif} si, pour tous objets $A, B$ et tous
morphismes $f, g \colon A \to B$,
\[
  F(f + g) = F(f) + F(g).
\]
Autrement dit, $F_{A,B} \colon \Hom_{\mathcal{A}}(A, B) \to
\Hom_{\mathcal{B}}(FA, FB)$ est un morphisme de groupes abéliens.
\end{definition}


%% ============================================================
\section{Catégories additives et biproduits}\label{sec:additive}
%% ============================================================

\begin{definition}[Objet nul]\label{def:zero-object}
\index{objet nul}
\index{objet initial}
\index{objet terminal}
Un objet $0 \in \Ob(\mathcal{A})$ est un \textbf{objet nul}
s'il est à la fois initial et terminal, c'est-à-dire si
$\Hom(0, A)$ et $\Hom(A, 0)$ sont des singletons pour tout objet $A$.
\end{definition}

\begin{definition}[Biproduit]\label{def:biproduct}
\index{biproduit}
\index{somme directe!catégorique}
Soit $\mathcal{A}$ une catégorie pré-additive possédant un objet nul.
Un \textbf{biproduit}\index{biproduit|textbf} de deux objets $A, B$ est un
diagramme
\[
\begin{tikzcd}
A \arrow[r, "\iota_A"] & A \oplus B
  \arrow[l, bend right, "\pi_A"']
  \arrow[r, bend left, "\pi_B"]
& B \arrow[l, "\iota_B"']
\end{tikzcd}
\]
vérifiant :
\begin{enumerate}[label=(\alph*)]
\item $\pi_A \circ \iota_A = \id_A$, \quad $\pi_B \circ \iota_B = \id_B$;
\item $\pi_B \circ \iota_A = 0$, \quad $\pi_A \circ \iota_B = 0$;
\item $\iota_A \circ \pi_A + \iota_B \circ \pi_B = \id_{A \oplus B}$.
\end{enumerate}
\end{definition}

\begin{proposition}\label{prop:biproduct-product-coproduct}
\index{biproduit!produit et coproduit}
Dans une catégorie pré-additive, un biproduit $A \oplus B$ est
simultanément un produit $(A \oplus B, \pi_A, \pi_B)$ et un
coproduit $(A \oplus B, \iota_A, \iota_B)$.
\end{proposition}

\begin{proof}
Montrons que $A \oplus B$ est un produit. Soit $X$ un objet muni de
morphismes $f \colon X \to A$ et $g \colon X \to B$. Posons
$h = \iota_A \circ f + \iota_B \circ g \colon X \to A \oplus B$.
Alors
\[
  \pi_A \circ h = \pi_A \circ \iota_A \circ f + \pi_A \circ \iota_B \circ g
  = f + 0 = f,
\]
et de même $\pi_B \circ h = g$. Pour l'unicité, si $h'$ vérifie
$\pi_A \circ h' = f$ et $\pi_B \circ h' = g$, alors
\[
  h' = (\iota_A \circ \pi_A + \iota_B \circ \pi_B) \circ h'
  = \iota_A \circ f + \iota_B \circ g = h.
\]
L'argument dual montre que c'est un coproduit.
\end{proof}

\begin{definition}[Catégorie additive]\label{def:additive-cat}
\index{catégorie!additive}
Une catégorie $\mathcal{A}$ est \textbf{additive} si :
\begin{enumerate}[label=(\roman*)]
\item $\mathcal{A}$ est pré-additive ;
\item $\mathcal{A}$ possède un objet nul ;
\item tout couple d'objets $(A, B)$ admet un biproduit $A \oplus B$.
\end{enumerate}
\end{definition}

\begin{exemple}\label{ex:additive-cats}
\index{catégorie!des espaces vectoriels}
\begin{enumerate}[label=(\roman*)]
\item $R\text{-}\Mod$ est additive pour tout anneau $R$, avec
  $M \oplus N$ la somme directe usuelle.
\item $\Ab$ est additive.
\item La catégorie des espaces vectoriels de dimension finie
  $\mathsf{Vect}_k^{\mathrm{fd}}$\index{Vect@$\mathsf{Vect}_k$}
  est additive.
\end{enumerate}
\end{exemple}

\begin{proposition}[Matrices de morphismes]\label{prop:matrix-morphisms}
\index{matrice de morphismes}
Dans une catégorie additive, un morphisme
$f \colon A_1 \oplus A_2 \to B_1 \oplus B_2$ est entièrement déterminé
par la matrice
\[
  f \longleftrightarrow
  \begin{pmatrix} f_{11} & f_{12} \\ f_{21} & f_{22} \end{pmatrix},
  \qquad f_{ij} = \pi_{B_i} \circ f \circ \iota_{A_j}
  \colon A_j \to B_i,
\]
et la composition correspond au produit matriciel.
\end{proposition}

\begin{proof}
On a $f = \sum_{i,j} \iota_{B_i} \circ f_{ij} \circ \pi_{A_j}$
par la condition (c) du biproduit appliquée au domaine et au codomaine.
La vérification du produit matriciel pour la composition est un calcul
direct utilisant les relations d'orthogonalité.
\end{proof}


%% ============================================================
\section{Catégories abéliennes}\label{sec:abelian}
%% ============================================================

\begin{definition}[Noyau et conoyau]\label{def:kernel-cokernel}
\index{noyau!catégorique}
\index{conoyau}
Soit $f \colon A \to B$ un morphisme dans une catégorie pré-additive
avec objet nul.
\begin{enumerate}[label=(\roman*)]
\item Un \textbf{noyau}\index{noyau|textbf} de $f$ est un
  égalisateur\index{égalisateur} de $f$ et $0_{A,B}$, c'est-à-dire un
  morphisme $\ker f \colon K \to A$ tel que $f \circ \ker f = 0$ et,
  pour tout $g \colon X \to A$ avec $f \circ g = 0$, il existe un unique
  $\bar{g} \colon X \to K$ avec $\ker f \circ \bar{g} = g$ :
\[
\begin{tikzcd}
X \arrow[dr, "g"'] \arrow[r, dashed, "\exists!\,\bar{g}"]
& K \arrow[d, "\ker f"] \\
& A \arrow[r, "f"] & B
\end{tikzcd}
\]

\item Un \textbf{conoyau}\index{conoyau|textbf} de $f$ est un
  coégalisateur\index{coégalisateur} de $f$ et $0_{A,B}$, c'est-à-dire un
  morphisme $\coker f \colon B \to C$ tel que $\coker f \circ f = 0$ et
  vérifiant la propriété universelle duale.
\end{enumerate}
\end{definition}

\begin{definition}[Image et coimage]\label{def:image-coimage}
\index{image!catégorique}
\index{coimage}
Soit $f \colon A \to B$ un morphisme dans une catégorie possédant
noyaux et conoyaux.
\begin{enumerate}[label=(\roman*)]
\item L'\textbf{image} de $f$ est $\im f = \ker(\coker f)$.
\item La \textbf{coimage} de $f$ est $\operatorname{coim} f = \coker(\ker f)$.
\end{enumerate}
\end{definition}

\begin{definition}[Catégorie abélienne]\label{def:abelian-cat}
\index{catégorie!abélienne|textbf}
\index{axiomes de Grothendieck}
Une catégorie $\mathcal{A}$ est \textbf{abélienne} si :
\begin{enumerate}[label=(AB\arabic*)]
\item\label{AB0} $\mathcal{A}$ est additive ;
\item\label{AB1} tout morphisme admet un noyau et un conoyau ;
\item\label{AB2} pour tout morphisme $f \colon A \to B$, le morphisme
  canonique
  \[
    \operatorname{coim} f \longrightarrow \im f
  \]
  est un isomorphisme.
\end{enumerate}
\end{definition}

\begin{remarque}\label{rem:abelian-reformulation}
\index{monomorphisme!dans une catégorie abélienne}
\index{épimorphisme!dans une catégorie abélienne}
La condition \ref{AB2} équivaut à : tout monomorphisme est un noyau et tout
épimorphisme est un conoyau. Plus précisément, dans une catégorie
abélienne, un morphisme $f$ se factorise comme
\[
\begin{tikzcd}
A \arrow[r, two heads, "\operatorname{coim} f"]
& \im f \arrow[r, hook, "\im f"]
& B
\end{tikzcd}
\]
où la première flèche est un épimorphisme et la seconde un monomorphisme.
C'est l'analogue catégorique du premier théorème d'isomorphisme\index{théorème d'isomorphisme!premier}.
\end{remarque}

\begin{theoreme}[Théorème de plongement de Freyd--Mitchell]%
\label{thm:freyd-mitchell}
\index{théorème!de Freyd--Mitchell}
\index{Freyd, Peter}
\index{Mitchell, Barry}
Toute petite catégorie abélienne $\mathcal{A}$ admet un foncteur
exact et pleinement fidèle
\[
  F \colon \mathcal{A} \longrightarrow R\text{-}\Mod
\]
pour un certain anneau $R$. En particulier, toute relation
diagrammatique\index{chasse au diagramme} valable dans les catégories
de modules est valable dans toute catégorie abélienne.
\end{theoreme}

\begin{proof}
La démonstration complète, due à Freyd (1964) et Mitchell (1965),
est longue et technique. L'idée est la suivante :
\begin{enumerate}[label=\arabic*.]
\item On fixe un générateur $U$ de $\mathcal{A}$
  (qui existe car $\mathcal{A}$ est petite abélienne).
\item On pose $R = \End_{\mathcal{A}}(U)^{\op}$ et on considère le
  foncteur $F = \Hom_{\mathcal{A}}(U, -) \colon \mathcal{A} \to R\text{-}\Mod$.
\item On montre que $F$ est exact et pleinement fidèle en utilisant
  le fait que $U$ est un générateur.
\end{enumerate}
Nous renvoyons à \cite{freyd1964} ou \cite{mitchell1965} pour les
détails.
\end{proof}

\begin{remarque}\label{rem:freyd-mitchell-usage}
Le théorème de Freyd--Mitchell\index{théorème!de Freyd--Mitchell}
justifie la technique de la \emph{chasse au diagramme}\index{chasse au diagramme|textbf} :
pour démontrer un résultat dans une catégorie abélienne abstraite,
il suffit de le vérifier pour les modules, où l'on peut travailler
avec des éléments.
\end{remarque}


%% ============================================================
\section{Exemples de catégories abéliennes}\label{sec:abelian-examples}
%% ============================================================

\begin{exemple}[Modules sur un anneau]\label{ex:R-mod-abelian}
\index{R-Mod@$R\text{-}\Mod$!catégorie abélienne}
Pour tout anneau $R$ (commutatif ou non), la catégorie
$R\text{-}\Mod$ est abélienne. Les noyaux et conoyaux sont les
noyaux et conoyaux de morphismes de modules au sens usuel,
et la condition image $=$ coimage est le premier théorème
d'isomorphisme : $M/\ker f \cong \im f$.
\end{exemple}

\begin{exemple}[Groupes abéliens]\label{ex:ab-abelian}
\index{Ab@$\Ab$!catégorie abélienne}
$\Ab = \Z\text{-}\Mod$ est le cas particulier $R = \Z$.
\end{exemple}

\begin{exemple}[Faisceaux de modules]\label{ex:sheaves-abelian}
\index{faisceau!de modules}
\index{catégorie!des faisceaux}
Soit $(X, \mathcal{O}_X)$ un espace annelé\index{espace annelé}.
La catégorie $\mathcal{O}_X\text{-}\mathsf{Mod}$ des faisceaux de
$\mathcal{O}_X$-modules est abélienne. Les noyaux sont calculés
section par section\index{noyau!de faisceaux} (le noyau préfaisceautique
est déjà un faisceau), mais les conoyaux\index{conoyau!de faisceaux}
nécessitent une faisceautisation.
\end{exemple}

\begin{exemple}[Foncteurs vers $\Ab$]\label{ex:functor-category-abelian}
\index{catégorie!de foncteurs!abélienne}
Si $\mathcal{C}$ est une petite catégorie, alors la catégorie des foncteurs
$\Fun(\mathcal{C}, \Ab)$\index{Fun@$\Fun(\mathcal{C}, \Ab)$} est abélienne,
les noyaux et conoyaux étant calculés argument par argument
(« point par point »\index{point par point}).
\end{exemple}

\begin{exemple}[Contre-exemple : groupes non abéliens]\label{ex:grp-not-abelian}
\index{Grp@$\Grp$!pas abélienne}
La catégorie $\Grp$ des groupes (non nécessairement abéliens)
n'est \emph{pas} abélienne. Elle est pointée (le groupe trivial
est un objet nul) et admet noyaux et conoyaux, mais la condition
\ref{AB2} échoue : le morphisme canonique
$\operatorname{coim} f \to \im f$ n'est pas toujours un isomorphisme.
En effet, $\im f$ au sens catégorique est le plus petit sous-groupe
normal contenant l'image ensembliste, qui peut différer du quotient
$G/\ker f$.
\end{exemple}


%% ============================================================
\section{Suites exactes}\label{sec:exact-sequences}
%% ============================================================

\begin{definition}[Suite exacte en un point]\label{def:exact-at}
\index{suite exacte!en un point}
Soit $\mathcal{A}$ une catégorie abélienne. Une suite de morphismes
\[
  A \xrightarrow{f} B \xrightarrow{g} C
\]
est dite \textbf{exacte en $B$} si $\im f = \ker g$
(comme sous-objets\index{sous-objet} de $B$).
De manière équivalente, $g \circ f = 0$ et, pour tout $h \colon X \to B$
avec $g \circ h = 0$, le morphisme $h$ se factorise (uniquement) par
$\im f$.
\end{definition}

\begin{definition}[Suite exacte courte]\label{def:ses}
\index{suite exacte!courte}
\index{SES|see{suite exacte courte}}
Une \textbf{suite exacte courte} est une suite
\[
\begin{tikzcd}
0 \arrow[r] & A \arrow[r, "f"] & B \arrow[r, "g"] & C \arrow[r] & 0
\end{tikzcd}
\]
qui est exacte en $A$, $B$ et $C$. Cela signifie :
\begin{enumerate}[label=(\roman*)]
\item $f$ est un monomorphisme ($\ker f = 0$) ;
\item $g$ est un épimorphisme ($\coker g = 0$) ;
\item $\im f = \ker g$.
\end{enumerate}
On note souvent une telle suite $0 \to A \to B \to C \to 0$.
\end{definition}

\begin{exemple}\label{ex:ses-Z}
\index{suite exacte!courte!exemples}
\begin{enumerate}[label=(\roman*)]
\item Dans $\Ab$ : $0 \to \Z \xrightarrow{\times n} \Z \to \Z/n\Z \to 0$
  pour tout $n \geq 1$.

\item Dans $R\text{-}\Mod$ : si $N \leq M$ est un sous-module, alors
  $0 \to N \hookrightarrow M \twoheadrightarrow M/N \to 0$ est exacte.

\item La suite $0 \to \Z \xrightarrow{2} \Z \to \Z/2\Z \to 0$
  est une suite exacte courte qui \emph{ne scinde pas}\index{suite exacte!scindée}.
\end{enumerate}
\end{exemple}

\begin{definition}[Suite exacte scindée]\label{def:split-ses}
\index{suite exacte!scindée|textbf}
Une suite exacte courte
$0 \to A \xrightarrow{f} B \xrightarrow{g} C \to 0$
est \textbf{scindée} si l'une des conditions équivalentes suivantes
est satisfaite :
\begin{enumerate}[label=(\roman*)]
\item il existe $s \colon C \to B$ tel que $g \circ s = \id_C$
  (section\index{section} de $g$) ;
\item il existe $r \colon B \to A$ tel que $r \circ f = \id_A$
  (rétraction\index{rétraction} de $f$) ;
\item $B \cong A \oplus C$ via un isomorphisme compatible avec $f$ et $g$.
\end{enumerate}
\end{definition}


%% ============================================================
\section{Foncteurs exacts}\label{sec:exact-functors}
%% ============================================================

\begin{definition}[Foncteur exact à gauche, à droite, exact]\label{def:exact-functor}
\index{foncteur!exact}
\index{foncteur!exact à gauche}
\index{foncteur!exact à droite}
Soit $F \colon \mathcal{A} \to \mathcal{B}$ un foncteur additif entre
catégories abéliennes.
\begin{enumerate}[label=(\roman*)]
\item $F$ est \textbf{exact à gauche}\index{foncteur!exact à gauche|textbf}
  s'il transforme toute suite exacte courte
  $0 \to A \to B \to C \to 0$ en une suite exacte
  $0 \to FA \to FB \to FC$.

\item $F$ est \textbf{exact à droite}\index{foncteur!exact à droite|textbf}
  s'il transforme toute suite exacte courte
  $0 \to A \to B \to C \to 0$ en une suite exacte
  $FA \to FB \to FC \to 0$.

\item $F$ est \textbf{exact}\index{foncteur!exact|textbf}
  s'il est exact à gauche et à droite, c'est-à-dire s'il préserve
  les suites exactes courtes.
\end{enumerate}
\end{definition}

\begin{proposition}\label{prop:hom-left-exact}
\index{Hom@$\Hom$!exactitude à gauche}
Pour tout objet $A$ d'une catégorie abélienne $\mathcal{A}$, les
foncteurs
\[
  \Hom_{\mathcal{A}}(A, -) \colon \mathcal{A} \to \Ab
  \qquad \text{et} \qquad
  \Hom_{\mathcal{A}}(-, A) \colon \mathcal{A}^{\op} \to \Ab
\]
sont exacts à gauche.
\end{proposition}

\begin{proof}
Soit $0 \to B' \xrightarrow{f} B \xrightarrow{g} B'' \to 0$ une suite
exacte courte. On doit montrer que
\[
  0 \to \Hom(A, B') \xrightarrow{f_*} \Hom(A, B) \xrightarrow{g_*} \Hom(A, B'')
\]
est exacte.

\emph{Injectivité de $f_*$.}
Si $f \circ h = 0$ pour $h \colon A \to B'$, alors $h$ se factorise
par $\ker f = 0$, donc $h = 0$.

\emph{Exactitude en $\Hom(A, B)$.}
On a $g_* \circ f_* = (g \circ f)_* = 0$. Réciproquement, si
$g \circ \varphi = 0$ pour $\varphi \colon A \to B$, alors $\varphi$
se factorise par $\ker g = \im f$, donc il existe
$\psi \colon A \to B'$ avec $f \circ \psi = \varphi$.
L'argument pour $\Hom(-, A)$ est dual.
\end{proof}

\begin{exemple}\label{ex:tensor-right-exact}
\index{produit tensoriel!exactitude à droite}
Le foncteur $- \otimes_R N \colon R\text{-}\Mod \to \Ab$ est exact
à droite mais pas exact en général. La mesure de sa non-exactitude
à gauche est le foncteur $\operatorname{Tor}$\index{Tor@$\operatorname{Tor}$},
que nous étudierons au chapitre~\ref{ch:derived}.
\end{exemple}


%% ============================================================
\section{Le lemme du serpent}\label{sec:snake-lemma}
%% ============================================================

\begin{theoreme}[Lemme du serpent]\label{thm:snake-lemma}
\index{lemme!du serpent|textbf}
\index{morphisme connectant}
Soit $\mathcal{A}$ une catégorie abélienne. Étant donné un diagramme
commutatif à lignes exactes
\[
\begin{tikzcd}[column sep=large]
  & A' \arrow[r, "f'"] \arrow[d, "a"]
  & B' \arrow[r, "g'"] \arrow[d, "b"]
  & C' \arrow[r] \arrow[d, "c"]
  & 0 \\
0 \arrow[r]
  & A \arrow[r, "f"]
  & B \arrow[r, "g"]
  & C &
\end{tikzcd}
\]
il existe un morphisme connectant\index{morphisme connectant|textbf}
$\delta \colon \ker c \to \coker a$ et une suite exacte
\[
\begin{tikzcd}[column sep=small]
\ker a \arrow[r] & \ker b \arrow[r] & \ker c
  \arrow[r, "\delta"] & \coker a \arrow[r] & \coker b \arrow[r] & \coker c.
\end{tikzcd}
\]
De plus, si $A' \to B'$ est un monomorphisme (c'est-à-dire la ligne supérieure
est exacte aussi à gauche), alors $\ker a \to \ker b$ est un
monomorphisme. Dualement, si $B \to C$ est un épimorphisme, alors
$\coker b \to \coker c$ est un épimorphisme.
\end{theoreme}

\begin{proof}
Par le théorème de Freyd--Mitchell\index{théorème!de Freyd--Mitchell},
il suffit de démontrer ce résultat dans $R\text{-}\Mod$, où l'on peut
raisonner sur les éléments. Nous donnons néanmoins une preuve
diagrammatique qui peut être formalisée dans toute catégorie abélienne.

\medskip
\textbf{Étape 1 : Morphismes entre noyaux.}
\index{lemme!du serpent!construction}
Puisque $f \circ a \circ \ker a' = b \circ f' \circ \ker a' = b \circ f' \circ \ker a'$
et le carré commute ($b \circ f' = f \circ a$), la restriction de $b$ à
$\ker a$ donne un morphisme $\ker a \to \ker b$ (car
$b \circ (f' \circ \ker a') = f \circ (a \circ \ker a') = f \circ 0 = 0$,
et $f$ est un monomorphisme dans la ligne inférieure lorsque la suite est
exacte à gauche). Plus précisément, pour tout $x \in \ker a$, on a
$a(x) = 0$, donc $b(f'(x)) = f(a(x)) = 0$, ce qui signifie $f'(x) \in \ker b$.
De même, on construit $\ker b \to \ker c$.

\medskip
\textbf{Étape 2 : Construction du morphisme connectant $\delta$.}
Soit $z \in \ker c$. Comme $g'$ est surjectif, il existe $y \in B'$ tel que
$g'(y) = z$. Considérons $b(y) \in B$. On a
\[
  g(b(y)) = c(g'(y)) = c(z) = 0,
\]
donc $b(y) \in \ker g = \im f$, et il existe $x \in A$ (unique car $f$ est
injectif) tel que $f(x) = b(y)$. On pose
\[
  \delta(z) = \text{classe de } x \text{ dans } \coker a = A / \im a.
\]

\emph{Bonne définition.}
Si $y_1$ est un autre relèvement de $z$, alors $y - y_1 \in \ker g' = \im f'$,
donc $y - y_1 = f'(w)$ pour un $w \in A'$. Alors
$b(y) - b(y_1) = b(f'(w)) = f(a(w))$, d'où $x - x_1 = a(w) \in \im a$,
et les classes coïncident.

\medskip
\textbf{Étape 3 : Exactitude en $\ker c$.}
Soit $z \in \ker c$ dans l'image de $\ker b \to \ker c$, c'est-à-dire
$z = g'(y)$ avec $b(y) = 0$. Alors on peut prendre $x = 0$
(car $f(0) = 0 = b(y)$), donc $\delta(z) = 0$.

Réciproquement, si $\delta(z) = 0$, alors $x = a(w)$ pour un $w \in A'$.
Posons $y_0 = y - f'(w)$. Alors $g'(y_0) = g'(y) = z$ et
$b(y_0) = b(y) - b(f'(w)) = f(x) - f(a(w)) = f(x - x) = 0$,
donc $y_0 \in \ker b$ et $z$ provient de $\ker b$.

\medskip
\textbf{Étape 4 : Exactitude en $\coker a$.}
Soit $\bar{x} \in \coker a$ l'image d'un $x \in A$ par la projection.
Alors $\bar{x} \in \im \delta$ si et seulement si il existe $y \in B'$
avec $f(x) = b(y)$, c'est-à-dire $\bar{x}$ est dans le noyau de
$\coker a \to \coker b$ (car $f(x) = b(y)$ signifie que $f(x) \in \im b$).
D'où l'exactitude en $\coker a$.

Les vérifications de l'exactitude en $\ker a$, $\ker b$, $\coker b$ et
$\coker c$ sont similaires et laissées au lecteur (exercice~\ref{exo:snake-details}).
\end{proof}

Le diagramme complet du lemme du serpent est souvent représenté ainsi :

\[
\begin{tikzcd}[column sep=large, row sep=large]
  & \ker a \arrow[r] \arrow[d, hook]
  & \ker b \arrow[r] \arrow[d, hook]
  & \ker c \arrow[d, hook]
    \arrow[dll, "\delta"',
      rounded corners,
      to path={ -- ([xshift=2.5em]\tikztostart.east)
                |- ([yshift=-1.2em]\tikztostart.south) -| ([xshift=-2.5em]\tikztotarget.west)
                -- (\tikztotarget)}] \\
  & A' \arrow[r, "f'"] \arrow[d, "a"]
  & B' \arrow[r, "g'"] \arrow[d, "b"]
  & C' \arrow[r] \arrow[d, "c"]
  & 0 \\
0 \arrow[r]
  & A \arrow[r, "f"] \arrow[d, two heads]
  & B \arrow[r, "g"] \arrow[d, two heads]
  & C \arrow[d, two heads] \\
  & \coker a \arrow[r]
  & \coker b \arrow[r]
  & \coker c
\end{tikzcd}
\]


%% ============================================================
\section{Le lemme des cinq}\label{sec:five-lemma}
%% ============================================================

\begin{theoreme}[Lemme des cinq]\label{thm:five-lemma}
\index{lemme!des cinq|textbf}
\index{five lemma|see{lemme des cinq}}
Soit un diagramme commutatif dans une catégorie abélienne dont les lignes
sont exactes :
\[
\begin{tikzcd}
A_1 \arrow[r, "d_1"] \arrow[d, "\alpha_1"]
& A_2 \arrow[r, "d_2"] \arrow[d, "\alpha_2"]
& A_3 \arrow[r, "d_3"] \arrow[d, "\alpha_3"]
& A_4 \arrow[r, "d_4"] \arrow[d, "\alpha_4"]
& A_5 \arrow[d, "\alpha_5"] \\
B_1 \arrow[r, "e_1"]
& B_2 \arrow[r, "e_2"]
& B_3 \arrow[r, "e_3"]
& B_4 \arrow[r, "e_4"]
& B_5
\end{tikzcd}
\]
\begin{enumerate}[label=(\roman*)]
\item Si $\alpha_1$ est un épimorphisme et $\alpha_2, \alpha_4$ sont des
  monomorphismes, alors $\alpha_3$ est un monomorphisme.
\item Si $\alpha_5$ est un monomorphisme et $\alpha_2, \alpha_4$ sont des
  épimorphismes, alors $\alpha_3$ est un épimorphisme.
\item En particulier, si $\alpha_1, \alpha_2, \alpha_4, \alpha_5$ sont des
  isomorphismes, alors $\alpha_3$ est un isomorphisme.
\end{enumerate}
\end{theoreme}

\begin{proof}
Par le théorème de Freyd--Mitchell, il suffit de raisonner dans
$R\text{-}\Mod$.

\medskip
\textbf{(i) $\alpha_3$ est un monomorphisme.}
Soit $x \in A_3$ tel que $\alpha_3(x) = 0$. Alors
$e_3(\alpha_3(x)) = 0$, donc $\alpha_4(d_3(x)) = 0$
(par commutativité). Comme $\alpha_4$ est un monomorphisme,
$d_3(x) = 0$. Par exactitude de la ligne supérieure en $A_3$,
il existe $y \in A_2$ tel que $d_2(y) = x$.

Considérons $\alpha_2(y) \in B_2$. On a
$e_2(\alpha_2(y)) = \alpha_3(d_2(y)) = \alpha_3(x) = 0$.
Par exactitude de la ligne inférieure en $B_2$, il existe
$z \in B_1$ tel que $e_1(z) = \alpha_2(y)$.
Comme $\alpha_1$ est un épimorphisme, il existe $w \in A_1$
avec $\alpha_1(w) = z$. Alors
\[
  \alpha_2(d_1(w)) = e_1(\alpha_1(w)) = e_1(z) = \alpha_2(y).
\]
Comme $\alpha_2$ est un monomorphisme, $d_1(w) = y$, donc
$x = d_2(y) = d_2(d_1(w)) = 0$ par exactitude ($d_2 \circ d_1 = 0$
car $\im d_1 \subseteq \ker d_2$). Donc $\alpha_3$ est un monomorphisme.

\medskip
\textbf{(ii) $\alpha_3$ est un épimorphisme.}
Argument dual. Soit $b \in B_3$. On a
$\alpha_4^{-1}(e_3(b))$... plus précisément, $e_3(b) \in B_4$
et comme $\alpha_4$ est un épimorphisme, il existe $a_4 \in A_4$
avec $\alpha_4(a_4) = e_3(b)$.

Considérons $e_4(e_3(b)) = 0$ (par exactitude). Donc
$e_4(\alpha_4(a_4)) = 0$, ce qui donne $\alpha_5(d_4(a_4)) = 0$.
Comme $\alpha_5$ est un monomorphisme, $d_4(a_4) = 0$.
Par exactitude en $A_4$, il existe $a_3 \in A_3$ avec $d_3(a_3) = a_4$.

Posons $b' = b - \alpha_3(a_3)$. Alors
$e_3(b') = e_3(b) - e_3(\alpha_3(a_3)) = e_3(b) - \alpha_4(d_3(a_3))
= e_3(b) - \alpha_4(a_4) = 0$.
Par exactitude en $B_3$, il existe $b_2 \in B_2$ avec $e_2(b_2) = b'$.
Comme $\alpha_2$ est un épimorphisme, il existe $a_2 \in A_2$ avec
$\alpha_2(a_2) = b_2$. Alors
$\alpha_3(d_2(a_2)) = e_2(\alpha_2(a_2)) = e_2(b_2) = b'$.
D'où $b = \alpha_3(a_3) + \alpha_3(d_2(a_2)) = \alpha_3(a_3 + d_2(a_2))$,
et $\alpha_3$ est surjectif.
\end{proof}

\begin{corollaire}[Lemme court des cinq]\label{cor:short-five}
\index{lemme!des cinq!court}
Soit un diagramme commutatif à lignes exactes courtes :
\[
\begin{tikzcd}
0 \arrow[r] & A' \arrow[r] \arrow[d, "\alpha"] & B' \arrow[r] \arrow[d, "\beta"]
  & C' \arrow[r] \arrow[d, "\gamma"] & 0 \\
0 \arrow[r] & A \arrow[r] & B \arrow[r] & C \arrow[r] & 0
\end{tikzcd}
\]
Si $\alpha$ et $\gamma$ sont des isomorphismes, alors $\beta$ est un
isomorphisme.
\end{corollaire}

\begin{proof}
On complète par $0 \to 0$ à gauche et $0 \to 0$ à droite pour obtenir
un diagramme de cinq colonnes, puis on applique le lemme des cinq avec
$\alpha_1 = \alpha_5 = \id_0$, $\alpha_2 = \alpha$, $\alpha_4 = \gamma$.
\end{proof}


%% ============================================================
\section{Objets injectifs et projectifs}\label{sec:injective-projective}
%% ============================================================

\begin{definition}[Objet projectif]\label{def:projective-object}
\index{objet projectif|textbf}
Un objet $P$ d'une catégorie abélienne $\mathcal{A}$ est \textbf{projectif}
si le foncteur $\Hom_{\mathcal{A}}(P, -)$ est exact, c'est-à-dire si
pour tout épimorphisme $g \colon B \twoheadrightarrow C$ et tout morphisme
$f \colon P \to C$, il existe $\tilde{f} \colon P \to B$ tel que
$g \circ \tilde{f} = f$ :
\[
\begin{tikzcd}
& P \arrow[d, "f"] \arrow[dl, dashed, "\tilde{f}"'] \\
B \arrow[r, two heads, "g"] & C \arrow[r] & 0
\end{tikzcd}
\]
\end{definition}

\begin{definition}[Objet injectif]\label{def:injective-object}
\index{objet injectif|textbf}
Dualement, un objet $I$ est \textbf{injectif} si $\Hom_{\mathcal{A}}(-, I)$
est exact, c'est-à-dire si pour tout monomorphisme $f \colon A \hookrightarrow B$
et tout morphisme $g \colon A \to I$, il existe $\tilde{g} \colon B \to I$
avec $\tilde{g} \circ f = g$ :
\[
\begin{tikzcd}
0 \arrow[r] & A \arrow[r, hook, "f"] \arrow[d, "g"'] & B \arrow[dl, dashed, "\tilde{g}"] \\
& I &
\end{tikzcd}
\]
\end{definition}

\begin{proposition}\label{prop:free-projective}
\index{module libre!est projectif}
\index{objet projectif!modules libres}
Tout $R$-module libre est projectif. Plus généralement, les
facteurs directs de modules libres (modules projectifs) sont projectifs.
\end{proposition}

\begin{proof}
Soit $F = \bigoplus_{i \in I} R$ un module libre. Un morphisme
$f \colon F \to C$ est déterminé par les images $f(e_i) \in C$ des
générateurs. Si $g \colon B \twoheadrightarrow C$ est surjectif, on
choisit pour chaque $i$ un $b_i \in B$ avec $g(b_i) = f(e_i)$.
La propriété universelle du module libre donne un unique
$\tilde{f} \colon F \to B$ avec $\tilde{f}(e_i) = b_i$, et
$g \circ \tilde{f} = f$.
\end{proof}

\begin{theoreme}[Critère de Baer]\label{thm:baer-criterion}
\index{critère de Baer|textbf}
\index{Baer, Reinhold}
Un $R$-module $I$ est injectif si et seulement si pour tout
idéal $J \subseteq R$ et tout morphisme $f \colon J \to I$,
il existe une extension $\tilde{f} \colon R \to I$ avec
$\tilde{f}|_J = f$.
\end{theoreme}

\begin{proof}
La nécessité est immédiate (prendre le monomorphisme $J \hookrightarrow R$).
Pour la suffisance, soit $f \colon A \hookrightarrow B$ un monomorphisme
et $g \colon A \to I$ un morphisme. On utilise le lemme de Zorn sur
l'ensemble des extensions partielles de $g$ le long de sous-modules
de $B$ contenant $A$. Un élément maximal $(M, h)$ doit avoir $M = B$ :
sinon, on prend $b \in B \setminus M$ et on pose
$J = \{r \in R : rb \in M\}$. La composée
$J \xrightarrow{r \mapsto rb} M \xrightarrow{h} I$
s'étend à $R$ par hypothèse, ce qui permet d'étendre $h$ à $M + Rb$,
contredisant la maximalité.
\end{proof}

\begin{definition}[Assez d'injectifs / de projectifs]\label{def:enough-injectives}
\index{assez d'injectifs|textbf}
\index{assez de projectifs|textbf}
Une catégorie abélienne $\mathcal{A}$ a \textbf{assez d'injectifs}
si tout objet $A$ admet un monomorphisme $A \hookrightarrow I$ avec
$I$ injectif.

Dualement, $\mathcal{A}$ a \textbf{assez de projectifs} si tout
objet $A$ admet un épimorphisme $P \twoheadrightarrow A$ avec $P$
projectif.
\end{definition}

\begin{theoreme}\label{thm:R-mod-enough}
\index{R-Mod@$R\text{-}\Mod$!assez d'injectifs}
\index{R-Mod@$R\text{-}\Mod$!assez de projectifs}
Pour tout anneau $R$, la catégorie $R\text{-}\Mod$ a assez d'injectifs
et assez de projectifs.
\end{theoreme}

\begin{proof}
\emph{Assez de projectifs.} Tout module $M$ est quotient d'un module
libre (prendre le module libre sur un ensemble de générateurs de $M$),
et les modules libres sont projectifs.

\emph{Assez d'injectifs.} On utilise le plongement de Baer. Tout
groupe abélien $A$ se plonge dans un groupe abélien divisible\index{groupe abélien!divisible}
(donc injectif en tant que $\Z$-module). Pour un $R$-module $M$,
on plonge le $\Z$-module sous-jacent dans un $\Z$-module injectif $Q$,
puis on montre que $\Hom_\Z(R, Q)$ est un $R$-module injectif
et que l'application $M \to \Hom_\Z(R, Q)$ définie par
$m \mapsto (r \mapsto rm \in Q)$ est un monomorphisme de $R$-modules.
\end{proof}


%% ============================================================
\section{Résolutions}\label{sec:resolutions}
%% ============================================================

\begin{definition}[Résolution projective]\label{def:projective-resolution}
\index{résolution!projective|textbf}
Soit $A$ un objet d'une catégorie abélienne $\mathcal{A}$.
Une \textbf{résolution projective} de $A$ est une suite exacte
\[
\begin{tikzcd}
\cdots \arrow[r] & P_2 \arrow[r, "d_2"] & P_1 \arrow[r, "d_1"]
  & P_0 \arrow[r, "\varepsilon"] & A \arrow[r] & 0
\end{tikzcd}
\]
où chaque $P_n$ est projectif. On note $P_\bullet \to A$ une telle résolution.
\end{definition}

\begin{definition}[Résolution injective]\label{def:injective-resolution}
\index{résolution!injective|textbf}
Dualement, une \textbf{résolution injective} de $A$ est une suite exacte
\[
\begin{tikzcd}
0 \arrow[r] & A \arrow[r, "\eta"] & I^0 \arrow[r, "d^0"]
  & I^1 \arrow[r, "d^1"] & I^2 \arrow[r] & \cdots
\end{tikzcd}
\]
où chaque $I^n$ est injectif. On note $A \to I^\bullet$.
\end{definition}

\begin{proposition}\label{prop:resolution-existence}
\index{résolution!existence}
Si $\mathcal{A}$ a assez de projectifs (resp.\ d'injectifs), tout
objet admet une résolution projective (resp.\ injective).
\end{proposition}

\begin{proof}
Pour une résolution projective : on choisit un épimorphisme
$P_0 \twoheadrightarrow A$ avec $P_0$ projectif, puis un épimorphisme
$P_1 \twoheadrightarrow \ker(P_0 \to A)$ avec $P_1$ projectif, et
ainsi de suite par récurrence. L'argument pour les résolutions
injectives est dual.
\end{proof}

\begin{theoreme}[Lemme de comparaison]\label{thm:comparison-lemma}
\index{lemme!de comparaison|textbf}
Soit $f \colon A \to B$ un morphisme, $P_\bullet \to A$ une résolution
projective et $Q_\bullet \to B$ une suite exacte quelconque (pas
nécessairement projective). Alors il existe un morphisme de chaînes
$\varphi_\bullet \colon P_\bullet \to Q_\bullet$ relevant $f$, et
deux tels relèvements sont homotopes\index{homotopie de chaînes}.
\end{theoreme}

\begin{proof}
On construit $\varphi_n \colon P_n \to Q_n$ par récurrence.
Pour $n = 0$ : on a $\varepsilon_B \circ \varphi_0 = f \circ \varepsilon_A$,
et comme $P_0$ est projectif et $\varepsilon_B$ est surjectif, le
relèvement existe. Pour le pas de récurrence : on utilise la
projectivité de $P_n$ et l'exactitude de $Q_\bullet$ pour construire
$\varphi_n$ relevant le morphisme induit sur les noyaux.

L'unicité à homotopie près se démontre de manière similaire, en
construisant une homotopie $s_n \colon P_n \to Q_{n+1}$ par
récurrence, utilisant à chaque étape la projectivité de $P_n$.
\end{proof}


%% ============================================================
\section{Exercices}\label{sec:exos-abelian}
%% ============================================================

\begin{exercice}\label{exo:additive-functor-zero}
\index{foncteur!additif!préserve l'objet nul}
Montrer qu'un foncteur additif $F \colon \mathcal{A} \to \mathcal{B}$
entre catégories additives préserve l'objet nul et les biproduits.
\end{exercice}

\begin{exercice}\label{exo:abelian-mono-epi}
\index{monomorphisme!dans une catégorie abélienne}
Soit $\mathcal{A}$ une catégorie abélienne et $f \colon A \to B$ un
morphisme. Montrer que :
\begin{enumerate}[label=(\alph*)]
\item $f$ est un monomorphisme si et seulement si $\ker f = 0$ ;
\item $f$ est un épimorphisme si et seulement si $\coker f = 0$ ;
\item $f$ est un isomorphisme si et seulement si $\ker f = 0$ et
  $\coker f = 0$.
\end{enumerate}
\end{exercice}

\begin{exercice}\label{exo:snake-details}
\index{lemme!du serpent!détails}
Compléter les détails de la preuve du lemme du serpent
(théorème~\ref{thm:snake-lemma}) en vérifiant l'exactitude de la suite
$\ker a \to \ker b \to \ker c$ et $\coker a \to \coker b \to \coker c$.
\end{exercice}

\begin{exercice}\label{exo:horseshoe}
\index{lemme!du fer à cheval}
\textbf{(Lemme du fer à cheval.)}
Soit $0 \to A' \to A \to A'' \to 0$ une suite exacte courte dans
une catégorie abélienne ayant assez de projectifs. Si $P'_\bullet \to A'$
et $P''_\bullet \to A''$ sont des résolutions projectives, montrer qu'il
existe une résolution projective $P_\bullet \to A$ avec
$P_n = P'_n \oplus P''_n$ qui s'insère dans une suite exacte courte
de complexes $0 \to P'_\bullet \to P_\bullet \to P''_\bullet \to 0$.
\end{exercice}

\begin{exercice}\label{exo:ses-projective}
\index{suite exacte!courte!scindée!caractérisation projective}
Soit $0 \to A \xrightarrow{f} B \xrightarrow{g} C \to 0$ une suite
exacte courte dans $R\text{-}\Mod$. Montrer que les conditions
suivantes sont équivalentes :
\begin{enumerate}[label=(\roman*)]
\item la suite est scindée ;
\item $g$ admet une section ;
\item $f$ admet une rétraction ;
\item pour tout module $M$, la suite
  $0 \to \Hom(M, A) \to \Hom(M, B) \to \Hom(M, C) \to 0$ est exacte.
\end{enumerate}
En déduire qu'un module $C$ est projectif si et seulement si toute suite
exacte courte $0 \to A \to B \to C \to 0$ est scindée.
\end{exercice}

\begin{exercice}\label{exo:injective-divisible}
\index{groupe abélien!divisible}
\index{objet injectif!groupes abéliens divisibles}
Un groupe abélien $A$ est dit \textbf{divisible}\index{divisible} si pour
tout $n \geq 1$ et tout $a \in A$, il existe $b \in A$ avec $nb = a$.
\begin{enumerate}[label=(\alph*)]
\item Montrer que $\Q$ et $\Q/\Z$ sont divisibles.
\item Montrer qu'un groupe abélien est injectif (dans $\Ab$) si et seulement
  s'il est divisible.
  \textit{Indication : utiliser le critère de Baer.}
\item En déduire que $\Ab$ a assez d'injectifs.
\end{enumerate}
\end{exercice}

\begin{exercice}\label{exo:nine-lemma}
\index{lemme!des neuf}
\textbf{(Lemme $3 \times 3$.)}
Soit un diagramme commutatif dans une catégorie abélienne dont les
colonnes et les deux premières lignes sont exactes :
\[
\begin{tikzcd}
& 0 \arrow[d] & 0 \arrow[d] & 0 \arrow[d] \\
0 \arrow[r] & A' \arrow[r] \arrow[d] & B' \arrow[r] \arrow[d] & C' \arrow[r] \arrow[d] & 0 \\
0 \arrow[r] & A \arrow[r] \arrow[d] & B \arrow[r] \arrow[d] & C \arrow[r] \arrow[d] & 0 \\
0 \arrow[r] & A'' \arrow[r] \arrow[d] & B'' \arrow[r] \arrow[d] & C'' \arrow[r] \arrow[d] & 0 \\
& 0 & 0 & 0
\end{tikzcd}
\]
Montrer que la troisième ligne est également exacte.
\textit{Indication : appliquer deux fois le lemme du serpent.}
\end{exercice}

\begin{exercice}\label{exo:grothendieck-AB5}
\index{axiome AB5}
\index{catégorie!de Grothendieck}
Soit $\mathcal{A}$ une catégorie abélienne. On dit que $\mathcal{A}$
satisfait l'axiome \textbf{AB5}\index{AB5} si $\mathcal{A}$ admet
toutes les colimites filtrantes et si les colimites filtrantes
commutent aux suites exactes. Montrer que $R\text{-}\Mod$ satisfait AB5.
\end{exercice}

\index{suite exacte|)}
\index{catégorie!abélienne|)}


%%%%%%%%%%%%%%%%%%%%%%%%%%%%%%%%%%%%%%%%%%%%%%%%%%%%%%%%%%%%
%%           CHAPITRE 8 : CATÉGORIES DÉRIVÉES             %%
%%                   — INTRODUCTION                        %%
%%%%%%%%%%%%%%%%%%%%%%%%%%%%%%%%%%%%%%%%%%%%%%%%%%%%%%%%%%%%

\chapter{Catégories dérivées --- Introduction}\label{ch:derived}
\minitoc

\vspace{1em}

Les catégories dérivées\index{catégorie!dérivée|(}, introduites par
Grothendieck\index{Grothendieck, Alexander} et développées par
Verdier\index{Verdier, Jean-Louis} dans sa thèse (1967), fournissent le
cadre conceptuel adéquat pour l'algèbre homologique. L'idée fondatrice est
simple mais profonde : au lieu de travailler avec un objet $A$ et ses
résolutions, on travaille directement avec la catégorie des complexes de
chaînes, en « inversant » les quasi-isomorphismes. Ce processus de
\emph{localisation}\index{localisation!de catégories} produit la catégorie
dérivée $D(\mathcal{A})$, dans laquelle un objet et sa résolution deviennent
isomorphes. La structure supplémentaire de \emph{catégorie
triangulée}\index{catégorie!triangulée} encode l'information des suites
exactes longues.

\index{complexe de chaînes|(}

%% ============================================================
\section{Complexes de chaînes}\label{sec:chain-complexes}
%% ============================================================

\begin{definition}[Complexe de chaînes]\label{def:chain-complex}
\index{complexe de chaînes|textbf}
\index{différentielle}
Soit $\mathcal{A}$ une catégorie abélienne. Un \textbf{complexe de chaînes}
(ou \textbf{complexe cochain}) dans $\mathcal{A}$ est une suite
$(A^n, d^n)_{n \in \Z}$ d'objets et de morphismes
\[
\begin{tikzcd}
\cdots \arrow[r] & A^{n-1} \arrow[r, "d^{n-1}"] & A^n \arrow[r, "d^n"]
  & A^{n+1} \arrow[r] & \cdots
\end{tikzcd}
\]
telle que $d^{n} \circ d^{n-1} = 0$ pour tout $n \in \Z$.
Les morphismes $d^n$ sont appelés les \textbf{différentielles}\index{différentielle|textbf}.
On note le complexe $A^\bullet$ ou $(A^\bullet, d)$.
\end{definition}

\begin{definition}[Morphisme de complexes]\label{def:chain-map}
\index{morphisme!de complexes|textbf}
Un \textbf{morphisme de complexes} $f \colon A^\bullet \to B^\bullet$
est une famille de morphismes $(f^n \colon A^n \to B^n)_{n \in \Z}$
telle que le diagramme suivant commute pour tout $n$ :
\[
\begin{tikzcd}
A^n \arrow[r, "d_A^n"] \arrow[d, "f^n"'] & A^{n+1} \arrow[d, "f^{n+1}"] \\
B^n \arrow[r, "d_B^n"] & B^{n+1}
\end{tikzcd}
\]
c'est-à-dire $f^{n+1} \circ d_A^n = d_B^n \circ f^n$.
\end{definition}

\begin{notation}\label{not:Ch}
\index{Ch(A)@$\mathsf{Ch}(\mathcal{A})$}
On note $\mathsf{Ch}(\mathcal{A})$ la catégorie dont les objets sont les
complexes de chaînes dans $\mathcal{A}$ et les morphismes sont les
morphismes de complexes. On note aussi :
\begin{itemize}
\item $\mathsf{Ch}^+(\mathcal{A})$\index{Ch+@$\mathsf{Ch}^+(\mathcal{A})$} :
  complexes bornés inférieurement ($A^n = 0$ pour $n \ll 0$) ;
\item $\mathsf{Ch}^-(\mathcal{A})$\index{Ch-@$\mathsf{Ch}^-(\mathcal{A})$} :
  complexes bornés supérieurement ($A^n = 0$ pour $n \gg 0$) ;
\item $\mathsf{Ch}^b(\mathcal{A})$\index{Chb@$\mathsf{Ch}^b(\mathcal{A})$} :
  complexes bornés ($A^n = 0$ pour $|n| \gg 0$).
\end{itemize}
\end{notation}

\begin{definition}[Cohomologie d'un complexe]\label{def:cohomology}
\index{cohomologie!d'un complexe|textbf}
\index{cocycle}
\index{cobord}
Le $n$-ième \textbf{objet de cohomologie}\index{cohomologie|textbf}
du complexe $A^\bullet$ est
\[
  H^n(A^\bullet) = \ker d^n / \im d^{n-1}
  = \ker(d^n \colon A^n \to A^{n+1}) \big/ \im(d^{n-1} \colon A^{n-1} \to A^n).
\]
Ceci a un sens car $\im d^{n-1} \subseteq \ker d^n$ (puisque $d^n \circ d^{n-1} = 0$).
Les éléments de $\ker d^n$ sont les \textbf{$n$-cocycles}\index{cocycle|textbf}
et ceux de $\im d^{n-1}$ les \textbf{$n$-cobords}\index{cobord|textbf}.
\end{definition}

\begin{proposition}\label{prop:cohomology-functor}
\index{cohomologie!foncteur}
La cohomologie définit un foncteur
$H^n \colon \mathsf{Ch}(\mathcal{A}) \to \mathcal{A}$
pour chaque $n \in \Z$.
\end{proposition}

\begin{proof}
Soit $f \colon A^\bullet \to B^\bullet$ un morphisme de complexes.
Comme $f^n$ envoie $\ker d_A^n$ dans $\ker d_B^n$ (par commutativité)
et $\im d_A^{n-1}$ dans $\im d_B^{n-1}$, il induit un morphisme
$H^n(f) \colon H^n(A^\bullet) \to H^n(B^\bullet)$. La fonctorialité
($H^n(\id) = \id$ et $H^n(g \circ f) = H^n(g) \circ H^n(f)$) est
immédiate.
\end{proof}

\begin{proposition}\label{prop:Ch-abelian}
\index{Ch(A)@$\mathsf{Ch}(\mathcal{A})$!catégorie abélienne}
Si $\mathcal{A}$ est abélienne, alors $\mathsf{Ch}(\mathcal{A})$ est
abélienne, les noyaux, conoyaux et (co)images étant calculés degré
par degré.
\end{proposition}

\begin{proof}
On vérifie les axiomes. L'objet nul est le complexe nul.
Le biproduit est $(A \oplus B)^n = A^n \oplus B^n$ avec
$d_{A \oplus B}^n = d_A^n \oplus d_B^n$.
Le noyau de $f \colon A^\bullet \to B^\bullet$ est le complexe
$(\ker f^n, d^n|_{\ker f^n})_{n \in \Z}$, qui est bien un sous-complexe
car $d^n$ envoie $\ker f^n$ dans $\ker f^{n+1}$. De même pour le
conoyau. La condition image $=$ coimage se vérifie degré par degré
car elle est satisfaite dans $\mathcal{A}$.
\end{proof}


%% ============================================================
\section{Homotopie de chaînes et catégorie homotopique}%
\label{sec:chain-homotopy}
%% ============================================================

\begin{definition}[Homotopie de chaînes]\label{def:chain-homotopy}
\index{homotopie de chaînes|textbf}
Soient $f, g \colon A^\bullet \to B^\bullet$ deux morphismes de complexes.
Une \textbf{homotopie de chaînes} de $f$ vers $g$ est une famille de
morphismes $(s^n \colon A^n \to B^{n-1})_{n \in \Z}$ telle que
\[
  f^n - g^n = d_B^{n-1} \circ s^n + s^{n+1} \circ d_A^n
  \qquad \text{pour tout } n \in \Z.
\]
\[
\begin{tikzcd}
\cdots \arrow[r]
& A^{n-1} \arrow[r, "d_A^{n-1}"] \arrow[d]
& A^n \arrow[r, "d_A^n"] \arrow[d] \arrow[dl, dashed, "s^n"']
& A^{n+1} \arrow[r] \arrow[d] \arrow[dl, dashed, "s^{n+1}"']
& \cdots \\
\cdots \arrow[r]
& B^{n-1} \arrow[r, "d_B^{n-1}"]
& B^n \arrow[r, "d_B^n"]
& B^{n+1} \arrow[r]
& \cdots
\end{tikzcd}
\]
On dit que $f$ et $g$ sont \textbf{homotopes}\index{homotope} et on
note $f \sim g$.
\end{definition}

\begin{proposition}\label{prop:homotopy-equivalence}
\index{homotopie de chaînes!relation d'équivalence}
La relation d'homotopie $\sim$ est une relation d'équivalence sur
$\Hom_{\mathsf{Ch}(\mathcal{A})}(A^\bullet, B^\bullet)$, compatible
avec la composition.
\end{proposition}

\begin{proof}
\emph{Réflexivité :} $f \sim f$ avec $s^n = 0$.
\emph{Symétrie :} si $s$ est une homotopie de $f$ vers $g$, alors $-s$ est
une homotopie de $g$ vers $f$.
\emph{Transitivité :} si $s$ est une homotopie de $f$ vers $g$ et $t$ de
$g$ vers $h$, alors $s + t$ est une homotopie de $f$ vers $h$.

\emph{Compatibilité.} Si $f \sim g \colon A^\bullet \to B^\bullet$
via $s$ et $h \colon B^\bullet \to C^\bullet$, alors $h \circ f \sim h \circ g$
via $h \circ s$. Si $h \colon C^\bullet \to A^\bullet$, alors
$f \circ h \sim g \circ h$ via $s \circ h$.
\end{proof}

\begin{proposition}\label{prop:homotopic-maps-cohomology}
\index{homotopie de chaînes!et cohomologie}
Si $f \sim g \colon A^\bullet \to B^\bullet$, alors
$H^n(f) = H^n(g)$ pour tout $n$.
\end{proposition}

\begin{proof}
Soit $[a] \in H^n(A^\bullet)$, c'est-à-dire $a \in \ker d_A^n$.
Alors
\[
  f^n(a) - g^n(a) = d_B^{n-1}(s^n(a)) + s^{n+1}(d_A^n(a))
  = d_B^{n-1}(s^n(a)) + 0 \in \im d_B^{n-1},
\]
donc $[f^n(a)] = [g^n(a)]$ dans $H^n(B^\bullet)$.
\end{proof}

\begin{definition}[Catégorie homotopique]\label{def:homotopy-category}
\index{catégorie!homotopique|textbf}
\index{K(A)@$K(\mathcal{A})$}
La \textbf{catégorie homotopique} $K(\mathcal{A})$ est la catégorie dont :
\begin{itemize}
\item les objets sont les complexes de chaînes dans $\mathcal{A}$ ;
\item les morphismes sont les classes d'homotopie :
  $\Hom_{K(\mathcal{A})}(A^\bullet, B^\bullet)
  = \Hom_{\mathsf{Ch}(\mathcal{A})}(A^\bullet, B^\bullet) / {\sim}$.
\end{itemize}
On définit de même $K^+(\mathcal{A})$, $K^-(\mathcal{A})$,
$K^b(\mathcal{A})$\index{K+@$K^+(\mathcal{A})$}%
\index{K-@$K^-(\mathcal{A})$}\index{Kb@$K^b(\mathcal{A})$}.
\end{definition}

\begin{definition}[Équivalence d'homotopie]\label{def:homotopy-equivalence}
\index{équivalence d'homotopie|textbf}
Un morphisme de complexes $f \colon A^\bullet \to B^\bullet$ est une
\textbf{équivalence d'homotopie} s'il devient un isomorphisme dans
$K(\mathcal{A})$, c'est-à-dire s'il existe $g \colon B^\bullet \to A^\bullet$
avec $g \circ f \sim \id_{A^\bullet}$ et $f \circ g \sim \id_{B^\bullet}$.
\end{definition}


%% ============================================================
\section{Quasi-isomorphismes et localisation}\label{sec:quasi-iso}
%% ============================================================

\begin{definition}[Quasi-isomorphisme]\label{def:quasi-iso}
\index{quasi-isomorphisme|textbf}
Un morphisme de complexes $f \colon A^\bullet \to B^\bullet$ est un
\textbf{quasi-isomorphisme} si
$H^n(f) \colon H^n(A^\bullet) \xrightarrow{\;\sim\;} H^n(B^\bullet)$
est un isomorphisme pour tout $n \in \Z$.
\end{definition}

\begin{remarque}\label{rem:htpy-eq-is-qiso}
\index{équivalence d'homotopie!est un quasi-isomorphisme}
Toute équivalence d'homotopie est un quasi-isomorphisme (par la
proposition~\ref{prop:homotopic-maps-cohomology}), mais la réciproque
est fausse en général.
\end{remarque}

\begin{exemple}\label{ex:quasi-iso-not-htpy}
\index{quasi-isomorphisme!qui n'est pas une équivalence d'homotopie}
Dans $\Ab$, considérons les complexes
\[
  A^\bullet : \quad 0 \to \Z \xrightarrow{2} \Z \to 0
  \qquad \text{et} \qquad
  B^\bullet : \quad 0 \to 0 \to \Z/2\Z \to 0
\]
(concentrés en degrés 0 et 1). Le morphisme $f \colon A^\bullet \to B^\bullet$
donné par $f^0 = 0$ et $f^1 = \text{projection}$ est un quasi-isomorphisme
($H^0 = 0$, $H^1 = \Z/2\Z$ pour les deux), mais \emph{pas} une
équivalence d'homotopie.
\end{exemple}

\begin{definition}[Localisation d'une catégorie]\label{def:localization}
\index{localisation!de catégories|textbf}
\index{Gabriel, Pierre}
\index{Zisman, Michel}
Soit $\mathcal{C}$ une catégorie et $S$ une classe de morphismes.
La \textbf{localisation} $\mathcal{C}[S^{-1}]$ est une catégorie munie
d'un foncteur $Q \colon \mathcal{C} \to \mathcal{C}[S^{-1}]$ tel que :
\begin{enumerate}[label=(\roman*)]
\item $Q(s)$ est un isomorphisme pour tout $s \in S$ ;
\item $Q$ est universel pour cette propriété : pour tout foncteur
  $F \colon \mathcal{C} \to \mathcal{D}$ tel que $F(s)$ est un
  isomorphisme pour tout $s \in S$, il existe un unique foncteur
  $\bar{F} \colon \mathcal{C}[S^{-1}] \to \mathcal{D}$ avec
  $\bar{F} \circ Q = F$.
\end{enumerate}
\end{definition}

\begin{remarque}\label{rem:localization-existence}
\index{localisation!existence}
La localisation existe toujours (modulo des problèmes ensemblistes),
mais les morphismes dans $\mathcal{C}[S^{-1}]$ sont en général
des « fractions » formelles\index{fraction!formelle}
$s_1^{-1} \circ f_1 \circ s_2^{-1} \circ f_2 \circ \cdots$,
ce qui rend la catégorie difficile à manipuler. La notion de
\emph{système multiplicatif}\index{système multiplicatif} permet de
simplifier cette description.
\end{remarque}

\begin{definition}[Système multiplicatif]\label{def:multiplicative-system}
\index{système multiplicatif|textbf}
\index{conditions d'Ore}
Une classe $S$ de morphismes dans une catégorie $\mathcal{C}$ est un
\textbf{système multiplicatif} si :
\begin{enumerate}[label=(MS\arabic*)]
\item\label{MS1} $S$ contient toutes les identités et est stable par
  composition ;
\item\label{MS2} \textbf{(Condition d'Ore à gauche et à droite.)}
  Pour tout diagramme $A \xrightarrow{f} C \xleftarrow{s} B$ avec
  $s \in S$, il existe un diagramme $A \xleftarrow{t} D \xrightarrow{g} B$
  avec $t \in S$ et $f \circ t = s \circ g$ (et dualement) :
\[
\begin{tikzcd}
D \arrow[r, "g"] \arrow[d, "t"', "{\in S}"] & B \arrow[d, "s", "{\in S}"'] \\
A \arrow[r, "f"] & C
\end{tikzcd}
\]
\item\label{MS3} Pour $f, g \colon A \to B$, il existe $s \in S$ avec
  $s \circ f = s \circ g$ si et seulement s'il existe $t \in S$ avec
  $f \circ t = g \circ t$.
\end{enumerate}
\end{definition}

\begin{proposition}\label{prop:qiso-multiplicative}
\index{quasi-isomorphisme!système multiplicatif}
La classe des quasi-isomorphismes dans $K(\mathcal{A})$ forme un
système multiplicatif.
\end{proposition}

\begin{proof}
La vérification de \ref{MS1} est immédiate.
Pour \ref{MS2} et \ref{MS3}, on utilise la construction du
\emph{cône de mapping}\index{cône!de mapping} (voir ci-dessous) et
les propriétés de la catégorie homotopique. Nous omettons les détails
techniques.
\end{proof}


%% ============================================================
\section{La catégorie dérivée}\label{sec:derived-category}
%% ============================================================

\begin{definition}[Catégorie dérivée]\label{def:derived-category}
\index{catégorie!dérivée|textbf}
\index{D(A)@$D(\mathcal{A})$}
La \textbf{catégorie dérivée} $D(\mathcal{A})$ est la localisation de la
catégorie homotopique $K(\mathcal{A})$ par rapport aux quasi-isomorphismes :
\[
  D(\mathcal{A}) = K(\mathcal{A})[\mathsf{Qis}^{-1}].
\]
On définit de même les variantes bornées :
\begin{itemize}
\item $D^+(\mathcal{A}) = K^+(\mathcal{A})[\mathsf{Qis}^{-1}]$\index{D+@$D^+(\mathcal{A})$} ;
\item $D^-(\mathcal{A}) = K^-(\mathcal{A})[\mathsf{Qis}^{-1}]$\index{D-@$D^-(\mathcal{A})$} ;
\item $D^b(\mathcal{A}) = K^b(\mathcal{A})[\mathsf{Qis}^{-1}]$\index{Db@$D^b(\mathcal{A})$}.
\end{itemize}
\end{definition}

\begin{remarque}\label{rem:derived-category-morphisms}
\index{catégorie!dérivée!morphismes}
\index{toit}
Un morphisme $A^\bullet \to B^\bullet$ dans $D(\mathcal{A})$ est
représenté par un \textbf{toit}\index{toit|textbf} :
\[
\begin{tikzcd}
& C^\bullet \arrow[dl, "s"', "\sim"] \arrow[dr, "f"] \\
A^\bullet & & B^\bullet
\end{tikzcd}
\]
où $s$ est un quasi-isomorphisme. La composition de deux toits
nécessite la condition d'Ore (MS2).
\end{remarque}

\begin{theoreme}\label{thm:derived-via-injectives}
\index{catégorie!dérivée!et résolutions injectives}
Si $\mathcal{A}$ a assez d'injectifs, le foncteur naturel
\[
  K^+(\mathcal{I}) \longrightarrow D^+(\mathcal{A})
\]
est une équivalence de catégories, où $\mathcal{I}$ désigne la
sous-catégorie pleine des objets injectifs de $\mathcal{A}$.
Autrement dit, tout objet de $D^+(\mathcal{A})$ est isomorphe à
un complexe d'injectifs.
\end{theoreme}

\begin{proof}
\emph{Essentielle surjectivité.} Soit $A^\bullet \in D^+(\mathcal{A})$.
Par récurrence, on construit une résolution injective
$I^\bullet$ du complexe $A^\bullet$ : c'est un complexe d'injectifs
muni d'un quasi-isomorphisme $A^\bullet \to I^\bullet$. Cette
construction généralise les résolutions injectives d'un seul objet
(proposition~\ref{prop:resolution-existence}).

\emph{Pleine fidélité.} Il faut montrer que si $I^\bullet, J^\bullet$
sont des complexes bornés inférieurement d'injectifs, alors
\[
  \Hom_{K^+(\mathcal{A})}(I^\bullet, J^\bullet)
  \xrightarrow{\;\sim\;} \Hom_{D^+(\mathcal{A})}(I^\bullet, J^\bullet).
\]
L'injectivité vient du fait qu'un morphisme homotope à zéro vers un
complexe d'injectifs est nul dans la catégorie dérivée. La surjectivité
utilise le fait que tout toit peut être « redressé » en utilisant
l'injectivité des composantes.
\end{proof}


%% ============================================================
\section{Le foncteur de décalage et le cône de mapping}%
\label{sec:shift-cone}
%% ============================================================

\begin{definition}[Foncteur de décalage]\label{def:shift-functor}
\index{foncteur!de décalage|textbf}
\index{décalage}
\index{translation}
Le \textbf{foncteur de décalage} (ou \textbf{translation})
$[1] \colon \mathsf{Ch}(\mathcal{A}) \to \mathsf{Ch}(\mathcal{A})$
est défini par
\[
  (A[1])^n = A^{n+1}, \qquad d_{A[1]}^n = -d_A^{n+1}.
\]
Pour un morphisme $f \colon A^\bullet \to B^\bullet$, on pose
$(f[1])^n = f^{n+1}$. Plus généralement, $A[k]^n = A^{n+k}$ avec
$d_{A[k]}^n = (-1)^k d_A^{n+k}$.
\end{definition}

\begin{definition}[Cône de mapping]\label{def:mapping-cone}
\index{cône!de mapping|textbf}
\index{mapping cone|see{cône de mapping}}
Soit $f \colon A^\bullet \to B^\bullet$ un morphisme de complexes.
Le \textbf{cône de mapping} de $f$, noté $\operatorname{Cone}(f)$,
est le complexe défini par
\[
  \operatorname{Cone}(f)^n = A^{n+1} \oplus B^n
\]
avec la différentielle
\[
  d_{\operatorname{Cone}(f)}^n =
  \begin{pmatrix} -d_A^{n+1} & 0 \\ f^{n+1} & d_B^n \end{pmatrix}
  \colon A^{n+1} \oplus B^n \longrightarrow A^{n+2} \oplus B^{n+1}.
\]
\end{definition}

\begin{proposition}\label{prop:cone-exact-triangle}
\index{cône!de mapping!suite exacte}
\index{triangle!distingué}
Pour tout morphisme de complexes $f \colon A^\bullet \to B^\bullet$,
on a une suite de morphismes de complexes
\[
\begin{tikzcd}
A^\bullet \arrow[r, "f"] & B^\bullet \arrow[r, "\iota"]
  & \operatorname{Cone}(f) \arrow[r, "\pi"] & A[1]^\bullet
\end{tikzcd}
\]
où $\iota^n(b) = (0, b)$ et $\pi^n(a, b) = a$. De plus, cette suite
induit une suite exacte longue en cohomologie :
\[
\begin{tikzcd}[column sep=small]
\cdots \arrow[r]
& H^n(A^\bullet) \arrow[r, "H^n(f)"]
& H^n(B^\bullet) \arrow[r]
& H^n(\operatorname{Cone}(f)) \arrow[r]
& H^{n+1}(A^\bullet) \arrow[r]
& \cdots
\end{tikzcd}
\]
\end{proposition}

\begin{proof}
On vérifie d'abord que $d_{\operatorname{Cone}(f)}^2 = 0$ :
\[
  \begin{pmatrix} -d_A^{n+2} & 0 \\ f^{n+2} & d_B^{n+1} \end{pmatrix}
  \begin{pmatrix} -d_A^{n+1} & 0 \\ f^{n+1} & d_B^n \end{pmatrix}
  = \begin{pmatrix} d_A^{n+2} d_A^{n+1} & 0 \\
    -f^{n+2} d_A^{n+1} + d_B^{n+1} f^{n+1} & d_B^{n+1} d_B^n \end{pmatrix}
  = 0,
\]
utilisant $d^2 = 0$ et $f \circ d_A = d_B \circ f$.

Pour la suite exacte courte de complexes, on a
$0 \to B^\bullet \xrightarrow{\iota} \operatorname{Cone}(f) \xrightarrow{\pi} A[1]^\bullet \to 0$
qui est exacte degré par degré. La suite exacte longue en cohomologie
en découle (voir théorème~\ref{thm:long-exact-sequence}).
Le morphisme connectant $H^n(A[1]^\bullet) = H^{n+1}(A^\bullet) \to H^{n+1}(B^\bullet)$
s'identifie à $H^{n+1}(f)$.
\end{proof}


%% ============================================================
\section{Structure triangulée}\label{sec:triangulated}
%% ============================================================

\begin{definition}[Catégorie triangulée]\label{def:triangulated}
\index{catégorie!triangulée|textbf}
\index{axiomes TR}
Une \textbf{catégorie triangulée} est un triplet
$(\mathcal{T}, [1], \Delta)$ où :
\begin{itemize}
\item $\mathcal{T}$ est une catégorie additive ;
\item $[1] \colon \mathcal{T} \to \mathcal{T}$ est un automorphisme
  (le \textbf{foncteur de translation}\index{foncteur!de translation}) ;
\item $\Delta$ est une classe de suites
  $A \xrightarrow{u} B \xrightarrow{v} C \xrightarrow{w} A[1]$
  appelées \textbf{triangles distingués}\index{triangle!distingué|textbf},
\end{itemize}
satisfaisant les axiomes suivants.
\end{definition}

\begin{definition}[Axiomes TR1--TR4]\label{def:TR-axioms}
\index{TR1}\index{TR2}\index{TR3}\index{TR4}
\begin{enumerate}[label=\textbf{(TR\arabic*)}]
\item\label{TR1} \textbf{(Identité et complétion.)}
  \begin{enumerate}[label=(\alph*)]
  \item Pour tout objet $A$, le triangle
    $A \xrightarrow{\id} A \to 0 \to A[1]$ est distingué.
  \item Tout morphisme $u \colon A \to B$ peut être complété en un
    triangle distingué $A \xrightarrow{u} B \to C \to A[1]$.
  \item Tout triangle isomorphe à un triangle distingué est distingué.
  \end{enumerate}

\item\label{TR2} \textbf{(Rotation.)}
  Un triangle $A \xrightarrow{u} B \xrightarrow{v} C \xrightarrow{w} A[1]$
  est distingué si et seulement si
  $B \xrightarrow{v} C \xrightarrow{w} A[1] \xrightarrow{-u[1]} B[1]$
  est distingué.

\item\label{TR3} \textbf{(Morphisme de triangles.)}
  Étant donnés deux triangles distingués et des morphismes $f, g$
  formant un carré commutatif comme ci-dessous, il existe un morphisme $h$
  complétant le diagramme en un morphisme de triangles :
\[
\begin{tikzcd}
A \arrow[r, "u"] \arrow[d, "f"] & B \arrow[r, "v"] \arrow[d, "g"]
  & C \arrow[r, "w"] \arrow[d, dashed, "h"] & A{[1]} \arrow[d, "f{[1]}"] \\
A' \arrow[r, "u'"] & B' \arrow[r, "v'"] & C' \arrow[r, "w'"] & A'{[1]}
\end{tikzcd}
\]

\item\label{TR4} \textbf{(Axiome de l'octaèdre.)}
\index{axiome!de l'octaèdre}
  Étant donnés des morphismes $u \colon A \to B$ et $v \colon B \to C$,
  et des triangles distingués
  \begin{align*}
    & A \xrightarrow{u} B \to D \to A[1], \\
    & B \xrightarrow{v} C \to E \to B[1], \\
    & A \xrightarrow{v \circ u} C \to F \to A[1],
  \end{align*}
  il existe un triangle distingué $D \to F \to E \to D[1]$ tel que
  le diagramme complet commute.
\end{enumerate}
\end{definition}

\begin{theoreme}\label{thm:K-triangulated}
\index{K(A)@$K(\mathcal{A})$!structure triangulée}
\index{catégorie!homotopique!triangulée}
La catégorie homotopique $K(\mathcal{A})$, munie du foncteur de décalage
$[1]$ et des triangles de la forme
\[
  A^\bullet \xrightarrow{f} B^\bullet \xrightarrow{\iota}
  \operatorname{Cone}(f) \xrightarrow{\pi} A[1]^\bullet
\]
(et de leurs isomorphes), est une catégorie triangulée.
\end{theoreme}

\begin{proof}
Nous vérifions les axiomes.

\ref{TR1} : (a) Le cône de $\id \colon A^\bullet \to A^\bullet$ est
homotopiquement trivial : $\operatorname{Cone}(\id)^n = A^{n+1} \oplus A^n$
avec la différentielle qui en fait un complexe acyclique (exercice).
(b) Par définition, tout $f$ donne un triangle via le cône.
(c) Par définition de la classe $\Delta$.

\ref{TR2} : La rotation des triangles se vérifie par un calcul explicite
d'homotopie entre $\operatorname{Cone}(\iota)$ et $A[1]^\bullet$.

\ref{TR3} : On construit $h$ composante par composante en utilisant les
propriétés universelles dans la catégorie homotopique.

\ref{TR4} : Vérification par construction explicite, voir par exemple
Weibel \cite[1.5]{weibel1994}.
\end{proof}

\begin{corollaire}\label{cor:D-triangulated}
\index{D(A)@$D(\mathcal{A})$!structure triangulée}
La catégorie dérivée $D(\mathcal{A})$ hérite d'une structure triangulée.
De même pour $D^+(\mathcal{A})$, $D^-(\mathcal{A})$ et $D^b(\mathcal{A})$.
\end{corollaire}


%% ============================================================
\section{Foncteurs dérivés}\label{sec:derived-functors}
%% ============================================================

\begin{definition}[Foncteur dérivé à droite]\label{def:right-derived}
\index{foncteur!dérivé à droite|textbf}
\index{RF@$RF$}
Soit $F \colon \mathcal{A} \to \mathcal{B}$ un foncteur exact à gauche
entre catégories abéliennes, avec $\mathcal{A}$ ayant assez d'injectifs.
Le \textbf{foncteur dérivé à droite} $RF \colon D^+(\mathcal{A}) \to D^+(\mathcal{B})$
est défini comme suit :
\begin{enumerate}[label=\arabic*.]
\item Pour $A^\bullet \in D^+(\mathcal{A})$, on choisit une résolution
  injective $A^\bullet \xrightarrow{\sim} I^\bullet$ ;
\item On pose $RF(A^\bullet) = F(I^\bullet)$
  (application de $F$ degré par degré).
\end{enumerate}
Ceci est bien défini (indépendant du choix de résolution, à isomorphisme
canonique près dans $D^+(\mathcal{B})$) par le lemme de comparaison
(théorème~\ref{thm:comparison-lemma}).
\end{definition}

\begin{definition}[Foncteur dérivé à gauche]\label{def:left-derived}
\index{foncteur!dérivé à gauche|textbf}
\index{LF@$LF$}
Dualement, si $F \colon \mathcal{A} \to \mathcal{B}$ est exact à droite
et $\mathcal{A}$ a assez de projectifs, le \textbf{foncteur dérivé à gauche}
$LF \colon D^-(\mathcal{A}) \to D^-(\mathcal{B})$ est défini par
$LF(A^\bullet) = F(P^\bullet)$ où $P^\bullet \xrightarrow{\sim} A^\bullet$
est une résolution projective.
\end{definition}

\begin{definition}[Foncteurs dérivés classiques]\label{def:classical-derived}
\index{foncteur!dérivé!classique}
Pour un objet $A$ vu comme complexe concentré en degré $0$, on pose :
\[
  R^n F(A) = H^n(RF(A)) \quad \text{pour tout } n \geq 0.
\]
On a $R^0 F = F$ (car $F$ est exact à gauche). De même,
$L_n F(A) = H^{-n}(LF(A))$ avec $L_0 F = F$.
\end{definition}

\begin{definition}[$\operatorname{Ext}$ et $\operatorname{Tor}$]%
\label{def:ext-tor}
\index{Ext@$\operatorname{Ext}$|textbf}
\index{Tor@$\operatorname{Tor}$|textbf}
Soit $R$ un anneau et $A, B$ des $R$-modules.
\begin{enumerate}[label=(\roman*)]
\item Les \textbf{groupes d'extension}\index{groupe d'extension} sont
\[
  \operatorname{Ext}_R^n(A, B) = R^n \Hom_R(A, -)(B)
  = H^n(\Hom_R(A, I^\bullet))
\]
où $B \to I^\bullet$ est une résolution injective.
Alternativement,
$\operatorname{Ext}_R^n(A, B) = H^n(\Hom_R(P_\bullet, B))$
où $P_\bullet \to A$ est une résolution projective.

\item Les \textbf{groupes de torsion}\index{groupe de torsion!Tor} sont
\[
  \operatorname{Tor}_n^R(A, B) = L_n(A \otimes_R -)(B)
  = H_n(A \otimes_R P_\bullet)
\]
où $P_\bullet \to B$ est une résolution projective.
\end{enumerate}
\end{definition}

\begin{remarque}\label{rem:ext-tor-balanced}
\index{Ext@$\operatorname{Ext}$!balancé}
\index{Tor@$\operatorname{Tor}$!balancé}
Les foncteurs $\operatorname{Ext}$ et $\operatorname{Tor}$ sont
\emph{balancés} : on peut calculer $\operatorname{Ext}_R^n(A, B)$
soit via une résolution projective de $A$, soit via une résolution
injective de $B$, et on obtient le même résultat (à isomorphisme
canonique près). De même, $\operatorname{Tor}_n^R(A, B)$ peut être
calculé en résolvant soit $A$, soit $B$.
\end{remarque}

\begin{proposition}\label{prop:ext0-tor0}
\index{Ext@$\operatorname{Ext}$!en degré zéro}
\index{Tor@$\operatorname{Tor}$!en degré zéro}
On a les identifications :
\[
  \operatorname{Ext}_R^0(A, B) \cong \Hom_R(A, B), \qquad
  \operatorname{Tor}_0^R(A, B) \cong A \otimes_R B.
\]
\end{proposition}

\begin{proof}
Pour $\operatorname{Ext}^0$ : $H^0(\Hom_R(P_\bullet, B))
= \ker(\Hom(P_0, B) \to \Hom(P_1, B))$, qui s'identifie
à $\Hom(A, B)$ par exactitude de
$0 \to \Hom(A, B) \to \Hom(P_0, B) \to \Hom(P_1, B)$
(le foncteur $\Hom(-, B)$ est exact à gauche).
L'argument pour $\operatorname{Tor}_0$ est analogue.
\end{proof}


%% ============================================================
\section{Suite exacte longue en cohomologie}\label{sec:long-exact}
%% ============================================================

\begin{theoreme}[Suite exacte longue]\label{thm:long-exact-sequence}
\index{suite exacte!longue|textbf}
\index{morphisme connectant!suite exacte longue}
Soit $0 \to A^\bullet \xrightarrow{f} B^\bullet \xrightarrow{g} C^\bullet \to 0$
une suite exacte courte de complexes dans une catégorie abélienne. Alors
il existe des morphismes connectants
$\delta^n \colon H^n(C^\bullet) \to H^{n+1}(A^\bullet)$ et une suite
exacte longue
\[
\begin{tikzcd}[column sep=small]
\cdots \arrow[r]
& H^n(A^\bullet) \arrow[r, "H^n(f)"]
& H^n(B^\bullet) \arrow[r, "H^n(g)"]
& H^n(C^\bullet) \arrow[r, "\delta^n"]
& H^{n+1}(A^\bullet) \arrow[r]
& \cdots
\end{tikzcd}
\]
De plus, $\delta^n$ est naturel en les suites exactes courtes.
\end{theoreme}

\begin{proof}
\textbf{Construction de $\delta^n$.}
Soit $[c] \in H^n(C^\bullet)$, représenté par $c \in C^n$ avec
$d_C^n(c) = 0$. Comme $g^n$ est surjectif, il existe $b \in B^n$
avec $g^n(b) = c$. Considérons $d_B^n(b) \in B^{n+1}$.
On a
\[
  g^{n+1}(d_B^n(b)) = d_C^n(g^n(b)) = d_C^n(c) = 0,
\]
donc $d_B^n(b) \in \ker g^{n+1} = \im f^{n+1}$, et il existe un unique
$a \in A^{n+1}$ avec $f^{n+1}(a) = d_B^n(b)$.

\emph{$a$ est un cocycle.} On a $f^{n+2}(d_A^{n+1}(a)) = d_B^{n+1}(f^{n+1}(a))
= d_B^{n+1}(d_B^n(b)) = 0$, et comme $f^{n+2}$ est injectif,
$d_A^{n+1}(a) = 0$.

\emph{Bonne définition.} Si $b'$ est un autre relèvement de $c$, alors
$b - b' \in \ker g^n = \im f^n$, donc $b - b' = f^n(a_0)$ et
$d_B^n(b) - d_B^n(b') = d_B^n(f^n(a_0)) = f^{n+1}(d_A^n(a_0))$,
d'où $a - a' = d_A^n(a_0) \in \im d_A^n$.

On pose $\delta^n([c]) = [a] \in H^{n+1}(A^\bullet)$.

\medskip
\textbf{Exactitude en $H^n(C^\bullet)$.}
$\delta^n \circ H^n(g) = 0$ : si $[c] = H^n(g)([b_0])$ avec $d_B^n(b_0) = 0$,
on peut prendre $b = b_0$ et alors $d_B^n(b) = 0$, donc $a = 0$.

Réciproquement, si $\delta^n([c]) = 0$, alors $a = d_A^n(a_0)$ et
$b' = b - f^n(a_0)$ vérifie $d_B^n(b') = f^{n+1}(a - d_A^n(a_0)) = 0$ et
$g^n(b') = c$, donc $[c] = H^n(g)([b'])$.

\medskip
\textbf{Exactitude en $H^{n+1}(A^\bullet)$.}
$H^{n+1}(f) \circ \delta^n = 0$ : $H^{n+1}(f)([a]) = [f^{n+1}(a)]
= [d_B^n(b)] = 0$ dans $H^{n+1}(B^\bullet)$.

Réciproquement, si $[a] \in \ker H^{n+1}(f)$, alors $f^{n+1}(a) = d_B^n(b)$
pour un $b \in B^n$, et $c = g^n(b)$ est un cocycle (car
$d_C^n(c) = g^{n+1}(d_B^n(b)) = g^{n+1}(f^{n+1}(a)) = 0$), avec
$\delta^n([c]) = [a]$.

L'exactitude en $H^n(A^\bullet)$ et $H^n(B^\bullet)$ se vérifie de manière
analogue.
\end{proof}

\begin{corollaire}\label{cor:long-exact-derived}
\index{foncteur!dérivé!suite exacte longue}
\index{Ext@$\operatorname{Ext}$!suite exacte longue}
Soit $F \colon \mathcal{A} \to \mathcal{B}$ un foncteur exact à gauche
et $0 \to A \to B \to C \to 0$ une suite exacte courte dans $\mathcal{A}$.
Alors on a une suite exacte longue
\[
\begin{tikzcd}[column sep=small]
0 \arrow[r]
& FA \arrow[r]
& FB \arrow[r]
& FC \arrow[r, "\delta^0"]
& R^1 F(A) \arrow[r]
& R^1 F(B) \arrow[r]
& R^1 F(C) \arrow[r]
& \cdots
\end{tikzcd}
\]
En particulier, pour $F = \Hom_R(M, -)$ :
\[
\begin{tikzcd}[column sep=small]
0 \arrow[r]
& \Hom(M, A) \arrow[r]
& \Hom(M, B) \arrow[r]
& \Hom(M, C) \arrow[r]
& \operatorname{Ext}^1(M, A) \arrow[r]
& \cdots
\end{tikzcd}
\]
\end{corollaire}

\begin{proof}
On applique $RF$ à la suite exacte courte (vue comme triangle distingué
dans $D^+(\mathcal{A})$), ce qui donne un triangle distingué dans
$D^+(\mathcal{B})$. La suite exacte longue en cohomologie associée
à ce triangle donne le résultat.
\end{proof}


%% ============================================================
\section{Exercices}\label{sec:exos-derived}
%% ============================================================

\begin{exercice}\label{exo:cone-id-acyclic}
\index{cône!de mapping!de l'identité}
Montrer que $\operatorname{Cone}(\id_{A^\bullet})$ est un complexe acyclique
(c'est-à-dire $H^n(\operatorname{Cone}(\id)) = 0$ pour tout $n$),
et construire explicitement une homotopie contractante.
\end{exercice}

\begin{exercice}\label{exo:quasi-iso-cone}
\index{quasi-isomorphisme!et cône}
Montrer qu'un morphisme de complexes $f \colon A^\bullet \to B^\bullet$
est un quasi-isomorphisme si et seulement si
$\operatorname{Cone}(f)$ est acyclique.
\end{exercice}

\begin{exercice}\label{exo:ext-cyclic}
\index{Ext@$\operatorname{Ext}$!groupes cycliques}
Calculer $\operatorname{Ext}_\Z^n(\Z/m\Z, \Z/k\Z)$ pour tous $m, k \geq 1$
et tout $n \geq 0$.
\textit{Indication : utiliser la résolution projective
$0 \to \Z \xrightarrow{\times m} \Z \to \Z/m\Z \to 0$.}
\end{exercice}

\begin{exercice}\label{exo:tor-cyclic}
\index{Tor@$\operatorname{Tor}$!groupes cycliques}
Calculer $\operatorname{Tor}_n^\Z(\Z/m\Z, \Z/k\Z)$ pour tous $m, k \geq 1$
et tout $n \geq 0$.
\end{exercice}

\begin{exercice}\label{exo:ext-ses-interpretation}
\index{Ext@$\operatorname{Ext}$!interprétation par extensions}
\index{extension!de modules}
Soit $R$ un anneau et $A, C$ des $R$-modules. Montrer que
$\operatorname{Ext}_R^1(C, A)$ classifie les suites exactes courtes
$0 \to A \to B \to C \to 0$ à équivalence de Yoneda près.

\textit{Indication : à une telle suite, associer l'image de $\id_C$ par
le morphisme connectant $\Hom(C, C) \to \operatorname{Ext}^1(C, A)$.}
\end{exercice}

\begin{exercice}\label{exo:triangulated-abelian}
\index{catégorie!triangulée!pas abélienne}
Montrer qu'une catégorie triangulée non triviale n'est pas abélienne.

\textit{Indication : dans une catégorie triangulée, $\Hom(A, B)$ est
un groupe abélien et les suites exactes sont remplacées par des triangles.
Montrer que les biproduits ne coïncident pas avec les produits en général
dans une catégorie triangulée.}
\end{exercice}

\begin{exercice}\label{exo:derived-category-Z}
\index{D(Ab)@$D(\Ab)$}
\index{catégorie!dérivée!de $\Ab$}
Soit $A$ un groupe abélien. Montrer que dans $D(\Ab)$, $A$ (vu comme
complexe concentré en degré $0$) est isomorphe à toute résolution
projective ou injective de $A$.
\end{exercice}

\begin{exercice}\label{exo:horseshoe-derived}
\index{lemme!du fer à cheval!et foncteurs dérivés}
En utilisant le lemme du fer à cheval (exercice~\ref{exo:horseshoe}),
donner une preuve directe de l'existence de la suite exacte longue
des foncteurs dérivés (corollaire~\ref{cor:long-exact-derived})
sans passer par les catégories dérivées.
\end{exercice}

\begin{exercice}\label{exo:dimension-homologique}
\index{dimension homologique}
\index{dimension projective}
La \textbf{dimension projective}\index{dimension projective|textbf} d'un
$R$-module $M$, notée $\operatorname{pd}_R(M)$, est la longueur minimale
d'une résolution projective de $M$ (ou $+\infty$ si aucune n'est finie).
\begin{enumerate}[label=(\alph*)]
\item Montrer que $\operatorname{pd}_R(M) \leq n$ si et seulement si
  $\operatorname{Ext}_R^{n+1}(M, N) = 0$ pour tout $N$.
\item La \textbf{dimension globale}\index{dimension!globale} de $R$ est
  $\operatorname{gldim}(R) = \sup_M \operatorname{pd}_R(M)$.
  Montrer que $\operatorname{gldim}(\Z) = 1$.
\item Montrer que $\operatorname{gldim}(k) = 0$ pour tout corps $k$.
\end{enumerate}
\end{exercice}

\begin{exercice}\label{exo:rotation-triangle}
\index{triangle!distingué!rotation}
Vérifier en détail l'axiome \ref{TR2} pour la catégorie homotopique
$K(\mathcal{A})$ : si $A^\bullet \xrightarrow{f} B^\bullet \xrightarrow{\iota}
\operatorname{Cone}(f) \xrightarrow{\pi} A[1]^\bullet$ est un triangle
distingué, construire un isomorphisme explicite entre
$\operatorname{Cone}(\iota)$ et $A[1]^\bullet$ dans $K(\mathcal{A})$.
\end{exercice}

\begin{exercice}\label{exo:derived-adjunction}
\index{foncteur!dérivé!et adjonction}
Soit $(F, G)$ un couple de foncteurs adjoints entre catégories abéliennes
avec assez d'injectifs et de projectifs. Si $F$ est exact à droite et $G$
est exact à gauche, montrer que $(LF, RG)$ est un couple de foncteurs
adjoints au niveau des catégories dérivées :
\[
  \Hom_{D^-(\mathcal{B})}(LF(A^\bullet), B^\bullet)
  \cong \Hom_{D^-(\mathcal{A})}(A^\bullet, RG(B^\bullet)).
\]
\end{exercice}

\index{complexe de chaînes|)}
\index{catégorie!dérivée|)}
%%%%%%%%%%%%%%%%%%%%%%%%%%%%%%%%%%%%%%%%%%%%%%%%%%%%%%%%%%%%%
%% THÉORIE DES CATÉGORIES — Chapitres 9–11 + Fin
%% Catégories monoïdales / Topos / ∞-Catégories / Appendice
%%%%%%%%%%%%%%%%%%%%%%%%%%%%%%%%%%%%%%%%%%%%%%%%%%%%%%%%%%%%%

%%%%%%%%%%%%%%%%%%%%%%%%%%%%%%%%%%%%%%%%%%%%%%%%%%%%%%%%%%%%
%%       CHAPITRE 9 : CATÉGORIES MONOÏDALES ET TRESSÉES   %%
%%%%%%%%%%%%%%%%%%%%%%%%%%%%%%%%%%%%%%%%%%%%%%%%%%%%%%%%%%%%

\chapter{Catégories monoïdales et tressées}\label{ch:monoidales}
\minitoc

\vspace{1em}

Les catégories monoïdales formalisent la notion de \og produit tensoriel\fg{}
dans un cadre catégorique général.  Elles interviennent en algèbre
(modules sur un anneau), en topologie (espaces de lacets),
en informatique théorique (types linéaires, sémantique catégorique)
et en physique mathématique (théories topologiques des champs).
Ce chapitre introduit les structures monoïdales, symétriques et tressées,
puis développe les notions de catégorie enrichie et de catégorie
cartésienne fermée, avec l'application fondamentale à la correspondance
de Curry--Howard--Lambek.

\index{catégorie monoïdale|(}

%% ============================================================
\section{Catégories monoïdales}\label{sec:cat-monoidale}
%% ============================================================

\begin{definition}[Catégorie monoïdale]\label{def:cat-monoidale}
\index{catégorie monoïdale!définition}
\index{produit tensoriel}
\index{objet unité}
\index{associateur}
\index{uniteur}
Une \textbf{catégorie monoïdale} est un sextuplet
$(\mathcal{C}, \otimes, I, \alpha, \lambda, \rho)$ où :
\begin{enumerate}[label=(\roman*)]
\item $\mathcal{C}$ est une catégorie ;
\item $\otimes \colon \mathcal{C} \times \mathcal{C} \to \mathcal{C}$
      est un bifoncteur, appelé \textbf{produit tensoriel} ;
\item $I \in \Ob(\mathcal{C})$ est l'\textbf{objet unité} ;
\item $\alpha_{A,B,C} \colon (A \otimes B) \otimes C
      \xrightarrow{\;\sim\;} A \otimes (B \otimes C)$
      est un isomorphisme naturel en $A, B, C$,
      appelé l'\textbf{associateur} ;
\item $\lambda_A \colon I \otimes A \xrightarrow{\;\sim\;} A$
      est un isomorphisme naturel en $A$,
      appelé l'\textbf{uniteur à gauche} ;
\item $\rho_A \colon A \otimes I \xrightarrow{\;\sim\;} A$
      est un isomorphisme naturel en $A$,
      appelé l'\textbf{uniteur à droite} ;
\end{enumerate}
soumis aux conditions de cohérence suivantes.

\medskip
\noindent\textbf{Diagramme du pentagone.}\index{pentagone (axiome du)}
Pour tous $A, B, C, D \in \Ob(\mathcal{C})$ :
\[
\begin{tikzcd}[column sep=small]
& (A \otimes B) \otimes (C \otimes D)
  \arrow[dr, "\alpha_{A,B,C \otimes D}"] & \\
((A \otimes B) \otimes C) \otimes D
  \arrow[ur, "\alpha_{A \otimes B,C,D}"]
  \arrow[d, "\alpha_{A,B,C} \otimes \id_D"'] & &
A \otimes (B \otimes (C \otimes D)) \\
(A \otimes (B \otimes C)) \otimes D
  \arrow[rr, "\alpha_{A,B \otimes C,D}"'] & &
A \otimes ((B \otimes C) \otimes D)
  \arrow[u, "\id_A \otimes \alpha_{B,C,D}"']
\end{tikzcd}
\]

\noindent\textbf{Diagramme du triangle.}\index{triangle (axiome du)}
Pour tous $A, B \in \Ob(\mathcal{C})$ :
\[
\begin{tikzcd}
(A \otimes I) \otimes B
  \arrow[rr, "\alpha_{A,I,B}"]
  \arrow[dr, "\rho_A \otimes \id_B"'] & &
A \otimes (I \otimes B)
  \arrow[dl, "\id_A \otimes \lambda_B"] \\
& A \otimes B &
\end{tikzcd}
\]
\end{definition}

\begin{definition}[Catégorie monoïdale stricte]\label{def:monoidale-stricte}
\index{catégorie monoïdale!stricte}
Une catégorie monoïdale est dite \textbf{stricte} si les isomorphismes
$\alpha$, $\lambda$ et $\rho$ sont tous des identités, c'est-à-dire
si $(A \otimes B) \otimes C = A \otimes (B \otimes C)$,
$I \otimes A = A$ et $A \otimes I = A$ pour tous objets $A, B, C$.
\end{definition}

\begin{theoreme}[Cohérence de Mac Lane]\label{thm:coherence-maclane}
\index{théorème de cohérence!Mac Lane}
\index{Mac Lane!théorème de cohérence}
Soit $(\mathcal{C}, \otimes, I, \alpha, \lambda, \rho)$ une catégorie
monoïdale.
\begin{enumerate}[label=(\alph*)]
\item Tout diagramme dont les flèches sont des composées d'instances de
      $\alpha$, $\lambda$, $\rho$ et de leurs inverses commute.
\item Toute catégorie monoïdale est monoïdalement équivalente à une
      catégorie monoïdale stricte.
\end{enumerate}
\end{theoreme}

\begin{proof}
La partie (a) est le théorème original de Mac Lane (1963).
On montre d'abord que tout diagramme de cohérence se réduit à un diagramme
dans la catégorie libre monoïdale sur un générateur, puis on vérifie
la commutativité dans ce cas par induction sur la complexité des
expressions.  La partie (b) se déduit en construisant une catégorie
monoïdale stricte $\mathcal{C}_s$ avec une équivalence monoïdale
$\mathcal{C} \simeq \mathcal{C}_s$ (cf.\ Joyal--Street).
\end{proof}

\begin{remarque}\label{rem:coherence-pratique}
\index{théorème de cohérence!conséquence pratique}
En pratique, le théorème de cohérence nous autorise à omettre les
parenthèses et à identifier $I \otimes A$, $A \otimes I$ et $A$
sans ambiguïté.  On écrit simplement $A \otimes B \otimes C$
au lieu de $(A \otimes B) \otimes C$ ou $A \otimes (B \otimes C)$.
\end{remarque}

%% ============================================================
\section{Exemples fondamentaux}\label{sec:monoidale-exemples}
%% ============================================================

\begin{exemple}[Catégories monoïdales classiques]\label{ex:monoidales-classiques}
\index{catégorie monoïdale!exemples}
\begin{enumerate}[label=(\roman*)]
\item \textbf{$(\Set, \times, \{*\})$} :\index{Set@$\Set$!monoïdale}
      la catégorie des ensembles avec le produit cartésien.
      L'objet unité est un singleton.

\item \textbf{$(\Ab, \otimes_\Z, \Z)$} :\index{Ab@$\Ab$!monoïdale}
      la catégorie des groupes abéliens avec le produit tensoriel
      sur $\Z$.

\item \textbf{$(R\text{-}\Mod, \otimes_R, R)$} :\index{Mod@$\Mod$!monoïdale}
      pour un anneau commutatif $R$, la catégorie des $R$-modules
      avec le produit tensoriel.

\item \textbf{$(\Vect_\K, \otimes_\K, \K)$} :\index{Vect@$\Vect$!monoïdale}
      la catégorie des $\K$-espaces vectoriels.

\item \textbf{$(\Cat, \times, \mathbf{1})$} :\index{Cat@$\Cat$!monoïdale}
      la catégorie des petites catégories avec le produit.

\item \textbf{$(\Set, \sqcup, \varnothing)$} :
      les ensembles avec l'union disjointe.  L'objet unité
      est l'ensemble vide.

\item \textbf{$([\mathcal{C}, \mathcal{C}], \circ, \id_{\mathcal{C}})$} :
      \index{catégorie d'endofoncteurs}
      la catégorie des endofoncteurs de $\mathcal{C}$ avec la
      composition.  Cette catégorie monoïdale est stricte.
\end{enumerate}
\end{exemple}

\begin{exemple}[Structures monoïdales sur les espaces vectoriels gradués]
\label{ex:ev-gradues-monoidal}
\index{espace vectoriel gradué!structure monoïdale}
Soit $\Vect_\K^{\Z}$ la catégorie des espaces vectoriels $\Z$-gradués.
Le produit tensoriel gradué est défini par
\[
  (V \otimes W)_n = \bigoplus_{p+q=n} V_p \otimes_\K W_q.
\]
L'objet unité est $\K$ concentré en degré $0$.
Cette structure monoïdale joue un rôle central en algèbre homologique
et en topologie algébrique.
\end{exemple}

%% ============================================================
\section{Catégories monoïdales symétriques et tressées}
\label{sec:symetrique-tressee}
%% ============================================================

\begin{definition}[Catégorie monoïdale tressée]\label{def:monoidale-tressee}
\index{catégorie monoïdale!tressée}
\index{tressage}
Une \textbf{catégorie monoïdale tressée} est une catégorie monoïdale
$(\mathcal{C}, \otimes, I)$ munie d'un isomorphisme naturel
\[
  \sigma_{A,B} \colon A \otimes B \xrightarrow{\;\sim\;} B \otimes A,
\]
appelé \textbf{tressage}, satisfaisant les deux diagrammes hexagonaux
suivants.

\medskip
\noindent\textbf{Premier hexagone} :\index{hexagone (axiome de l')}
\[
\begin{tikzcd}[column sep=1.8em]
(A \otimes B) \otimes C
  \arrow[r, "\sigma_{A,B} \otimes \id"]
  \arrow[d, "\alpha"'] &
(B \otimes A) \otimes C
  \arrow[r, "\alpha"] &
B \otimes (A \otimes C)
  \arrow[d, "\id \otimes \sigma_{A,C}"] \\
A \otimes (B \otimes C)
  \arrow[r, "\sigma_{A,B \otimes C}"'] &
(B \otimes C) \otimes A
  \arrow[r, "\alpha"'] &
B \otimes (C \otimes A)
\end{tikzcd}
\]

\noindent\textbf{Second hexagone} : le diagramme analogue avec
$\sigma^{-1}$ et les associateurs inversés.
\end{definition}

\begin{definition}[Catégorie monoïdale symétrique]\label{def:monoidale-symetrique}
\index{catégorie monoïdale!symétrique}
\index{symétrie (tressage)}
Une catégorie monoïdale tressée est dite \textbf{symétrique} si
$\sigma_{B,A} \circ \sigma_{A,B} = \id_{A \otimes B}$
pour tous $A, B$.
\end{definition}

\begin{exemple}\label{ex:symetrique-vs-tressee}
\begin{enumerate}[label=(\roman*)]
\item $(\Set, \times, \{*\})$, $(\Ab, \otimes_\Z, \Z)$,
      $(\Vect_\K, \otimes_\K, \K)$ sont symétriques.
\item La catégorie des représentations d'un groupe quantique
      $U_q(\mathfrak{g})$ pour $q$ générique est tressée mais
      \emph{pas} symétrique en général.
      \index{groupe quantique}
\end{enumerate}
\end{exemple}

%% ============================================================
\section{Objets monoïdes}\label{sec:objets-monoides}
%% ============================================================

\begin{definition}[Monoïde dans une catégorie monoïdale]
\label{def:objet-monoide}
\index{monoïde!dans une catégorie monoïdale}
\index{objet monoïde}
Soit $(\mathcal{C}, \otimes, I)$ une catégorie monoïdale.
Un \textbf{objet monoïde} (ou \textbf{monoïde dans $\mathcal{C}$})
est un triplet $(M, \mu, \eta)$ où :
\begin{enumerate}[label=(\roman*)]
\item $M \in \Ob(\mathcal{C})$ ;
\item $\mu \colon M \otimes M \to M$ est la \textbf{multiplication} ;
\item $\eta \colon I \to M$ est l'\textbf{unité} ;
\end{enumerate}
satisfaisant les axiomes d'associativité et d'unité :
\[
\begin{tikzcd}
(M \otimes M) \otimes M
  \arrow[r, "\alpha"]
  \arrow[d, "\mu \otimes \id"'] &
M \otimes (M \otimes M)
  \arrow[d, "\id \otimes \mu"] \\
M \otimes M \arrow[r, "\mu"'] &
M
\end{tikzcd}
\qquad
\begin{tikzcd}
I \otimes M
  \arrow[r, "\eta \otimes \id"]
  \arrow[dr, "\lambda"'] &
M \otimes M
  \arrow[d, "\mu"] &
M \otimes I
  \arrow[l, "\id \otimes \eta"']
  \arrow[dl, "\rho"] \\
& M &
\end{tikzcd}
\]
\end{definition}

\begin{exemple}[Monoïdes dans diverses catégories]\label{ex:objets-monoides}
\index{objet monoïde!exemples}
\begin{enumerate}[label=(\roman*)]
\item Un monoïde dans $(\Set, \times, \{*\})$ est un monoïde
      au sens usuel.
\item Un monoïde dans $(\Ab, \otimes_\Z, \Z)$ est un anneau.
      \index{anneau!comme monoïde dans Ab}
\item Un monoïde dans $(R\text{-}\Mod, \otimes_R, R)$ est une
      $R$-algèbre.\index{algèbre!comme monoïde}
\item Un monoïde dans
      $([\mathcal{C}, \mathcal{C}], \circ, \id_\mathcal{C})$
      est une monade.\index{monade!comme monoïde}
\end{enumerate}
\end{exemple}

\begin{definition}[Comonoïde]\label{def:comonoide}
\index{comonoïde}
Dualement, un \textbf{comonoïde} dans $(\mathcal{C}, \otimes, I)$
est un triplet $(C, \Delta, \varepsilon)$ avec
$\Delta \colon C \to C \otimes C$ (comultiplication) et
$\varepsilon \colon C \to I$ (coünité), satisfaisant les axiomes duaux.
Un comonoïde dans $(\Vect_\K, \otimes_\K, \K)$ est une coalgèbre.
\index{coalgèbre}
\end{definition}

%% ============================================================
\section{Catégories enrichies}\label{sec:categories-enrichies}
%% ============================================================

\begin{definition}[Catégorie enrichie]\label{def:cat-enrichie}
\index{catégorie enrichie}
\index{catégorie-$\mathcal{V}$}
Soit $(\mathcal{V}, \otimes, I)$ une catégorie monoïdale.
Une \textbf{catégorie enrichie sur $\mathcal{V}$} (ou
\textbf{$\mathcal{V}$-catégorie}) $\mathcal{C}$ consiste en :
\begin{enumerate}[label=(\roman*)]
\item une collection d'objets $\Ob(\mathcal{C})$ ;
\item pour chaque paire $(A, B)$, un objet
      $\mathcal{C}(A,B) \in \Ob(\mathcal{V})$
      (l'\textbf{objet de morphismes}) ;
\item pour chaque triplet $(A, B, C)$, un morphisme de composition
      $\circ_{A,B,C} \colon \mathcal{C}(B,C) \otimes \mathcal{C}(A,B)
      \to \mathcal{C}(A,C)$ dans $\mathcal{V}$ ;
\item pour chaque objet $A$, un morphisme identité
      $j_A \colon I \to \mathcal{C}(A,A)$ dans $\mathcal{V}$ ;
\end{enumerate}
satisfaisant les axiomes d'associativité et d'unité
(exprimés comme commutativité de diagrammes dans $\mathcal{V}$).
\end{definition}

\begin{exemple}[Enrichissements classiques]\label{ex:enrichissements}
\index{catégorie enrichie!exemples}
\begin{enumerate}[label=(\roman*)]
\item $\Set$-catégories $=$ catégories ordinaires
      (localement petites).
\item $\Ab$-catégories $=$ catégories préadditives.
      \index{catégorie préadditive}
\item $\Vect_\K$-catégories $=$ catégories $\K$-linéaires.
\item $\Cat$-catégories $=$ $2$-catégories (strictes).
      \index{2-catégorie}
\item $([0,\infty], \geq, +, 0)$-catégories $=$ espaces métriques
      généralisés (au sens de Lawvere).
      \index{espace métrique!de Lawvere}
\end{enumerate}
\end{exemple}

%% ============================================================
\section{Catégories fermées}\label{sec:categories-fermees}
%% ============================================================

\begin{definition}[Catégorie monoïdale fermée]\label{def:monoidale-fermee}
\index{catégorie monoïdale!fermée}
\index{Hom interne}
\index{exponentiation}
Une catégorie monoïdale $(\mathcal{C}, \otimes, I)$ est dite
\textbf{fermée (à droite)} si pour chaque objet $B$,
le foncteur $- \otimes B \colon \mathcal{C} \to \mathcal{C}$
admet un adjoint à droite, noté $[B, -]$ ou $\underline{\Hom}(B, -)$ :
\[
  \Hom_\mathcal{C}(A \otimes B, C)
  \;\cong\;
  \Hom_\mathcal{C}(A, [B, C])
\]
naturellement en $A$ et $C$.
L'objet $[B, C]$ est appelé le \textbf{Hom interne}.
\end{definition}

\begin{definition}[Catégorie cartésienne fermée]\label{def:cartesienne-fermee}
\index{catégorie cartésienne fermée}
\index{CCC|see{catégorie cartésienne fermée}}
Une catégorie est \textbf{cartésienne fermée} (CCC) si :
\begin{enumerate}[label=(\roman*)]
\item elle possède un objet terminal $1$ ;
\item elle possède tous les produits binaires $A \times B$ ;
\item elle est monoïdale fermée pour la structure
      $(\mathcal{C}, \times, 1)$, c'est-à-dire les exponentielles
      $B^A = [A, B]$ existent :
      \[
        \Hom(A \times B, C) \cong \Hom(A, C^B).
      \]
\end{enumerate}
\end{definition}

\begin{exemple}[Catégories cartésiennes fermées]\label{ex:ccc}
\index{catégorie cartésienne fermée!exemples}
\begin{enumerate}[label=(\roman*)]
\item $\Set$ avec $B^A = \Set(A, B)$.
\item $\Cat$ avec $\mathcal{D}^\mathcal{C} = \Fun(\mathcal{C}, \mathcal{D})$.
\item La catégorie des graphes orientés.
\item La catégorie des ensembles simpliciaux
      $\mathbf{sSet} = \Fun(\Delta^\op, \Set)$.
      \index{ensemble simplicial}
\item Tout topos élémentaire (cf.\ chapitre~\ref{ch:topos}).
\end{enumerate}
\end{exemple}

\begin{proposition}[Évaluation et curryfication]\label{prop:eval-curry}
\index{évaluation (morphisme d')}
\index{curryfication}
Dans une catégorie cartésienne fermée, on dispose de :
\begin{enumerate}[label=(\alph*)]
\item un morphisme d'\textbf{évaluation}
      $\mathrm{ev}_{A,B} \colon B^A \times A \to B$ ;
\item pour tout $f \colon A \times B \to C$, un unique
      $\tilde{f} \colon A \to C^B$ (le \textbf{curryfié} de $f$)
      tel que $\mathrm{ev} \circ (\tilde{f} \times \id_B) = f$.
\end{enumerate}
\end{proposition}

\begin{proof}
Le morphisme d'évaluation est la coünité de l'adjonction
$(-) \times A \dashv (-)^A$.  La curryfication est le passage au
transposé adjoint.
\end{proof}

%% ============================================================
\section{Correspondance de Curry--Howard--Lambek}\label{sec:curry-howard-lambek}
%% ============================================================

\index{correspondance de Curry--Howard--Lambek|(}
\index{Curry--Howard--Lambek|see{correspondance de Curry--Howard--Lambek}}

La correspondance de Curry--Howard--Lambek établit un dictionnaire
entre trois domaines :

\medskip
\begin{center}
\renewcommand{\arraystretch}{1.3}
\begin{tabular}{lll}
\toprule
\textbf{Logique} & \textbf{Typage} & \textbf{Catégories} \\
\midrule
Proposition $A$ & Type $A$ & Objet $A$ \\
Preuve de $A$ & Terme de type $A$ & Morphisme $1 \to A$ \\
$A \Rightarrow B$ & Type fonction $A \to B$ & Exponentielle $B^A$ \\
$A \wedge B$ & Type produit $A \times B$ & Produit $A \times B$ \\
Vrai ($\top$) & Type unité & Objet terminal $1$ \\
Modus ponens & Application & Évaluation \\
Introduction de $\Rightarrow$ & $\lambda$-abstraction & Curryfication \\
\bottomrule
\end{tabular}
\end{center}

\begin{theoreme}[Lambek]\label{thm:lambek}
\index{Lambek!théorème de}
Il y a une équivalence de catégories entre :
\begin{enumerate}[label=(\alph*)]
\item les théories de type du $\lambda$-calcul simplement typé
      (avec produits) ;
\item les catégories cartésiennes fermées.
\end{enumerate}
Plus précisément, la catégorie syntaxique d'un $\lambda$-calcul
simplement typé est cartésienne fermée, et réciproquement,
le langage interne d'une CCC est un $\lambda$-calcul simplement typé.
\index{lambda-calcul@$\lambda$-calcul!simplement typé}
\end{theoreme}

\begin{proof}[Esquisse]
On construit la catégorie syntaxique $\mathcal{C}_T$ d'une théorie $T$ :
les objets sont les types, les morphismes $A \to B$ sont les termes
$x : A \vdash t : B$ modulo $\beta\eta$-équivalence.
Le produit de types donne le produit catégorique,
et le type fonction donne l'exponentielle.
Réciproquement, à une CCC $\mathcal{C}$ on associe la théorie
dont les types sont les objets et dont les termes sont les morphismes.
\end{proof}

\index{correspondance de Curry--Howard--Lambek|)}

%% ============================================================
\section{Exercices}\label{sec:exo-monoidales}
%% ============================================================

\begin{exercice}\label{exo:monoidale-strict}
\index{catégorie monoïdale!stricte}
Montrer que la catégorie $(\mathbf{Mat}_\K, \otimes, 1)$, dont les
objets sont les entiers naturels et les morphismes $m \to n$ sont
les matrices $n \times m$ à coefficients dans $\K$, avec le produit
de Kronecker, est monoïdale stricte.
\end{exercice}

\begin{exercice}\label{exo:monoide-dans-Ab}
Vérifier en détail que les monoïdes dans $(\Ab, \otimes_\Z, \Z)$
sont exactement les anneaux (unitaires).
\end{exercice}

\begin{exercice}\label{exo:ccc-pas-Set}
Montrer que la catégorie $\Top$ des espaces topologiques n'est
\emph{pas} cartésienne fermée.
\textit{Indication} : considérer le produit $\Q \times \Q$ et la
topologie compacte-ouverte.
\index{Top@$\Top$!pas cartésienne fermée}
\end{exercice}

\begin{exercice}\label{exo:enrichie-metrique}
\index{espace métrique!de Lawvere}
Soit $\mathcal{V} = ([0, \infty], \geq)$ munie de la structure
monoïdale $(+, 0)$.  Montrer qu'une $\mathcal{V}$-catégorie est
exactement un espace quasi-métrique généralisé au sens de Lawvere :
$d(x,x) = 0$ et $d(x,z) \leq d(x,y) + d(y,z)$.
\end{exercice}

\begin{exercice}[Groupe de tresses et catégorie tressée]
\label{exo:tresses}
\index{groupe de tresses}
Soit $\mathcal{B}$ la catégorie dont les objets sont les entiers
$n \geq 0$ et dont les morphismes $n \to n$ forment le groupe de
tresses $B_n$.  Montrer que $\mathcal{B}$ est monoïdale (pour
l'addition) et tressée.
\end{exercice}

\begin{exercice}\label{exo:curry-logique}
En utilisant la correspondance de Curry--Howard--Lambek, interpréter
la tautologie $(A \Rightarrow (B \Rightarrow C)) \Leftrightarrow
((A \wedge B) \Rightarrow C)$ comme un isomorphisme dans une CCC.
\end{exercice}

\index{catégorie monoïdale|)}


%%%%%%%%%%%%%%%%%%%%%%%%%%%%%%%%%%%%%%%%%%%%%%%%%%%%%%%%%%%%
%%      CHAPITRE 10 : TOPOS DE GROTHENDIECK               %%
%%%%%%%%%%%%%%%%%%%%%%%%%%%%%%%%%%%%%%%%%%%%%%%%%%%%%%%%%%%%

\chapter{Topos de Grothendieck --- Introduction}\label{ch:topos}
\minitoc

\vspace{1em}

La théorie des topos, initiée par Grothendieck dans le cadre de la
géométrie algébrique et reformulée par Lawvere et Tierney dans un cadre
axiomatique, constitue l'un des sommets de la théorie des catégories.
Un topos unifie la géométrie (faisceaux sur un espace),
l'algèbre (faisceaux sur un site) et la logique (univers de
mathématiques intuitionnistes).  Ce chapitre présente les fondements :
sites, topologies de Grothendieck, faisceaux, topos élémentaires,
classificateur de sous-objets et logique interne.

\index{topos|(}

%% ============================================================
\section{Sites et topologies de Grothendieck}\label{sec:sites}
%% ============================================================

\begin{definition}[Crible]\label{def:crible}
\index{crible}
\index{crible!sous-foncteur}
Soit $\mathcal{C}$ une petite catégorie et $C \in \Ob(\mathcal{C})$.
Un \textbf{crible} sur $C$ est un sous-foncteur
$S \hookrightarrow \Hom_\mathcal{C}(-, C)$
du préfaisceau représentable, c'est-à-dire une famille de morphismes
de but $C$ stable par précomposition : si $f \colon D \to C$ est dans
$S$ et $g \colon E \to D$ est un morphisme quelconque, alors
$f \circ g \in S$.
\end{definition}

\begin{definition}[Topologie de Grothendieck]\label{def:topologie-grothendieck}
\index{topologie de Grothendieck}
\index{site}
Une \textbf{topologie de Grothendieck} $J$ sur une petite catégorie
$\mathcal{C}$ est la donnée, pour chaque objet $C \in \Ob(\mathcal{C})$,
d'une collection $J(C)$ de cribles sur $C$ (appelés \textbf{cribles
couvrants}), satisfaisant les axiomes :
\begin{enumerate}[label=(T\arabic*)]
\item \textbf{Maximalité} : le crible maximal
      $\Hom_\mathcal{C}(-, C)$ est dans $J(C)$.
\item \textbf{Stabilité} : si $S \in J(C)$ et $f \colon D \to C$,
      alors le crible tiré en arrière
      $f^*(S) = \{g \colon E \to D \mid f \circ g \in S\}$
      est dans $J(D)$.
\item \textbf{Transitivité} : si $S \in J(C)$ et $R$ est un crible
      sur $C$ tel que pour tout $f \colon D \to C$ dans $S$ on a
      $f^*(R) \in J(D)$, alors $R \in J(C)$.
\end{enumerate}
Un couple $(\mathcal{C}, J)$ est appelé un \textbf{site}.
\end{definition}

\begin{exemple}[Sites fondamentaux]\label{ex:sites}
\index{site!exemples}
\begin{enumerate}[label=(\roman*)]
\item \textbf{Site ouvert.}\index{site!ouvert}
      Soit $X$ un espace topologique.
      La catégorie $\mathsf{Ouv}(X)$ des ouverts de $X$
      (avec les inclusions) munie de la topologie où un crible
      sur $U$ est couvrant si et seulement si les ouverts qu'il
      contient recouvrent $U$.

\item \textbf{Topologie étale.}\index{topologie étale}
      Sur la catégorie des schémas de type fini sur un corps $k$,
      les familles couvrantes sont les familles surjectives de
      morphismes étales.

\item \textbf{Topologie fppf.}\index{topologie fppf}
      Les familles couvrantes sont les familles surjectives de
      morphismes fidèlement plats de présentation finie.

\item \textbf{Topologie de Zariski.}\index{topologie de Zariski}
      Les familles couvrantes sont les recouvrements par des
      ouverts de Zariski.

\item \textbf{Topologie grossière.}\index{topologie grossière}
      Seul le crible maximal est couvrant.  Tout préfaisceau est
      un faisceau pour cette topologie.

\item \textbf{Topologie canonique.}\index{topologie canonique}
      La plus fine topologie pour laquelle tout foncteur
      représentable est un faisceau.
\end{enumerate}
\end{exemple}

%% ============================================================
\section{Préfaisceaux et faisceaux}\label{sec:faisceaux-site}
%% ============================================================

\begin{definition}[Préfaisceau]\label{def:prefaisceau-site}
\index{préfaisceau!sur une catégorie}
Un \textbf{préfaisceau} (d'ensembles) sur une petite catégorie
$\mathcal{C}$ est un foncteur $F \colon \mathcal{C}^\op \to \Set$.
La catégorie des préfaisceaux est notée
$\widehat{\mathcal{C}} = \Fun(\mathcal{C}^\op, \Set)$.
\index{catégorie de préfaisceaux}
\end{definition}

\begin{definition}[Condition de faisceau]\label{def:faisceau-site}
\index{faisceau!sur un site}
\index{condition de recollement}
Soit $(\mathcal{C}, J)$ un site.  Un préfaisceau
$F \colon \mathcal{C}^\op \to \Set$ est un \textbf{faisceau}
(pour la topologie $J$) si pour tout objet $C$ et tout crible
couvrant $S \in J(C)$, le morphisme canonique
\[
  F(C) \longrightarrow \lim_{\substack{(f \colon D \to C) \in S}} F(D)
\]
est un isomorphisme.  Autrement dit, toute famille compatible d'éléments
indexée par $S$ se recolle de manière unique.
\end{definition}

\begin{remarque}\label{rem:faisceau-egalisateur}
\index{diagramme égalisateur}
Lorsque la topologie est engendrée par des familles couvrantes
$\{f_i \colon U_i \to C\}_{i \in I}$, la condition de faisceau
se reformule par l'exactitude du diagramme égalisateur classique :
\[
\begin{tikzcd}
F(C) \arrow[r] &
\displaystyle\prod_{i \in I} F(U_i)
  \arrow[r, shift left=0.5ex]
  \arrow[r, shift right=0.5ex] &
\displaystyle\prod_{i,j \in I} F(U_i \times_C U_j)
\end{tikzcd}
\]
\end{remarque}

%% ============================================================
\section{Topos de Grothendieck}\label{sec:topos-grothendieck}
%% ============================================================

\begin{definition}[Topos de Grothendieck]\label{def:topos-grothendieck}
\index{topos!de Grothendieck}
\index{Sh@$\mathrm{Sh}(\mathcal{C}, J)$}
Un \textbf{topos de Grothendieck} est une catégorie équivalente
à la catégorie $\mathrm{Sh}(\mathcal{C}, J)$ des faisceaux d'ensembles
sur un site $(\mathcal{C}, J)$.
\end{definition}

\begin{proposition}[Faisceautisation]\label{prop:faisceautisation}
\index{faisceautisation}
\index{foncteur faisceau associé}
L'inclusion $\iota \colon \mathrm{Sh}(\mathcal{C}, J) \hookrightarrow
\widehat{\mathcal{C}}$ admet un adjoint à gauche
$a \colon \widehat{\mathcal{C}} \to \mathrm{Sh}(\mathcal{C}, J)$,
appelé \textbf{foncteur faisceau associé} ou \textbf{faisceautisation}.
L'adjonction $a \dashv \iota$ vérifie que $a$ est exact à gauche.
\end{proposition}

\begin{proof}[Esquisse]
On construit $a$ en deux étapes d'une procédure de type «~plus~» ($+$).
Pour un préfaisceau $F$, on définit
$F^+(C) = \colim_{S \in J(C)} \Hom(S, F)$.
Le foncteur $F \mapsto (F^+)^+$ donne le faisceau associé.
\end{proof}

\begin{exemple}[Topos fondamentaux]\label{ex:topos}
\index{topos!exemples}
\begin{enumerate}[label=(\roman*)]
\item \textbf{Topos des faisceaux sur un espace.}
      \index{topos!de faisceaux sur un espace}
      Pour un espace topologique $X$,
      $\mathrm{Sh}(X) = \mathrm{Sh}(\mathsf{Ouv}(X), J_{\text{ouv}})$
      est le topos classique.

\item \textbf{Topos étale.}
      \index{topos!étale}
      Pour un schéma $X$, le petit site étale donne le topos étale
      $X_{\text{ét}}$, central en cohomologie étale.

\item \textbf{Topos des préfaisceaux.}
      \index{topos!de préfaisceaux}
      $\widehat{\mathcal{C}} = \Fun(\mathcal{C}^\op, \Set)$ est un
      topos (pour la topologie grossière).

\item \textbf{Topos classifiant.}
      \index{topos!classifiant}
      Pour un groupe $G$, la catégorie $G\text{-}\Set$ des $G$-ensembles
      est un topos.
\end{enumerate}
\end{exemple}

%% ============================================================
\section{Topos élémentaire}\label{sec:topos-elementaire}
%% ============================================================

\begin{definition}[Topos élémentaire]\label{def:topos-elementaire}
\index{topos!élémentaire}
\index{Lawvere}
\index{Tierney}
Un \textbf{topos élémentaire} est une catégorie $\mathcal{E}$ vérifiant :
\begin{enumerate}[label=(\roman*)]
\item $\mathcal{E}$ possède toutes les limites finies ;
\item $\mathcal{E}$ est cartésienne fermée ;
\item $\mathcal{E}$ possède un \textbf{classificateur de sous-objets}
      $\Omega$ (cf.\ définition~\ref{def:classificateur}).
\end{enumerate}
\end{definition}

\begin{definition}[Classificateur de sous-objets]\label{def:classificateur}
\index{classificateur de sous-objets}
\index{Omega@$\Omega$ (classificateur)}
\index{morphisme caractéristique}
Un objet $\Omega$ muni d'un monomorphisme
$\mathsf{vrai} \colon 1 \rightarrowtail \Omega$
est un \textbf{classificateur de sous-objets} si pour tout
monomorphisme $m \colon S \rightarrowtail A$, il existe un unique
morphisme $\chi_m \colon A \to \Omega$ (le \textbf{morphisme
caractéristique}) tel que le carré suivant est un tiré en arrière :
\[
\begin{tikzcd}
S \arrow[r] \arrow[d, tail, "m"'] & 1 \arrow[d, tail, "\mathsf{vrai}"] \\
A \arrow[r, "\chi_m"'] & \Omega
\end{tikzcd}
\]
\end{definition}

\begin{exemple}[Classificateur dans $\Set$]\label{ex:omega-set}
\index{classificateur de sous-objets!dans Set}
Dans $\Set$, le classificateur est $\Omega = \{0, 1\}$ avec
$\mathsf{vrai} \colon \{*\} \to \{0,1\}$ envoyant $*$ sur $1$.
Pour un sous-ensemble $S \subseteq A$, le morphisme caractéristique
est la fonction indicatrice $\chi_S$.
\end{exemple}

\begin{proposition}\label{prop:topos-groth-elementaire}
\index{topos!de Grothendieck!est élémentaire}
Tout topos de Grothendieck est un topos élémentaire.  La réciproque
est fausse en général.
\end{proposition}

\begin{proof}[Esquisse]
Un topos de Grothendieck possède toutes les limites (et colimites),
est cartésien fermé, et le foncteur $\mathrm{Sub}(-) \colon
\mathcal{E}^\op \to \Set$ est représentable par un objet $\Omega$.
Pour la non-réciproque, la catégorie des ensembles finis est un
topos élémentaire mais pas de Grothendieck (pas de colimites
quelconques).
\end{proof}

%% ============================================================
\section{Logique interne d'un topos}\label{sec:logique-interne}
%% ============================================================

\index{logique interne|(}
\index{logique intuitionniste}

La structure de topos élémentaire induit une logique interne
qui est, en général, \emph{intuitionniste} : le tiers exclu
$A \vee \neg A$ n'est pas nécessairement valide.

\begin{definition}[Algèbre de Heyting interne]\label{def:heyting-interne}
\index{algèbre de Heyting!interne}
\index{sous-objet!treillis de}
Dans un topos élémentaire $\mathcal{E}$, pour chaque objet $A$,
le treillis des sous-objets $\mathrm{Sub}(A)$ est une algèbre de Heyting.
Les opérations logiques sont :
\begin{enumerate}[label=(\roman*)]
\item $\wedge$ : intersection de sous-objets (tiré en arrière le
      long de la diagonale) ;
\item $\vee$ : union de sous-objets (image de la somme
      amalgamée) ;
\item $\Rightarrow$ : implication de Heyting (adjoint à droite de
      $\wedge$) ;
\item $\neg$ : négation définie par $\neg S = (S \Rightarrow \bot)$.
\end{enumerate}
\end{definition}

\begin{proposition}[Structure algébrique de $\Omega$]\label{prop:omega-heyting}
\index{Omega@$\Omega$!algèbre de Heyting}
Dans un topos élémentaire, l'objet $\Omega$ est un objet d'algèbre
de Heyting interne.  Les morphismes
$\wedge, \vee, \Rightarrow \colon \Omega \times \Omega \to \Omega$
et $\neg \colon \Omega \to \Omega$ satisfont les axiomes d'algèbre
de Heyting.  L'objet $\Omega$ est une algèbre de Boole interne si et
seulement si la logique du topos est classique.
\end{proposition}

\begin{remarque}\label{rem:topos-logique-classique}
\index{topos!booléen}
Un topos $\mathcal{E}$ est dit \textbf{booléen} si $\Omega$ est
une algèbre de Boole interne, c'est-à-dire si
$\neg \neg = \id_\Omega$.  Les conditions suivantes sont
équivalentes :
\begin{enumerate}[label=(\roman*)]
\item $\mathcal{E}$ est booléen ;
\item tout sous-objet admet un complément ;
\item $\mathrm{Sub}(A)$ est une algèbre de Boole pour tout $A$.
\end{enumerate}
Le topos $\Set$ est booléen.  Le topos $\mathrm{Sh}(X)$ pour un
espace $X$ est booléen si et seulement si $X$ est
extrêmement disconnecté (pour les faisceaux de Stone).
\end{remarque}

%% ============================================================
\section{Morphismes géométriques}\label{sec:morphismes-geometriques}
%% ============================================================

\begin{definition}[Morphisme géométrique]\label{def:morphisme-geometrique}
\index{morphisme géométrique}
\index{image directe}
\index{image inverse}
Un \textbf{morphisme géométrique} $f \colon \mathcal{F} \to \mathcal{E}$
entre topos est un couple de foncteurs
$(f^* \colon \mathcal{E} \to \mathcal{F},\;
  f_* \colon \mathcal{F} \to \mathcal{E})$
tel que :
\begin{enumerate}[label=(\roman*)]
\item $f^* \dashv f_*$ (adjonction) ;
\item $f^*$ est exact à gauche (préserve les limites finies).
\end{enumerate}
Le foncteur $f^*$ est l'\textbf{image inverse} et $f_*$ est
l'\textbf{image directe}.
\end{definition}

\begin{exemple}\label{ex:morphismes-geo}
\index{morphisme géométrique!exemples}
\begin{enumerate}[label=(\roman*)]
\item Pour une application continue $f \colon X \to Y$,
      l'image directe et l'image inverse des faisceaux donnent
      un morphisme géométrique $\mathrm{Sh}(X) \to \mathrm{Sh}(Y)$.
\item Le morphisme global $\Gamma \colon \mathcal{E} \to \Set$
      défini par $\Gamma = \Hom_\mathcal{E}(1, -)$ est l'image
      directe d'un morphisme géométrique
      $\gamma \colon \mathcal{E} \to \Set$
      (le \textbf{point global}).
      \index{point global}
\end{enumerate}
\end{exemple}

%% ============================================================
\section{Théorème de Giraud}\label{sec:giraud}
%% ============================================================

\begin{theoreme}[Giraud]\label{thm:giraud}
\index{théorème de Giraud}
\index{Giraud!théorème de}
Soit $\mathcal{E}$ une catégorie.  Les conditions suivantes sont
équivalentes :
\begin{enumerate}[label=(\alph*)]
\item $\mathcal{E}$ est un topos de Grothendieck (i.e.\ $\mathcal{E}
      \simeq \mathrm{Sh}(\mathcal{C}, J)$ pour un petit site
      $(\mathcal{C}, J)$) ;
\item $\mathcal{E}$ vérifie :
  \begin{enumerate}[label=(\roman*)]
  \item $\mathcal{E}$ possède toutes les petites colimites ;
  \item les colimites sont universelles (stables par tiré en arrière) ;
  \item les relations d'équivalence sont effectives ;
  \item les sommes amalgamées sont disjointes ;
  \item $\mathcal{E}$ possède une famille génératrice (petit ensemble).
  \end{enumerate}
\end{enumerate}
\end{theoreme}

\begin{proof}[Idée de preuve]
$(a) \Rightarrow (b)$ : on vérifie chaque condition pour les faisceaux.
$(b) \Rightarrow (a)$ : on prend pour $\mathcal{C}$ une petite
sous-catégorie génératrice et on munit $\mathcal{C}$ de la topologie
canonique.  Le plongement de Yoneda suivi de la faisceautisation
donne une équivalence $\mathcal{E} \simeq \mathrm{Sh}(\mathcal{C}, J)$.
\end{proof}

\index{logique interne|)}

%% ============================================================
\section{Exercices}\label{sec:exo-topos}
%% ============================================================

\begin{exercice}\label{exo:crible}
\index{crible!exercice}
Soit $\mathcal{C}$ une petite catégorie et $C \in \Ob(\mathcal{C})$.
Montrer que les cribles sur $C$ sont en bijection avec les
sous-foncteurs du foncteur représentable $\Hom(-, C)$.
\end{exercice}

\begin{exercice}\label{exo:topologie-triviale}
\index{topologie grossière!exercice}
Vérifier que la topologie grossière (seul le crible maximal est
couvrant) satisfait les axiomes (T1)--(T3).  Montrer que tout
préfaisceau est un faisceau pour cette topologie.
\end{exercice}

\begin{exercice}\label{exo:omega-Sh}
\index{classificateur de sous-objets!dans Sh(X)}
Soit $X$ un espace topologique.  Montrer que dans $\mathrm{Sh}(X)$,
le classificateur de sous-objets est le faisceau
$\Omega(U) = \{V \subseteq U \mid V \text{ ouvert}\}$
pour tout ouvert $U \subseteq X$.
En déduire que $\mathrm{Sh}(X)$ n'est en général pas booléen.
\end{exercice}

\begin{exercice}\label{exo:G-set-topos}
\index{topos!classifiant!exercice}
Soit $G$ un groupe.  Montrer que la catégorie $G\text{-}\Set$ des
$G$-ensembles (à gauche) est un topos de Grothendieck.
Décrire explicitement le classificateur $\Omega$.
\end{exercice}

\begin{exercice}[Morphismes géométriques et points]\label{exo:points-topos}
\index{point d'un topos}
Montrer qu'un morphisme géométrique $p \colon \Set \to \mathcal{E}$
(appelé \textbf{point} du topos $\mathcal{E}$) est déterminé par
un foncteur exact à gauche et qui commute aux petites colimites
$p^* \colon \mathcal{E} \to \Set$.
\end{exercice}

\begin{exercice}\label{exo:giraud-conditions}
Montrer que dans $\Set$ :
(a) les colimites sont universelles,
(b) les relations d'équivalence sont effectives,
(c) les sommes sont disjointes.
\end{exercice}

\index{topos|)}


%%%%%%%%%%%%%%%%%%%%%%%%%%%%%%%%%%%%%%%%%%%%%%%%%%%%%%%%%%%%
%%      CHAPITRE 11 : THÉORIE DES ∞-CATÉGORIES            %%
%%%%%%%%%%%%%%%%%%%%%%%%%%%%%%%%%%%%%%%%%%%%%%%%%%%%%%%%%%%%

\chapter{Théorie des $\infty$-catégories --- Aperçu}\label{ch:infty-cat}
\minitoc

\vspace{1em}

La théorie des catégories classique traite des objets et des morphismes
entre eux.  Mais dans de nombreuses situations (topologie algébrique,
algèbre homotopique, géométrie algébrique dérivée), il est naturel
de considérer des morphismes entre morphismes, des morphismes entre
ceux-ci, et ainsi de suite.  Les \textbf{$\infty$-catégories}
(ou plus précisément les $(\infty,1)$-catégories) formalisent cette
idée : ce sont des catégories avec des morphismes à tous les niveaux,
où les morphismes de niveau $\geq 2$ sont tous inversibles (à homotopie
près).  Ce chapitre final offre un aperçu des fondements, sans
prétendre à l'exhaustivité.

\index{infty-catégorie@$\infty$-catégorie|(}

%% ============================================================
\section{Motivation}\label{sec:infty-motivation}
%% ============================================================

Les catégories ordinaires ne capturent pas toute l'information
homotopique.  Considérons les exemples suivants :

\begin{enumerate}[label=(\roman*)]
\item \textbf{Espaces topologiques.}\index{type d'homotopie}
      La catégorie homotopique $\mathsf{Ho}(\Top)$ (où l'on a inversé
      les équivalences d'homotopie faible) perd de l'information :
      les espaces de morphismes ne sont plus des ensembles mais des
      \emph{espaces}.

\item \textbf{Complexes de chaînes.}\index{catégorie dérivée}
      La catégorie dérivée $D(\mathcal{A})$ d'une catégorie abélienne
      est obtenue en inversant les quasi-isomorphismes, mais cette
      localisation détruit des structures (les triangles distingués
      ne forment pas une bonne notion de limites).

\item \textbf{Champs.}\index{champ (stack)}
      Un champ sur un site est un \og faisceau de catégories\fg{}, ce
      qui demande naturellement des $2$-morphismes et des cohérences
      supérieures.
\end{enumerate}

\begin{remarque}\label{rem:modeles-infty}
\index{modèles d'$\infty$-catégories}
Il existe plusieurs modèles équivalents pour les $(\infty,1)$-catégories :
\begin{enumerate}[label=(\alph*)]
\item les quasi-catégories (Boardman--Vogt, Joyal, Lurie) ;
\item les catégories de Segal complètes (Rezk) ;
\item les catégories simpliciales (Dwyer--Kan) ;
\item les catégories de modèles (Quillen) --- comme présentation.
\end{enumerate}
Ces modèles sont tous reliés par des équivalences de Quillen ou
des équivalences catégoriques appropriées.
\end{remarque}

%% ============================================================
\section{Quasi-catégories}\label{sec:quasi-categories}
%% ============================================================

\index{quasi-catégorie|(}

\begin{notation}\label{not:simplicial}
\index{catégorie simpliciale!Delta@$\Delta$}
\index{ensemble simplicial}
On note $\Delta$ la catégorie simpliciale dont les objets sont les
ensembles ordonnés finis $[n] = \{0, 1, \ldots, n\}$ pour $n \geq 0$,
et les morphismes sont les applications croissantes.
Un \textbf{ensemble simplicial} est un foncteur
$X \colon \Delta^\op \to \Set$.
On note $\mathbf{sSet} = \Fun(\Delta^\op, \Set)$ la catégorie
des ensembles simpliciaux.
\end{notation}

\begin{definition}[Corne]\label{def:corne}
\index{corne}
\index{Lambda@$\Lambda^n_k$ (corne)}
Pour $0 \leq k \leq n$, la \textbf{corne}
$\Lambda^n_k \subset \Delta^n$ est le sous-ensemble simplicial
obtenu en retirant la face intérieure et la $k$-ème face de
$\Delta^n$.  La corne est dite :
\begin{enumerate}[label=(\roman*)]
\item \textbf{intérieure} si $0 < k < n$ ;
\item \textbf{extérieure} si $k = 0$ ou $k = n$.
\end{enumerate}
\end{definition}

\begin{definition}[Quasi-catégorie]\label{def:quasi-categorie}
\index{quasi-catégorie!définition}
\index{complexe de Kan faible}
\index{Joyal}
Un ensemble simplicial $X$ est une \textbf{quasi-catégorie}
(ou \textbf{$\infty$-catégorie} au sens de Joyal--Lurie) s'il
satisfait la \textbf{condition de relèvement de Kan interne} :
pour tout $n \geq 2$ et tout $0 < k < n$, toute application
$\Lambda^n_k \to X$ admet une extension à $\Delta^n$ :
\[
\begin{tikzcd}
\Lambda^n_k \arrow[r] \arrow[d, hook] & X \\
\Delta^n \arrow[ur, dashed] &
\end{tikzcd}
\]
Contrairement à un complexe de Kan (qui exige le relèvement pour
\emph{toutes} les cornes, y compris extérieures), une quasi-catégorie
ne requiert le relèvement que pour les cornes intérieures.
\end{definition}

\begin{exemple}[Quasi-catégories fondamentales]\label{ex:quasi-categories}
\index{quasi-catégorie!exemples}
\begin{enumerate}[label=(\roman*)]
\item \textbf{Nerf d'une catégorie.}\index{nerf d'une catégorie}
      Si $\mathcal{C}$ est une catégorie ordinaire, son nerf
      $N(\mathcal{C})$ est une quasi-catégorie.
      Les $0$-simplexes sont les objets, les $1$-simplexes sont les
      morphismes, et les $n$-simplexes ($n \geq 2$) sont les suites
      composables de $n$ morphismes.
      Le relèvement interne est \emph{unique} car la composition dans
      $\mathcal{C}$ est strictement définie.

\item \textbf{Complexes de Kan.}\index{complexe de Kan}
      Tout complexe de Kan est une quasi-catégorie.
      Les complexes de Kan correspondent aux $\infty$-groupoïdes
      (tous les morphismes sont inversibles à homotopie près).
      L'ensemble simplicial singulier $\mathrm{Sing}(X)$ d'un espace
      topologique $X$ est un complexe de Kan.

\item \textbf{Nerf de Dold--Kan.}
      \index{nerf de Dold--Kan}
      Le nerf de la catégorie dérivée d'une catégorie abélienne
      donne une quasi-catégorie qui encode l'information homotopique
      perdue par la localisation.
\end{enumerate}
\end{exemple}

\begin{proposition}[Catégorie homotopique]\label{prop:cat-homotopique}
\index{catégorie homotopique!d'une quasi-catégorie}
Soit $X$ une quasi-catégorie.  On définit sa \textbf{catégorie
homotopique} $\mathrm{h}X$ comme la catégorie ordinaire dont :
\begin{enumerate}[label=(\roman*)]
\item les objets sont les $0$-simplexes de $X$ ;
\item les morphismes de $x$ à $y$ sont les classes d'homotopie de
      $1$-simplexes de $x$ à $y$ ;
\item la composition est induite par le relèvement des cornes
      $\Lambda^2_1$.
\end{enumerate}
La catégorie $\mathrm{h}X$ est bien définie et
$\mathrm{h}(N(\mathcal{C})) \simeq \mathcal{C}$ pour toute catégorie
ordinaire $\mathcal{C}$.
\end{proposition}

\index{quasi-catégorie|)}

%% ============================================================
\section{$(\infty,1)$-Catégories}\label{sec:infty-1-categories}
%% ============================================================

\index{(infty,1)-catégorie@$(\infty,1)$-catégorie|(}

\begin{definition}[$(\infty,1)$-catégorie (informelle)]
\label{def:infty-1-cat}
\index{(infty,1)-catégorie@$(\infty,1)$-catégorie!définition}
Une \textbf{$(\infty,1)$-catégorie} est une structure possédant :
\begin{enumerate}[label=(\roman*)]
\item des objets ;
\item des $1$-morphismes entre objets ;
\item des $2$-morphismes entre $1$-morphismes ;
\item des $n$-morphismes entre $(n-1)$-morphismes, pour tout $n \geq 2$ ;
\end{enumerate}
où tous les $k$-morphismes pour $k \geq 2$ sont inversibles
(à homotopie près).
\end{definition}

\begin{remarque}\label{rem:infty-vs-quasi}
Le modèle standard pour les $(\infty,1)$-catégories est celui
des quasi-catégories.  Dans la suite, on utilise les termes
\og $\infty$-catégorie\fg{} et \og quasi-catégorie\fg{} de manière
interchangeable.
\end{remarque}

\begin{definition}[Espaces de morphismes]\label{def:mapping-space}
\index{espace de morphismes}
\index{mapping space}
Soit $\mathcal{C}$ une $\infty$-catégorie et $x, y \in \mathcal{C}$.
L'\textbf{espace de morphismes} $\mathrm{Map}_\mathcal{C}(x, y)$
est un espace de Kan (un complexe de Kan) défini par
\[
  \mathrm{Map}_\mathcal{C}(x, y)_n
  = \{f \colon \Delta^{n+1} \to \mathcal{C} \mid
      f|_{\Delta^{\{0\}}} = x,\;
      f|_{\Delta^{\{1,\ldots,n+1\}}} \text{ dégénéré sur } y\}.
\]
Les composantes connexes $\pi_0(\mathrm{Map}_\mathcal{C}(x,y))$
redonnent l'ensemble des morphismes dans $\mathrm{h}\mathcal{C}$.
\end{definition}

\begin{exemple}[$\infty$-catégories fondamentales]\label{ex:infty-cats}
\index{(infty,1)-catégorie@$(\infty,1)$-catégorie!exemples}
\begin{enumerate}[label=(\roman*)]
\item \textbf{$\mathcal{S}$ --- l'$\infty$-catégorie des espaces.}
      \index{S@$\mathcal{S}$ ($\infty$-catégorie des espaces)}
      Les objets sont les complexes de Kan (ou les espaces
      topologiques, à équivalence faible près).  Les espaces de
      morphismes sont les espaces de fonctions.

\item \textbf{$\mathcal{C}\mathrm{at}_\infty$ --- l'$\infty$-catégorie
      des $\infty$-catégories.}
      \index{Cat infty@$\mathcal{C}\mathrm{at}_\infty$}
      C'est l'analogue supérieur de $\Cat$.

\item \textbf{$\mathcal{D}(\mathcal{A})$ --- catégorie dérivée
      $\infty$-catégorique.}
      \index{catégorie dérivée!infty@$\infty$-catégorique}
      Pour une catégorie abélienne $\mathcal{A}$, on obtient une
      $\infty$-catégorie stable qui raffine la catégorie dérivée
      triangulée classique.
\end{enumerate}
\end{exemple}

%% ============================================================
\section{Yoneda $\infty$-catégorique}\label{sec:yoneda-infty}
%% ============================================================

\begin{theoreme}[Lemme de Yoneda $\infty$-catégorique]
\label{thm:yoneda-infty}
\index{lemme de Yoneda!infty-catégorique@$\infty$-catégorique}
\index{Lurie}
Soit $\mathcal{C}$ une petite $\infty$-catégorie.  Le plongement
de Yoneda
\[
  \mathbf{y} \colon \mathcal{C} \longrightarrow
  \Fun(\mathcal{C}^\op, \mathcal{S})
\]
est pleinement fidèle : pour tous $x, y \in \mathcal{C}$, l'application
\[
  \mathrm{Map}_\mathcal{C}(x, y)
  \xrightarrow{\;\sim\;}
  \mathrm{Map}_{\Fun(\mathcal{C}^\op, \mathcal{S})}(\mathbf{y}(x), \mathbf{y}(y))
\]
est une équivalence d'espaces.
\end{theoreme}

\begin{remarque}\label{rem:yoneda-infty-subtil}
La preuve du lemme de Yoneda $\infty$-catégorique est considérablement
plus subtile que dans le cas classique.
Lurie en donne une démonstration dans \emph{Higher Topos Theory}
(Proposition 5.1.3.1) en utilisant les joints d'ensembles
simpliciaux et les extensions de Kan à droite.
\end{remarque}

%% ============================================================
\section{Adjonctions $\infty$-catégoriques}\label{sec:adjonctions-infty}
%% ============================================================

\begin{definition}[Adjonction entre $\infty$-catégories]
\label{def:adjonction-infty}
\index{adjonction!infty-catégorique@$\infty$-catégorique}
Soient $\mathcal{C}$ et $\mathcal{D}$ deux $\infty$-catégories.
Un foncteur $F \colon \mathcal{C} \to \mathcal{D}$ est
\textbf{adjoint à gauche} d'un foncteur
$G \colon \mathcal{D} \to \mathcal{C}$ s'il existe une
équivalence d'espaces
\[
  \mathrm{Map}_\mathcal{D}(F(x), y)
  \simeq
  \mathrm{Map}_\mathcal{C}(x, G(y))
\]
naturelle en $x \in \mathcal{C}$ et $y \in \mathcal{D}$.
\end{definition}

\begin{proposition}[Caractérisation par unité/coünité]
\label{prop:adjonction-infty-unite}
\index{adjonction!unité et coünité ($\infty$)}
L'adjonction $F \dashv G$ dans le cadre $\infty$-catégorique est
équivalente à l'existence de transformations naturelles
$\eta \colon \id_\mathcal{C} \to G \circ F$ et
$\varepsilon \colon F \circ G \to \id_\mathcal{D}$
satisfaisant les identités triangulaires à homotopie (cohérente) près.
\end{proposition}

%% ============================================================
\section{Topos supérieurs}\label{sec:topos-superieurs}
%% ============================================================

\index{topos!supérieur|(}

\begin{definition}[Topos supérieur (d'après Lurie)]
\label{def:topos-superieur}
\index{topos!supérieur!définition}
\index{Lurie!topos supérieur}
Un \textbf{$\infty$-topos} (ou topos supérieur) est une
$\infty$-catégorie $\mathcal{X}$ qui est une localisation exacte
à gauche accessible d'une $\infty$-catégorie de préfaisceaux :
il existe une petite $\infty$-catégorie $\mathcal{C}$ et une
localisation
\[
  L \colon \Fun(\mathcal{C}^\op, \mathcal{S})
  \rightleftarrows \mathcal{X} :\! \iota
\]
avec $L \dashv \iota$ et $L$ exact à gauche.
\end{definition}

\begin{theoreme}[Giraud $\infty$-catégorique (Lurie)]
\label{thm:giraud-infty}
\index{théorème de Giraud!infty-catégorique@$\infty$-catégorique}
Soit $\mathcal{X}$ une $\infty$-catégorie présentable.
Les conditions suivantes sont équivalentes :
\begin{enumerate}[label=(\alph*)]
\item $\mathcal{X}$ est un $\infty$-topos ;
\item les colimites dans $\mathcal{X}$ sont universelles, et les
      objets de groupoïde sont effectifs ;
\item $\mathcal{X}$ vérifie la descente : pour tout diagramme
      simplicial $U_\bullet \to X$, le foncteur
      $\mathcal{X}_{/X} \to \lim_n \mathcal{X}_{/U_n}$ est une
      équivalence.
\end{enumerate}
\end{theoreme}

\begin{exemple}[$\infty$-topos fondamentaux]\label{ex:infty-topos}
\index{topos!supérieur!exemples}
\begin{enumerate}[label=(\roman*)]
\item $\mathcal{S}$, l'$\infty$-catégorie des espaces, est le
      $\infty$-topos terminal (analogue de $\Set$).
\item Pour un espace topologique $X$, la catégorie
      $\mathrm{Sh}_\infty(X)$ des faisceaux d'espaces sur $X$ est
      un $\infty$-topos.
\item Le $\infty$-topos étale d'un schéma intervient en géométrie
      algébrique dérivée.
      \index{géométrie algébrique dérivée}
\end{enumerate}
\end{exemple}

\begin{remarque}[Perspectives]\label{rem:perspectives}
La théorie des $\infty$-catégories et des $\infty$-topos est
aujourd'hui un domaine très actif.  Parmi les directions de
recherche :
\begin{enumerate}[label=(\roman*)]
\item la \textbf{géométrie algébrique dérivée} (Lurie, Toën--Vezzosi) ;
\item la \textbf{topologie algébrique chromatique}
      \index{homotopie chromatique}
      (spectres, localisation chromatique) ;
\item les \textbf{théories des champs topologiques}
      (TQFT, théorie de Cobordisme étendu,
      hypothèse du cobordisme de Baez--Dolan--Lurie) ;
      \index{hypothèse du cobordisme}
\item la \textbf{logique homotopique} et la théorie des types
      homotopiques (HoTT/UF) ;
      \index{HoTT (théorie des types homotopiques)}
\item les \textbf{$(\infty,n)$-catégories} pour $n > 1$ et les
      structures de dimension supérieure.
\end{enumerate}
\end{remarque}

\index{topos!supérieur|)}

%% ============================================================
\section{Exercices}\label{sec:exo-infty}
%% ============================================================

\begin{exercice}\label{exo:nerf-quasi-cat}
\index{nerf d'une catégorie!exercice}
Soit $\mathcal{C}$ une catégorie ordinaire.  Montrer que le nerf
$N(\mathcal{C})$ est une quasi-catégorie, et que les relèvements
de cornes intérieures sont \emph{uniques}.
\end{exercice}

\begin{exercice}\label{exo:kan-infty-groupoide}
\index{complexe de Kan!et $\infty$-groupoïde}
Soit $X$ un complexe de Kan.  Montrer que la catégorie homotopique
$\mathrm{h}X$ est un groupoïde (tout morphisme est un isomorphisme).
\end{exercice}

\begin{exercice}\label{exo:mapping-pi0}
Soit $\mathcal{C}$ une quasi-catégorie et $x, y \in \mathcal{C}$.
Montrer que $\pi_0(\mathrm{Map}_\mathcal{C}(x,y)) \cong
\Hom_{\mathrm{h}\mathcal{C}}(x, y)$.
\end{exercice}

\begin{exercice}\label{exo:infty-topos-S}
\index{S@$\mathcal{S}$!comme $\infty$-topos}
Montrer que l'$\infty$-catégorie $\mathcal{S}$ des espaces est le
$\infty$-topos final : pour tout $\infty$-topos $\mathcal{X}$, il
existe un morphisme géométrique essentiellement unique
$\mathcal{X} \to \mathcal{S}$.
\end{exercice}

\begin{exercice}[Nerf et catégorie homotopique]\label{exo:nerf-h}
Montrer que pour toute catégorie ordinaire $\mathcal{C}$,
le foncteur catégorie homotopique du nerf redonne la catégorie
d'origine : $\mathrm{h}(N(\mathcal{C})) \cong \mathcal{C}$.
\end{exercice}

\begin{exercice}\label{exo:sing-kan}
\index{ensemble simplicial singulier}
Soit $X$ un espace topologique.  Montrer que l'ensemble simplicial
singulier $\mathrm{Sing}(X)$ est un complexe de Kan.
\textit{Indication} : utiliser le fait que les cornes topologiques
sont des rétractes de déformation des simplexes standard.
\end{exercice}

\index{infty-catégorie@$\infty$-catégorie|)}
\index{(infty,1)-catégorie@$(\infty,1)$-catégorie|)}


%%%%%%%%%%%%%%%%%%%%%%%%%%%%%%%%%%%%%%%%%%%%%%%%%%%%%%%%%%%%
%%                     APPENDICE                           %%
%%%%%%%%%%%%%%%%%%%%%%%%%%%%%%%%%%%%%%%%%%%%%%%%%%%%%%%%%%%%

\appendix

\chapter{Rappels d'algèbre et de topologie}\label{app:rappels}
\minitoc

\vspace{1em}

Cet appendice rassemble sans démonstration les résultats d'algèbre
et de topologie utilisés dans le cours.

%% ============================================================
\section{Algèbre}\label{sec:rappels-algebre}
%% ============================================================

\index{rappels d'algèbre|(}

\begin{definition}[Groupe, anneau, module]\label{def:rappel-structures}
\index{groupe!rappel}
\index{anneau!rappel}
\index{module!rappel}
\begin{enumerate}[label=(\roman*)]
\item Un \textbf{groupe} $(G, \cdot, e)$ est un ensemble muni d'une
      loi associative, d'un élément neutre et d'inverses.

\item Un \textbf{anneau} $(R, +, \cdot, 0, 1)$ est un groupe abélien
      $(R, +, 0)$ muni d'une seconde loi associative avec unité $1$,
      distributive par rapport à l'addition.

\item Un \textbf{$R$-module à gauche} est un groupe abélien $(M, +)$
      muni d'une action $R \times M \to M$ compatible avec les
      structures de $R$ et de $M$.
\end{enumerate}
\end{definition}

\begin{proposition}[Produit tensoriel]\label{prop:rappel-tensoriel}
\index{produit tensoriel!rappel}
Soient $R$ un anneau commutatif, $M$ un $R$-module à droite et $N$
un $R$-module à gauche.  Le \textbf{produit tensoriel}
$M \otimes_R N$ est le quotient du $\Z$-module libre sur
$M \times N$ par les relations de bilinéarité et de $R$-équilibre :
\[
  (m \cdot r) \otimes n = m \otimes (r \cdot n).
\]
Il satisfait la propriété universelle :
$\Hom_\Z(M \otimes_R N, P) \cong \mathrm{Bil}_R(M \times N, P)$.
\end{proposition}

\begin{definition}[Suite exacte]\label{def:rappel-suite-exacte}
\index{suite exacte!rappel}
Une suite de morphismes de $R$-modules
$\cdots \to M_{n+1} \xrightarrow{f_{n+1}} M_n
\xrightarrow{f_n} M_{n-1} \to \cdots$
est \textbf{exacte} en $M_n$ si $\ker f_n = \mathrm{im}\, f_{n+1}$.
Une \textbf{suite exacte courte} est une suite exacte
$0 \to A \xrightarrow{f} B \xrightarrow{g} C \to 0$.
\end{definition}

\begin{definition}[Algèbre libre, quotient]\label{def:rappel-algebre-libre}
\index{algèbre libre!rappel}
\index{quotient!rappel}
Soit $R$ un anneau commutatif.
\begin{enumerate}[label=(\roman*)]
\item Le $R$-module libre sur un ensemble $S$ est
      $R^{(S)} = \bigoplus_{s \in S} R$.
\item Pour un sous-module $N \subseteq M$, le quotient $M/N$ est
      un $R$-module avec la surjection canonique $\pi \colon M \to M/N$.
\end{enumerate}
\end{definition}

\index{rappels d'algèbre|)}

%% ============================================================
\section{Topologie}\label{sec:rappels-topologie}
%% ============================================================

\index{rappels de topologie|(}

\begin{definition}[Espace topologique]\label{def:rappel-topologique}
\index{espace topologique!rappel}
Un \textbf{espace topologique} est un couple $(X, \tau)$ où $\tau$
est une collection de sous-ensembles de $X$ (les \textbf{ouverts})
stable par unions quelconques et intersections finies, contenant
$\varnothing$ et $X$.
\end{definition}

\begin{definition}[Application continue, homéomorphisme]
\label{def:rappel-continue}
\index{application continue!rappel}
\index{homéomorphisme!rappel}
Une application $f \colon X \to Y$ entre espaces topologiques est
\textbf{continue} si $f^{-1}(V)$ est ouvert dans $X$ pour tout
ouvert $V$ de $Y$.  Un \textbf{homéomorphisme} est une bijection
continue dont l'inverse est continue.
\end{definition}

\begin{definition}[Compacité, connexité]\label{def:rappel-compact-connexe}
\index{compacité!rappel}
\index{connexité!rappel}
\begin{enumerate}[label=(\roman*)]
\item $X$ est \textbf{compact} si de tout recouvrement ouvert on peut
      extraire un sous-recouvrement fini.
\item $X$ est \textbf{connexe} s'il ne peut pas s'écrire comme
      union de deux ouverts disjoints non vides.
\end{enumerate}
\end{definition}

\begin{definition}[Préfaisceau et faisceau sur un espace]
\label{def:rappel-faisceau}
\index{faisceau!rappel}
\index{préfaisceau!rappel}
Soit $X$ un espace topologique.
\begin{enumerate}[label=(\roman*)]
\item Un \textbf{préfaisceau} (d'ensembles) sur $X$ est un foncteur
      $F \colon \mathsf{Ouv}(X)^\op \to \Set$.
\item Un préfaisceau $F$ est un \textbf{faisceau} s'il satisfait la
      condition de recollement : pour tout recouvrement ouvert
      $\{U_i\}$ de $U$, le diagramme
      \[
        F(U) \to \prod_i F(U_i) \rightrightarrows
        \prod_{i,j} F(U_i \cap U_j)
      \]
      est un égalisateur.
\end{enumerate}
\end{definition}

\begin{definition}[Homotopie, équivalence d'homotopie]
\label{def:rappel-homotopie}
\index{homotopie!rappel}
\index{équivalence d'homotopie!rappel}
\begin{enumerate}[label=(\roman*)]
\item Deux applications continues $f, g \colon X \to Y$ sont
      \textbf{homotopes} s'il existe une application continue
      $H \colon X \times [0,1] \to Y$ avec $H(-, 0) = f$ et
      $H(-, 1) = g$.
\item $f \colon X \to Y$ est une \textbf{équivalence d'homotopie}
      s'il existe $g \colon Y \to X$ avec $g \circ f \simeq \id_X$
      et $f \circ g \simeq \id_Y$.
\end{enumerate}
\end{definition}

\begin{definition}[Groupe fondamental]\label{def:rappel-pi1}
\index{groupe fondamental!rappel}
Pour un espace topologique pointé $(X, x_0)$, le \textbf{groupe
fondamental} $\pi_1(X, x_0)$ est le groupe des classes d'homotopie
de lacets basés en $x_0$.
\end{definition}

\index{rappels de topologie|)}

%% ============================================================
%%                    BIBLIOGRAPHIE
%% ============================================================

\begin{thebibliography}{99}
\bibitem{maclane} Mac\,Lane, S., \emph{Categories for the Working Mathematician}, 2nd ed., Springer, 1998.
\bibitem{freyd1964} Freyd, P., \emph{Abelian Categories: An Introduction to the Theory of Functors}, Harper \& Row, 1964.
\bibitem{mitchell1965} Mitchell, B., \emph{Theory of Categories}, Academic Press, 1965.
\bibitem{weibel1994} Weibel, C.A., \emph{An Introduction to Homological Algebra}, Cambridge University Press, 1994.
\bibitem{borceux} Borceux, F., \emph{Handbook of Categorical Algebra} (3 volumes), Cambridge University Press, 1994.
\bibitem{riehl} Riehl, E., \emph{Category Theory in Context}, Dover, 2016.
\bibitem{awodey} Awodey, S., \emph{Category Theory}, 2nd ed., Oxford University Press, 2010.
\bibitem{kashiwara} Kashiwara, M. and Schapira, P., \emph{Categories and Sheaves}, Springer, 2006.
\bibitem{grothendieck} Grothendieck, A., \emph{Sur quelques points d'alg\`ebre homologique}, T\^ohoku Math.\ J., 9, 119--221, 1957.
\end{thebibliography}

%% ============================================================
%%                       INDEX
%% ============================================================

\printindex

\end{document}
